\chapter{Proofs for Finite-Depth Formation}
\label{app:finite-depth-formation-proofs}

This appendix proves the propositions of \cref{ch:finite-depth-formation},
the lemmas behind \Cref{thm:finite-depth-adaptive-formation},
\Cref{prop:rg-formation-endogenous-cycle}, and its extension
\Cref{prop:phil-seeded-contrast-formation} of \cref{ch:philosophy}, in the symbols of
\cref{tab:rg-formation-notation}; \Cref{tab:rg-formation-lemma-dependencies}
orders the lemmas. The depth \(L\ge2\) is fixed. Constants
\(\mathrm c_L,\mathrm C_L>0\) may change between displays; they depend on
\(L\), on the fixed data, and, once chosen, on \(k_L\) and \(\eta_L\), but not
on \(\epsilon\), \(\mathrm A\), \(\mathrm T\), or the strengths unless the
dependence is displayed, and all limits have \(\epsilon,\mathrm T>0\).
Densities are relative to Lebesgue measure times counting measure, values of
the complete assignment are written \(j,k\in\mathcal J\), and the independent
equilibrium root factor of every law is omitted from non-root integrals and
appended unchanged. On \(M_\ell\), \(H_\ell\) is the Gaussian Schur precision
and \(\mathrm d_\ell=\dim M_0-\dim M_\ell\) the number of eliminated
coordinates.

\begin{sciencetable}{Lemma dependencies of the finite-depth formation theorem.}
\label{tab:rg-formation-lemma-dependencies}
\footnotesize
\begin{tabularx}{\textwidth}{>{\raggedright\arraybackslash}p{0.11\textwidth}>{\raggedright\arraybackslash}p{0.27\textwidth}>{\raggedright\arraybackslash}X>{\raggedright\arraybackslash}p{0.23\textwidth}}
\toprule
\textbf{Lemma} & \textbf{Inputs} & \textbf{Output} & \textbf{Exponent or constant} \\
\midrule
\Cref{lem:formation-parent-score} & \Cref{lem:rg-formation-carrier} & parent potentials, rows, covector; force, Hessian, metric bounds & \(B_*=\delta\sqrt{2N}\), \(\delta_*<1/2\) \\
\Cref{lem:formation-preparation-margin} & \Cref{lem:rg-formation-carrier} & prepared direction, strict margin & \(\mathrm m_L>0\) \\
\Cref{lem:formation-prepared-continuation} & \Cref{lem:formation-parent-score} & well-posedness, moments, parameter bounds & \(\mathrm A^2\mathrm J_{\ell+1}+\mathrm T\varsigma_\epsilon^2\) \\
\Cref{lem:formation-parameter-comparison} & \Cref{lem:formation-prepared-continuation} & calibration-dominated parameter & \(r_\epsilon=O(\epsilon^{2\mathrm Q_1+20})\) \\
\Cref{lem:formation-entropy-tools} & cited inequalities & log-Sobolev, Herbst, duality, transport, Pinsker & curvature, \(\mathrm T\) explicit \\
\Cref{lem:formation-activation-entropy} & \Cref{lem:formation-parent-score,lem:formation-prepared-continuation,lem:formation-entropy-tools} & activation and Fisher budgets; score moments & \(\mathrm A^2\mathrm J_{\ell+1}+\mathrm T\varsigma_\epsilon^2\) \\
\Cref{lem:formation-conditional-entropy} & \Cref{lem:formation-entropy-tools,lem:formation-activation-entropy} & early discarded loss; covector error & \(\epsilon^{14}\); \(\tfrac12\mathrm c_{\mathrm{early},L}k_L>\mathrm Q_1+20\) \\
\Cref{lem:formation-terminal-joint-entropy} & \Cref{lem:formation-parent-score,lem:formation-activation-entropy} & terminal loss by joint relaxation & \(\epsilon^2\); \(\mathrm c_{\mathrm{joint},L}k_L>\mathrm Q_1+8\) \\
\Cref{lem:formation-fast-marginal} & \Cref{lem:formation-parent-score} & averaged row; environment errors & \(\mathrm e_\epsilon(y)\) \\
\Cref{lem:formation-fast-mixing} & \Cref{lem:formation-fast-marginal,lem:formation-prepared-continuation} & mixing to the moving fast equilibrium & rate \(\mathrm c_{\mathrm{fast},L}\mathrm J_1\) \\
\Cref{lem:formation-history-coupling} & \Cref{lem:formation-fast-marginal,lem:formation-fast-mixing,lem:formation-parameter-comparison} & assignment histories & \(\epsilon^6\); \(\mathrm c_{\mathrm{fast},L}k_L>\mathrm Q_1/2+12\) \\
\Cref{lem:formation-coherence} & \Cref{lem:formation-preparation-margin,lem:formation-prepared-continuation} & coherent components & \(e^{-\mathrm c_L/\epsilon^6}\); \(k_{\mathrm{geom},L}\), \(\eta_L^{\max}\) \\
\Cref{lem:formation-source-envelope} & \Cref{lem:formation-parameter-comparison} & source movement; kernel change & \(57/100\); \(1/200\) \\
\Cref{lem:formation-continuous-path} & \Cref{lem:formation-coherence,lem:formation-parameter-comparison} & common-path concentration & \(\mathrm A\epsilon^2\); \(e^{-\mathrm c_L/\epsilon^2}\) \\
\bottomrule
\end{tabularx}
\end{sciencetable}

\section{Exact marginals, rates, and parameter derivatives}
\label{sec:rg-formation-exact-proof}

\noindent\emph{Proof of \Cref{prop:rg-formation-normalized-experiment}.}
Under \(\nu_i^*\), \(k_i-m_i=\mathfrak d_i+\xi'_i\), whereas the physical
evaluator gives \(k_i-m_i=\xi'_i\) with the same law of \(m_i\). The
log-density ratio \([\mathfrak d_i(k_i-m_i)-\mathfrak d_i^2/2]/(1-\sigma^2_{\mathrm{ev}})\)
therefore has expectation \(\mathfrak d_i^2/(2(1-\sigma^2_{\mathrm{ev}}))\)
under every coupling of the two marginals, since
\(\E[k_i-m_i]=\bar k_i-\bar m_i\); this proves
\cref{eq:rg-formation-evaluator-optimum}, and summation gives
\(\mathbf x^\top R_{\mathrm{ev}}\mathbf x/2\). The rows \(\mathfrak d_i\) belong
to \(B_1\), so \(\mathrm J_1\mathfrak D^\top\mathfrak D\preceq H\), and
\(\mathrm T\le\mathrm J_1(1-\sigma^2_{\mathrm{ev}})/2\) gives
\(\mathrm TR_{\mathrm{ev}}\preceq H/2\); with \(\vartheta\ge1\) and the
orthogonal split of \Cref{lem:rg-formation-carrier},
\(\Omega_0(\vartheta)\succeq(\vartheta/2)H+P_{\mathcal K}\succ0\). Every factor
of \cref{eq:rg-formation-full-joint} is a positive normalized kernel of inputs
none of which is a recognition law, so iterated integration gives total mass
one and the records a positive density. \(\square\)

\noindent\emph{Proof of \Cref{prop:rg-formation-laws}.}
Since \(\widehat K_\vartheta-\mathrm TR_{\mathrm{ev}}\succeq(\vartheta/2)H
+(1+\kappa)P_{\mathcal K}\succ0\) and \(0\le\phi_S\le B_0\), both integrals are
finite and positive. At fixed \(\mathbf x\) both laws are proportional to
\(\pi_S(\vartheta)e^{-\phi_S(\mathbf x)}\) in \(S\), which gives
\cref{eq:rg-formation-full-source-row}, a product of receiver rows because
\(\phi_S\) is a sum and \(\pi_S\) a product over receivers. Define
\begin{equation}
\begin{aligned}
\mathrm n_{\mathcal K}&=\bigl(\kappa/(2\pi\mathrm T)\bigr)^{N/2},&
\mathrm n_{\mathrm{cal}}&=(2\pi\mathrm T)^{-1/2},\\
\mathrm B_{\mathrm T}(\vartheta)
 &=\mathrm T\log Z_0(\vartheta)
   -\mathrm T\log\bigl(\mathrm n_{\mathcal K}\mathrm n_{\mathrm{cal}}q_0^N\bigr)
   +\tfrac12\kappa|z_\star|^2+\tfrac12(\vartheta-2)^2,\\
\mathrm C_{\mathrm T}(\vartheta)&=\mathrm B_{\mathrm T}(\vartheta)
   -\mathrm T\log Z_{\mathrm{opt}}(\vartheta),&
p_\vartheta(\mathsf o)&=e^{-\mathrm B_{\mathrm T}(\vartheta)/\mathrm T}
                         Z_{\mathrm{phys}}(\vartheta).
\end{aligned}
\label{eq:rg-formation-evidence-scalar}
\end{equation}
As \(z_\star\in\mathcal K\), the log density of the joint at the records is
\(-[\mathrm E^{\mathrm{phys}}_\vartheta+\mathrm T\phi_S+\tfrac12\kappa|z_\star|^2
+\tfrac12(\vartheta-2)^2]/\mathrm T-\log Z_0
+\log(\mathrm n_{\mathcal K}\mathrm n_{\mathrm{cal}}q_0^N)+\log\pi_S\) plus the
physical evaluator density. Hence \(\mathrm T\E\log(q\prod_i\nu_i^*/p)\), in
which \cref{eq:rg-formation-evaluator-optimum} adds
\(\tfrac{\mathrm T}2\mathbf x^\top R_{\mathrm{ev}}\mathbf x
=\mathrm E_\vartheta-\mathrm E^{\mathrm{phys}}_\vartheta\), is the first line
of \cref{eq:rg-formation-restricted-vfe}, and completing the relative entropy
to \(q^\star_\vartheta\) gives the second. Integrating the joint over
\((\mathbf x,S)\) at the records gives \(p_\vartheta(\mathsf o)\), so
\(\mathrm C_{\mathrm T}+\mathrm T\log p_\vartheta(\mathsf o)
=\mathrm T\log(Z_{\mathrm{phys}}/Z_{\mathrm{opt}})\ge0\), as
\(\mathrm E^{\mathrm{phys}}_\vartheta\le\mathrm E_\vartheta\).

For the gap, \(Z_{\mathrm{phys}}/Z_{\mathrm{opt}}=
\E_{q^\star_\vartheta}e^{\mathbf x^\top R_{\mathrm{ev}}\mathbf x/2}\). The
Gaussian law \(\gamma^{\mathrm E}_\vartheta\propto e^{-\mathrm E_\vartheta/\mathrm T}\)
has its mean in \(\mathcal K\), which \(R_{\mathrm{ev}}\) annihilates, and since
\(R_{\mathrm{ev}}\) has rank \(N\) the eigenvalues \(\lambda_k\) of
\(\mathrm T\widehat K_\vartheta^{-1}R_{\mathrm{ev}}\) satisfy
\(\lambda_k\le\mathrm T/(\mathrm J_1(1-\sigma^2_{\mathrm{ev}}))\le1/2\) and
\(\sum_k\lambda_k\le N\mathrm T/(\mathrm J_1(1-\sigma^2_{\mathrm{ev}}))\le1\), so
\(\E_{\gamma^{\mathrm E}_\vartheta}e^{\mathbf x^\top R_{\mathrm{ev}}\mathbf x/2}
=\prod_k(1-\lambda_k)^{-1/2}\le e^{\sum_k\lambda_k}\). The density of
\(q^\star_\vartheta\) relative to \(\gamma^{\mathrm E}_\vartheta\otimes\pi(\vartheta)\)
is at most \(e^{B_0}\) and \(e^{\mathbf x^\top R_{\mathrm{ev}}\mathbf x/2}\ge1\),
so \(\E_{q^\star_\vartheta}[e^{\mathbf x^\top R_{\mathrm{ev}}\mathbf x/2}-1]
\le e^{B_0}(e^{\sum_k\lambda_k}-1)\le2e^{B_0}\sum_k\lambda_k\), and
\(\log(1+y)\le y\) gives the first bound of
\cref{eq:rg-formation-physical-optimizer-gap}. The second follows from
\(\KL(q^\star_\vartheta\Vert\boldsymbol\Pi_{\vartheta,\mathsf o})
=\log(Z_{\mathrm{phys}}/Z_{\mathrm{opt}})
-\tfrac12\E_{q^\star_\vartheta}\mathbf x^\top R_{\mathrm{ev}}\mathbf x\).
\(\square\)

\lemmaheading{Weighted geometry and exact parent score}{lem:formation-parent-score}
For the triangle-chain experiment, all marked marginal laws and inherited
proposal rows are normalized and compose exactly. With orthonormal flag
columns \(\mathsf Q_\ell\) spanning \(\mathcal U_{\ell-1}\cap\mathcal U_\ell^\perp\),
\(\mathsf B=B[\mathsf Q_1\cdots\mathsf Q_L]\), and
\(\mathbf x_{\mathcal U}=\sum_\ell\mathsf Q_\ell w_\ell\),
\begin{equation}
\|\mathsf B^{-1}\|_F^{-2}\sum_\ell\mathrm J_\ell|w_\ell|^2
\le \mathbf x_{\mathcal U}^{\top}H\mathbf x_{\mathcal U}
\le\|\mathsf B\|_F^2\sum_\ell\mathrm J_\ell|w_\ell|^2,
\label{eq:rg-formation-weighted-flag}
\end{equation}
and the same bounds hold for every Schur complement obtained by minimizing over
an initial group of normal layers. On \(1\le\vartheta\le21/10\), the parent
potentials are \cref{eq:rg-formation-parent-potential} with the tilt and
eliminated Gaussian scalar of
\cref{eq:rg-formation-parent-tilt-proof,eq:rg-formation-hidden-scalar-proof}
below, the inherited rows are those of \cref{eq:rg-formation-parent-generator}
and satisfy the detailed balance of \Cref{prop:rg-formation-rate-composition},
and the complete covector, including the source-prior score, is
\cref{eq:rg-formation-complete-score}. With \(B_*=\delta\sqrt{2N}\) and
\(\delta_*=2\delta+2N\delta^2<1/2\), the parent forces, Hessians, and scaled
metric obey
\begin{align}
 &|H_\ell^{-1/2}\mathrm T\nabla\psi_{\ell,j}|\le B_*\sqrt{\mathrm T},
 \qquad
 (\vartheta-\delta_*)H_\ell
 \preceq\nabla^2V_{\ell,j}
 \preceq(\vartheta+2\delta)H_\ell,
 \label{eq:rg-formation-parent-relative-bounds}\\
 &1\le\mathsf m(\vartheta)\le1+\mathrm C_L\mathrm T,\qquad
 |\mathsf m'(\vartheta)|\le\mathrm C_L\mathrm T,
 \label{eq:rg-formation-metric-bounds-proof}
\end{align}
for every fixed depth and all sufficiently small \(\epsilon\).
\status{ESTABLISHED}

\noindent\emph{Proof.}
By \Cref{lem:rg-formation-carrier}, rows of level at most \(\ell\) annihilate
\(\mathcal U_\ell\), so \(\mathsf B\) and \(\mathsf B^{-1}\) are block lower
triangular with invertible diagonal blocks. With
\(\mathsf D_{\mathrm J}=\operatorname{diag}(\sqrt{\mathrm J_\ell}I)\), each
nonzero block \((\ell,r)\) of
\(\mathsf D_{\mathrm J}\mathsf B^{\pm1}\mathsf D_{\mathrm J}^{-1}\) is that of
\(\mathsf B^{\pm1}\) times \(\sqrt{\mathrm J_\ell/\mathrm J_r}\le1\), which
proves \cref{eq:rg-formation-weighted-flag}, and \(H|_{M_0}=H_{\mathcal U}\oplus
\mathrm J_1I_{\mathcal H}\). The Schur quadratic is the infimum over the
eliminated coordinates, so minimizing both bounds proves the retained
versions; in particular \(H_\ell\succeq\mathrm h_LI\) for a fixed
\(\mathrm h_L>0\), since \(\mathrm J_L=1\), the normal part of \(H_\ell\) is at
most \(\mathrm C_L\mathrm J_{\ell+1}I\) with \(\mathrm J_{L+1}=0\), and the loop
part keeps the weight \(\mathrm J_1\).

For a cumulative hidden/retained split \((z,y)\) of \(M_0\) with blocks
\(H^{(\ell)}_{zz}\), \(H^{(\ell)}_{zy}\) of \(H\), completing the square gives
the lift in the first line below, with \(z\) embedded in the hidden subspace.
The centered hidden Gaussian \(\gamma^{\mathrm{hid}}_{\ell,\vartheta}\), of
covariance \((\mathrm T/\vartheta)(H^{(\ell)}_{zz})^{-1}\), has a mean
independent of \(\vartheta\), and integrating the source factor gives
\cref{eq:rg-formation-parent-potential} with the tilt, the scalar, and the
normalized optimizer of the other lines, \(\mathrm C_0^{\mathrm G}=0\):
\begin{align}
 \Xi_\ell y&=\bigl(-(H^{(\ell)}_{zz})^{-1}H^{(\ell)}_{zy}y,\ y\bigr),\qquad
 \Xi_\ell^\top H\Xi_\ell=H_\ell,\qquad \Xi_\ell^\top Hz=0,
 \label{eq:rg-formation-gaussian-lift}\\
 \psi_{\ell,j}(y,\vartheta)
 &=-\log\int e^{-\phi_j(\Xi_\ell y+z)}
                  \gamma^{\mathrm{hid}}_{\ell,\vartheta}(dz),
 \qquad 0\le\psi_{\ell,j}\le B_0,
 \label{eq:rg-formation-parent-tilt-proof}\\
 \mathrm C_\ell^{\mathrm G}(\vartheta)
 &=-\frac{\mathrm T\mathrm d_\ell}{2}\log\frac{2\pi\mathrm T}{\vartheta}
       +\frac{\mathrm T}2\log\det H^{(\ell)}_{zz},
 \label{eq:rg-formation-hidden-scalar-proof}\\
 q^{\star,\ell}_\vartheta(y,j)
 &=\frac{e^{-\psi_{\ell,j}(y,\vartheta)}
       \gamma_{\ell,\vartheta}(y)\pi_j(\vartheta)}{Z_\ell(\vartheta)},\quad
 \gamma_{\ell,\vartheta}=\mathcal N\bigl(0,(\mathrm T/\vartheta)H_\ell^{-1}\bigr),\quad
 e^{-B_0}\le Z_\ell\le1.
 \label{eq:rg-formation-normalized-parent-proof}
\end{align}
Positivity and Gaussian domination justify iterated integration and
differentiation on compact parameter intervals.
\Cref{prop:rg-formation-rate-composition} gives detailed balance, composition,
and \(\beta^\ell_{jk}\ge e^{-B_0}\pi_k(\vartheta)\) for the rows of
\cref{eq:rg-formation-parent-generator}, and the physical, optimizing, and
recognition laws compose by separate integration through the same channel.

Let \(\mathfrak B_{\mathrm{int}}\) have rows \(\sqrt{\mathrm J_1}\mathbf r_e\)
on internal edges, \(\mathbf r_e\) the covector of \(r_e\), so
\(\mathfrak B_{\mathrm{int}}^\top\mathfrak B_{\mathrm{int}}\preceq H\). An edge is
selected \(\mathrm n_e\in\{0,1,2\}\) times, \(\sum_e\mathrm n_e=N\), so the
force of \(\phi_j=\delta\sum_i\chi(\sqrt{\mathrm J_1/\mathrm T} r_{ij_i})\) is
\(\delta\sqrt{\mathrm T}\mathfrak B_{\mathrm{int}}^\top\mathfrak c\) with
\(\mathfrak c_e=\mathrm n_e\chi'(\sqrt{\mathrm J_1/\mathrm T} r_e)\) and, since
\(|\chi'|,|\chi''|\le1\), \(|\mathfrak c|^2\le2N\). This gives the first line
below on \(M_0\), and with the centered Gaussian variable of
\cref{eq:rg-formation-parent-tilt-proof} held fixed, the normalized tilted law
\(\nu_{\ell,y,j}\) gives the other two:
\begin{align}
 &|H^{-1/2}\mathrm T\nabla\phi_j|\le B_*\sqrt{\mathrm T},
 \qquad
 -2\delta H\preceq \mathrm T\nabla^2\phi_j\preceq2\delta H,
 \qquad B_*=\delta\sqrt{2N},
 \label{eq:rg-formation-original-relative-bounds}\\
 &\mathrm T\nabla\psi_{\ell,j}
 =\E_{\nu_{\ell,y,j}}\Xi_\ell^\top(\mathrm T\nabla\phi_j),
 \label{eq:rg-formation-marginal-force-proof}\\
 &\mathrm T\nabla^2\psi_{\ell,j}
 =\E_{\nu_{\ell,y,j}}\Xi_\ell^\top(\mathrm T\nabla^2\phi_j)\Xi_\ell
       -\mathrm T^{-1}\operatorname{Cov}_{\nu_{\ell,y,j}}
                    (\Xi_\ell^\top\mathrm T\nabla\phi_j).
 \label{eq:rg-formation-marginal-hessian-proof}
\end{align}
Through the isometry \(H^{1/2}\Xi_\ell H_\ell^{-1/2}\), the first line bounds
the force by \(B_*\sqrt{\mathrm T}\) in the \(H_\ell\) norm, the covariance by
\(B_*^2H_\ell=2N\delta^2H_\ell\), and the first Hessian term between
\(\pm2\delta H_\ell\), which proves
\cref{eq:rg-formation-parent-relative-bounds}. The source factor is even, so
centered Gaussian integration makes \(\psi_{\ell,j}\) even in \(y\), and every
non-root sheet has its unique minimum at zero.

The hidden Gaussian scale score
\(\mathfrak s_\ell=(\mathrm d_\ell-|\mathbf Z|^2)/(2\vartheta)\), with
\(\mathbf Z\) standard Gaussian, has mean zero and second moment
\(\mathrm d_\ell/(2\vartheta^2)\), and the tilted density ratio lies in
\([e^{-B_0},e^{B_0}]\); as the Gaussian conditional center is independent of
\(\vartheta\), uniformly in the retained input,
\begin{equation}
 \partial_\vartheta\psi_{\ell,j}=-\E_{\nu_{\ell,y,j}}\mathfrak s_\ell,\qquad
 |\partial_\vartheta\psi_{\ell,j}|
 \le(e^{B_0}-1)\frac{\sqrt{\mathrm d_\ell}}{\sqrt2 \vartheta}.
 \label{eq:rg-formation-tilt-score-proof}
\end{equation}
By \cref{eq:rg-formation-parent-potential} the parent free energy is
\(\mathrm T\E_q\log q+\E_q[V_{\ell,S}-\mathrm T\log\pi_S]+\mathrm B_{\mathrm T}\),
with \(\mathrm B_{\mathrm T}'=\vartheta-2-\tfrac{\mathrm T}2\Tr(\Omega_0^{-1}H)\),
and differentiation at a fixed recognition law, which does not differentiate
assignment masses, gives \cref{eq:rg-formation-complete-score}, whose source-prior term, present
although it has no continuous gradient, has absolute unscaled score at most
four for the distinguished binary prior. In the normalized joint the prior, calibration, and source-prior
scores have zero mean and zero cross expectations, the other kernels have no
explicit parameter dependence at fixed inputs, and the three scores have
variances \(\tfrac12\Tr[(\Omega_0^{-1}H)^2]\), \(1/\mathrm T\), and, for
\(4(\mathbf 1_{\{S_{i_0}=+\}}-\pi_{i_0,+})\), \(16\pi_{i_0,+}(1-\pi_{i_0,+})\),
which gives \cref{eq:rg-formation-model-fisher}. As \(R_{\mathrm{ev}}\) is
fixed, annihilates \(\mathcal K\), and is supported on the first normal layer,
\(\Omega_0\) is positive for small \(\mathrm T\), and conjugation by the
quadratic core bounds the trace terms and their derivatives uniformly;
calibration contributes exactly one to \(\mathsf m\) and the other two terms
are nonnegative and \(O(\mathrm T)\), which proves
\cref{eq:rg-formation-metric-bounds-proof}. \hfill\(\square\)

\noindent\emph{Proof of \Cref{prop:rg-formation-score-dissipation}.}
The metric and the covector are derived in \Cref{lem:formation-parent-score}.
The diffusion part of the generator is reversible with respect to
\(e^{-V_{\ell,S}/\mathrm T}\) on each sheet, which gives the first term of the
dissipation after integration by parts. For the jump part, detailed balance
(\Cref{prop:rg-formation-rate-composition}) turns
\(\mathrm T\int\sum_S\partial_tq_S\log\mathfrak r_S\) into
\(\mathrm T\int\sum_{S,S'}q^{\star,\ell}(y,S)\beta^\ell_{SS'}
(\mathfrak r_{S'}-\mathfrak r_S)\log\mathfrak r_S\), and symmetrizing in
\(S,S'\) gives the second term, whose integrand is nonnegative because the
logarithm is increasing. The parameter term is
\(g_\ell\dot\vartheta_\ell=-\epsilon\mathsf m^{-1}g_\ell^2\). At \(\ell=0\) the
row \(\beta_{S'}\) does not depend on \(S\), and with
\(\mathfrak r_S=(q_{\mathbf X}/q^\star_{\mathbf X}) q(S\mid\mathbf x)/\beta_S\)
the jump integrand sums to
\(2[\KL(q(\cdot\mid\mathbf x)\Vert\beta)+\KL(\beta\Vert q(\cdot\mid\mathbf x))]\).
These identities hold for positive smooth densities with finite entropy and
moments; the moment bounds of \Cref{lem:formation-prepared-continuation} and
the cutoff passage of \Cref{lem:formation-activation-entropy} justify them for
the admitted solutions. \(\square\)

\section{Prepared laws and coupled continuation}
\label{sec:rg-formation-continuation-proof}

\lemmaheading{Prepared direction with a strict margin}{lem:formation-preparation-margin}
Consider the list of linear functionals on \(\mathcal U\) consisting of each
internal edge covector and, for each bridge \(e\) and each \(\ell<\ell(e)\),
the functional \(r_e\circ P_\ell|_{\mathcal U}\). Every functional in the list
is nonzero, and there is a unit vector \(\mathbf v_\star\in\mathcal U\) with
\begin{equation}
\mathrm m_L=\min_{e,\ \ell<\ell(e)}
       |r_e(P_\ell\mathbf v_\star)|>0,
\label{eq:rg-formation-preparation-margin}
\end{equation}
where internal edges enter only at \(\ell=0\). A fixed neighborhood of
\(\mathbf v_\star\) keeps a margin of at least \(\mathrm m_L/2\).
\status{ESTABLISHED}

\noindent\emph{Proof.}
For a bridge \(e\) and \(\ell<\ell(e)\), \(r_e\) is independent of the rows of
\(B_1,\ldots,B_\ell\) (\Cref{lem:rg-formation-carrier}), so it does not vanish
on \(W_\ell\); it annihilates \(\mathcal H\oplus\mathcal K\), hence does not
vanish on \(\mathcal U_\ell\), which \(P_\ell\) fixes. An internal residual of
triangle \(q\) is a nonzero combination of the contrast rows
\(O_\triangle r_q\), which lie in \(\mathcal U\), plus a multiple of the loop
covector, so its restriction to \(\mathcal U\) is nonzero. Evaluated in an
orthonormal basis of \(\mathcal U\) on \((1,t,\ldots,t^{\dim\mathcal U-1})\),
each of the \(\mathrm n_0\) functionals is a nonzero polynomial of degree less
than \(\dim\mathcal U\), so some integer \(t\in[0,\mathrm n_0(\dim\mathcal U-1)]\)
is a simultaneous nonroot; normalizing gives \(\mathbf v_\star\), and
continuity gives the neighborhood. \(\square\)

\lemmaheading{Prepared continuation}{lem:formation-prepared-continuation}
For the preparation of \Cref{def:rg-formation-regime} and every finite
sequence of activations, the coupled diffusion, source proposals, and own-law
parameter equation have a unique nonexplosive solution through
\(t_*=2/\epsilon\), with positive smooth joint densities. On every active level
\(\ell\) and for all \(t\le t_*\),
\begin{align}
 \E_{q_t^\ell}y^\top H_\ell y
 &\le \mathrm C_L[\mathrm A^2\mathrm J_{\ell+1}+\mathrm T\varsigma_\epsilon^2],
 \label{eq:rg-formation-actual-energy}\\
 \E_{q_t^\ell}|y|^2
 &\le \mathrm C_L[\mathrm A^2+\mathrm T\varsigma_\epsilon^2],
 \qquad \mathrm J_1\E_{q_t^\ell}|h|^2\le \mathrm C_L\mathrm T,
 \label{eq:rg-formation-actual-moments}
\end{align}
and the parameter satisfies \(1\le\vartheta_\ell(t)\le21/10\) and
\(|\dot\vartheta_\ell(t)|\le\mathrm C_L\epsilon\); these bounds use no
conditional-entropy or history estimate. \status{ESTABLISHED}

\noindent\emph{Proof.}
For a fixed positive instance and a prescribed continuous parameter path near
\([1,21/10]\), each sheet drift has a bounded Hessian and linear growth by
\cref{eq:rg-formation-parent-relative-bounds}, and the proposal rows are
smooth with bounded derivatives and a positive lower bound, their
Gaussian-integral denominator being at least \(e^{-B_0}\); interlacing the
sheet diffusions on an independent rate-one Poisson grid gives a unique
nonexplosive process with all polynomial moments. For the parameter feedback
let \(\mathfrak u^\vartheta_j(t',\mathbf x)=\E_\vartheta[\mathsf f(\mathbf X_t,S_t)\mid
\mathbf X_{t'}=\mathbf x,S_{t'}=j]\), for smooth \(\mathsf f\) with polynomially
growing derivatives. Given the proposal times and labels, it composes
Lipschitz sheet flows with products of positive row weights; flow derivatives
are bounded on compact intervals, a differentiated weight is at most a
constant times the weight, and the proposal count has all exponential
moments, so \(\mathfrak u\) and \(\nabla\mathfrak u\) have polynomial bounds.
It\^o's formula for \(\mathfrak u^{\widetilde\vartheta}\) along the
\(\vartheta\) process, stopped in a ball and then using the moments, gives
\begin{equation}
 |\E_{\vartheta}\mathsf f(\mathbf X_t,S_t)
       -\E_{\widetilde\vartheta}\mathsf f(\mathbf X_t,S_t)|
 \le \mathrm C\int_0^t|\vartheta(t')-\widetilde\vartheta(t')| dt',
 \label{eq:rg-formation-local-law-dependence}
\end{equation}
since the second-order coefficients agree and the drift and row differences
are at most \(\mathrm C|\vartheta-\widetilde\vartheta|(1+|\mathbf x|)\) and
\(\mathrm C|\vartheta-\widetilde\vartheta|\). The observable terms of
\cref{eq:rg-formation-complete-score} meet these hypotheses, and their explicit
terminal parameter dependence adds a local Lipschitz term, so the integral map
for \(\dot\vartheta=-\epsilon\mathsf m(\vartheta)^{-1}g_\ell(q,\vartheta)\) is
a contraction on a short interval, which gives local coupled existence and
uniqueness. The step length is uniform: for parameter paths in
\([3/4,5/2]\), \(\vartheta-\delta_*\ge1/4\) and \(H_\ell\succeq\mathrm h_LI\)
make every non-root sheet drift uniformly confining toward zero, the root is a
fixed Ornstein--Uhlenbeck factor, proposals fix and projections contract the
configuration, and so the generator maps \(|\mathbf x|^{2k}\) below
\(\mathrm C_k-\mathrm c_k|\mathbf x|^{2k}\); all moments on \([0,t_*]\), hence
the constants of \cref{eq:rg-formation-local-law-dependence} and of the
covector, are uniform over such paths, and finitely many steps reach \(t_*\)
unless \(\vartheta\) leaves \([1,21/10]\), which the stopped argument below
excludes.

Stop at the first possible exit of \(\vartheta\) from \([1,21/10]\). The
Gaussian comparison
\(d\mathbf X^{\mathrm G}=-\vartheta(t)H_\ell\mathbf X^{\mathrm G}dt
+\sqrt{2\mathrm T} dW\), with the same parameter and Brownian increments on the
active level, starts from the reference of \Cref{def:rg-formation-regime}
coupled by \cref{eq:rg-formation-preparation-transport}; at an activation its
projected reference continues while the parent starts at the full retained
sample. For \(\mathbf e=\mathbf X-\mathbf X^{\mathrm G}\) the noise cancels and
proposals leave \(\mathbf e\) continuous, so
\cref{eq:rg-formation-parent-relative-bounds} gives, for every adapted label
history,
\begin{equation}
 \frac d{dt}|\mathbf e|^2\le-\mathbf e^\top H_\ell\mathbf e+\mathrm C_L\mathrm T,
 \qquad
 \sup_{t\le t_*}|\mathbf e_t|
 \le |\mathbf e_0|+\mathrm C_L\sqrt{\mathrm Tt_*}.
 \label{eq:rg-formation-path-force-bound}
\end{equation}
For the energy, variation of constants with the commuting propagator
\(\exp[-H_\ell\int_{t'}^t\vartheta]\) uses the first bound below, obtained by
splitting the integral at the inverse extreme eigenvalues, where the integrand
is at most \(\lambda_{\max}(H_\ell)\), \(1/(e\sigma)\), and
\(\lambda_{\min}(H_\ell)e^{-\lambda_{\min}(H_\ell)\sigma}\); Euclidean
projection and the Schur variational identity contract the two bounds at each
activation, and finitely many stages give the second:
\begin{equation}
 \int_0^\infty\|H_\ell e^{-H_\ell\sigma}\| d\sigma
 \le \mathrm C_L\left[1+\log
       \frac{\lambda_{\max}(H_\ell)}{\lambda_{\min}(H_\ell)}\right],
 \qquad
 \sup_t|H_\ell^{1/2}\mathbf e_t|\le \mathrm C_L\sqrt{\mathrm T}\varsigma_\epsilon.
 \label{eq:rg-formation-energy-force-bound}
\end{equation}
The Gaussian covariance barrier \(\Sigma^{\mathrm G}\preceq\mathrm C_L\mathrm TH_\ell^{-1}\)
needs only \(\vartheta\ge1\), by substitution into
\(\dot\Sigma^{\mathrm G}=-\vartheta(H_\ell\Sigma^{\mathrm G}+\Sigma^{\mathrm G}H_\ell)
+2\mathrm TI\), and survives retention because the retained block of
\(H_\ell^{-1}\) is the inverse Schur complement; the deterministic mean
contracts, has norm \(O_L(\mathrm A)\), and has zero loop part. This gives
\cref{eq:rg-formation-actual-energy,eq:rg-formation-actual-moments}, the loop
bound following from the core rate \(\vartheta\mathrm J_1\) and a force of at
most \(\mathrm C_L\sqrt{\mathrm T\mathrm J_1}\), whose forced error is at most
\(\mathrm C_L\sqrt{\mathrm T/\mathrm J_1}\).

Substituting these moments in the covector and metric bounds gives
\(\dot\vartheta=\epsilon(2-\vartheta)+O_L(\epsilon r_\epsilon)\), with
\(r_\epsilon\) of \cref{eq:rg-formation-parameter-remainder}, before any bound
on \(\dot\vartheta\) is used. The field points inward at both endpoints of
\([1,21/10]\) for small \(\epsilon\), which excludes the stopped exit; the
moments exclude explosion, so the solution reaches \(t_*\) with
\(|\dot\vartheta|\le\mathrm C_L\epsilon\). Positive-temperature elliptic
smoothing and positive finite rates keep densities positive and smooth, and
the construction restarts after an activation from the full retained law.
\hfill\(\square\)

\lemmaheading{Comparison of independently learned parameters}{lem:formation-parameter-comparison}
Every process or finite replacement lineage of
\Cref{lem:formation-prepared-continuation} satisfies
\cref{eq:rg-formation-parameter-bound} through \(t_*\), where every
independently evolving descendant uses its own law and copies only the
starting parameter. \status{ESTABLISHED}

\noindent\emph{Proof.}
The first two bounds are in \Cref{lem:formation-prepared-continuation}.
Subtracting \(\dot\vartheta_\circ=\epsilon(2-\vartheta_\circ)\),
\(\vartheta_\circ(0)=1\), from the equation of the process and integrating
with the factor \(e^{\epsilon t}\) bounds the forcing contribution by
\(\mathrm C_Lr_\epsilon\) through the horizon, where
\(r_\epsilon=\epsilon^{2\mathrm Q_1+20}(1+\epsilon^{\mathrm Q_1+6}\varsigma_\epsilon^2)\to0\)
by \cref{eq:rg-formation-parameter-remainder}, and copying the parameter at
finitely many activations adds a finite multiple of that error. The
comparison trajectory solves the calibration part of the own-law equation,
since the calibration record contributes \(\vartheta-2\) to the covector and
\(1/\mathrm T\) to \(\mathsf M\), and every law-dependent term is of order
\(\epsilon^{2\mathrm Q_1+20}\). As \(2-e^{-2}>7/4\), the coefficient exceeds
\(7/4\) on the last ten units of the horizon for small \(\epsilon\).
\hfill\(\square\)

\section{Entropy estimates}
\label{sec:rg-formation-entropy-proof}

\lemmaheading{Continuous entropy and score tools}{lem:formation-entropy-tools}
Let \(\mu(d\mathbf x)\propto e^{-\mathcal V(\mathbf x)/\mathrm T}d\mathbf x\) on
a finite-dimensional Euclidean space, where \(\mathcal V\) is smooth with a
globally bounded Hessian and \(\nabla^2\mathcal V\succeq\lambda I\),
\(\lambda>0\). For a positive density \(\mathsf f\) relative to \(\mu\), with
\(\operatorname{Ent}_\mu(\mathsf f)=\int\mathsf f\log(\mathsf f/\int\mathsf f d\mu) d\mu\),
for \(\mathfrak F\) with \(\operatorname{Lip}(\mathfrak F)\le\mathrm K\), for
every probability law \(\nu\), and whenever the right sides are finite,
\begin{align}
 \operatorname{Ent}_\mu(\mathsf f)
 &\le \frac{\mathrm T}{2\lambda}
           \int\frac{|\nabla\mathsf f|^2}{\mathsf f} d\mu,
 \label{eq:rg-formation-continuous-lsi}\\
 \log\E_\mu e^{\upsilon(\mathfrak F-\E_\mu\mathfrak F)}
 &\le\frac{\mathrm T\mathrm K^2\upsilon^2}{2\lambda},
 \qquad \upsilon\in\R,
 \label{eq:rg-formation-herbst-proof}\\
 \upsilon\E_\nu\mathfrak F&\le\KL(\nu\Vert\mu)+\log\E_\mu e^{\upsilon\mathfrak F},
 \label{eq:rg-formation-entropy-duality}
\end{align}
so \(|\E_\nu\mathfrak F-\E_\mu\mathfrak F|\le\mathrm K\sqrt{2\mathrm T\KL(\nu\Vert\mu)/\lambda}\).
A centered observable \(\xi\) that obeys the first line below obeys the second:
\begin{align}
 \log\E_\mu e^{\upsilon\xi}&\le
       \frac{\mathrm v\upsilon^2}{2(1-\mathrm b|\upsilon|)}\quad\text{for } |\upsilon|<\mathrm b^{-1},
 \label{eq:rg-formation-subgamma-proof}\\
 |\E_\nu\xi|
 &\le\sqrt{2\mathrm v\KL(\nu\Vert\mu)}+\mathrm b\KL(\nu\Vert\mu).
 \label{eq:rg-formation-score-duality}
\end{align}
Moreover \(\|\nu-\mu\|_{\mathrm{TV}}\le\sqrt{\KL(\nu\Vert\mu)/2}\). These
statements apply to every conditional sheet potential of the construction.
\status{ESTABLISHED}

\noindent\emph{Proof.}
For the semigroup \(P_t\) generated by
\(\mathrm T\Delta-\nabla\mathcal V\cdot\nabla\), the derivative of the
synchronous stochastic flow has operator norm at most \(e^{-\lambda t}\), so
weighted Cauchy--Schwarz gives
\(|\nabla P_t\mathsf f|^2/P_t\mathsf f\le e^{-2\lambda t}P_t(|\nabla\mathsf f|^2/\mathsf f)\)
and invariance of \(\mu\) gives
\(\mathcal I_\mu(P_t\mathsf f)\le e^{-2\lambda t}\mathcal I_\mu(\mathsf f)\) for
the Fisher information \(\mathcal I_\mu\). Integration by parts gives
\(d\operatorname{Ent}_\mu(P_t\mathsf f)/dt=-\mathrm T\mathcal I_\mu(P_t\mathsf f)\);
for bounded positive Lipschitz \(\mathsf f\), contraction against a stationary
initial sample gives \(P_t\mathsf f\to\int\mathsf f d\mu\) and, by bounded
convergence, entropy tending to zero, so integrating the Fisher bound proves
\cref{eq:rg-formation-continuous-lsi}. Cutoff, positive truncation, and short
semigroup smoothing extend it by lower semicontinuity, and Gaussian tails from
strong convexity justify the integrations. This is the curvature argument of
\citet[pp.~187, 194--200]{BakryEmery1985} with the diffusion coefficient made
explicit; its Gaussian case is \citet[Thm.~5.4]{BoucheronLugosiMassart2013}.
Applied to \(e^{\upsilon\mathfrak F}\), \cref{eq:rg-formation-continuous-lsi}
gives \(\operatorname{Ent}_\mu(e^{\upsilon\mathfrak F})\le
(\mathrm T\mathrm K^2\upsilon^2/(2\lambda))\E_\mu e^{\upsilon\mathfrak F}\), the
hypothesis of Herbst's argument \citep[Prop.~6.1]{BoucheronLugosiMassart2013}
with variance factor \(\mathrm T\mathrm K^2/\lambda\); applied to
\(\pm\mathfrak F\) it gives \cref{eq:rg-formation-herbst-proof}. The inequality
\cref{eq:rg-formation-entropy-duality} is the variational formula for relative
entropy \citep[Cor.~4.15]{BoucheronLugosiMassart2013}, and optimizing it over
\(\upsilon\) with \cref{eq:rg-formation-herbst-proof} gives the Lipschitz
consequence. Under \cref{eq:rg-formation-subgamma-proof}, the transportation
lemma \citep[Lemma~4.18]{BoucheronLugosiMassart2013} applied to \(\pm\xi\),
with the inverse convex conjugate \(\sqrt{2\mathrm vt}+\mathrm bt\) of the
sub-gamma function \citep[Sec.~2.4]{BoucheronLugosiMassart2013}, gives
\cref{eq:rg-formation-score-duality}; the case \(\KL=\infty\) is trivial. The
last inequality is Pinsker's inequality
\citep[Thm.~4.19]{BoucheronLugosiMassart2013}, with the constant \(1/2\) due to
\citet{Csiszar1967}. \hfill\(\square\)

\lemmaheading{Entropy at activations}{lem:formation-activation-entropy}
For the law of the active level \(\ell\), let
\(\mathfrak E_\ell(t)=\mathrm T\KL(q_t^\ell\Vert q^{\star,\ell}_{\vartheta_\ell(t)})\),
with the root factor canceled. At an activation time and on each ensuing
active interval,
\begin{align}
 \mathfrak E_\ell(t^{\mathrm{act}}_\ell)&\le \mathrm C_L[\mathrm A^2\mathrm J_{\ell+1}+\mathrm T\varsigma_\epsilon^2],
 \label{eq:rg-formation-activation-budget}\\
 \mathrm T^2\int_{t^{\mathrm{act}}_\ell}^{t}\mathcal I_\ell(\sigma) d\sigma
 &\le \mathrm C_L[\mathfrak E_\ell(t^{\mathrm{act}}_\ell)+\mathrm T],
 \label{eq:rg-formation-full-fisher-budget}
\end{align}
where \(\mathcal I_\ell\) is the full continuous relative Fisher information of
the joint law. The normalized parameter scores of the hidden Gaussian
conditionals and of the joint satisfy \cref{eq:rg-formation-subgamma-proof}
with constants \(\mathrm v_L,\mathrm b_L\) independent of the retained input,
of \(\mathrm T\), of the strengths, and of the source masses.
\status{ESTABLISHED}

\noindent\emph{Proof.}
The normalized hidden Gaussian conditionals have centers independent of
\(\vartheta\), and their logarithmic parameter scores, bounded marginal tilt
included, are centered versions of
\(-|\mathbf Z|^2/(2\vartheta)-\partial_\vartheta\psi_{\ell,j}\); the joint score
adds \(\partial_\vartheta\log\pi_j(\vartheta)\), of absolute value at most four.
A bounded tilt of a standard Gaussian satisfies
\begin{equation}
 \E|\mathbf Z|^{2m}
 \le e^{B_0}2^{\dim M_0/2}4^m m!,\qquad m\ge0,
 \label{eq:rg-formation-gaussian-score-moments}
\end{equation}
because \(r^m\le m!4^me^{r/4}\) and
\(\E e^{|\mathbf Z|^2/4}\le2^{\dim M_0/2}\), and the other score terms are
bounded by \cref{eq:rg-formation-tilt-score-proof}. Centering at most doubles
a geometric moment constant, and summing the exponential series without its
zero linear term, with \(\log(1+x)\le x\), gives
\cref{eq:rg-formation-subgamma-proof}, which conditional expectations keep by
Jensen's inequality.

For the differential and counting entropy
\(\mathcal S(q)=-\sum_j\int q_j\log q_j\) across changes of dimension, write
\(q_j=\lambda_jq_{\mid j}\), omitting terms of zero mass. The entropy--Fisher
bound \(\mathcal I(\mathsf f)\ge2\pi e \mathrm d e^{-2\mathcal S(\mathsf f)/\mathrm d}\)
of \citet{Stam1959}, equivalent to the Gaussian logarithmic Sobolev inequality
\citep{Carlen1991}, follows from the Gaussian case of
\cref{eq:rg-formation-continuous-lsi} by expanding both sides relative to
\(\mathcal N(0,\sigma^2I)\), integrating \(\int\mathbf x\cdot\nabla\mathsf f=-\mathrm d\),
canceling the second moments, and minimizing at
\(\sigma^2=\mathrm d/\mathcal I(\mathsf f)\); Jensen's inequality and
\(\mathcal S(q)=\sum_j\lambda_j\mathcal S(q_{\mid j})+\mathcal S(\lambda)\), with
\(\mathcal S(\lambda)\ge0\), extend it to
\(\sum_j\lambda_j\mathcal I(q_{\mid j})\ge2\pi e \mathrm d e^{-2\mathcal S(q)/\mathrm d}\).
At fixed \(\mathbf x\), a proposal moves the conditional label row
\(\mathfrak a\) in the direction \(\mathfrak a\beta^\ell-\mathfrak a\) and leaves
the continuous marginal unchanged, and concavity of categorical entropy bounds
the directional derivative below by
\(\mathcal S(\mathfrak a\beta^\ell)-\mathcal S(\mathfrak a)\ge-\log|\mathcal J|\).
On a \(\mathrm d\)-dimensional sheet, diffusion therefore gives the first line
below, and the two bounds together give the second:
\begin{align}
 \mathcal S'(q)_{\mathrm{diff}}
 &=\mathrm T\sum_j\lambda_j\mathcal I(q_{\mid j})-\E_q\Delta V_{\ell,S}
 \ge \mathrm T\sum_j\lambda_j\mathcal I(q_{\mid j})-\mathrm d\Lambda_\ell,
 \qquad \Lambda_\ell\le \mathrm C_L\mathrm J_1,
 \label{eq:rg-formation-actual-entropy-derivative}\\
 \mathcal S'(q)&\ge 2\pi e \mathrm d \mathrm T e^{-2\mathcal S(q)/\mathrm d}
                -\mathrm d\Lambda_\ell-\log|\mathcal J|.
 \label{eq:rg-formation-counting-entropy-floor}
\end{align}
Scalar comparison bounds \(\mathcal S(q)\) below by the lesser of its initial
value and
\(\frac{\mathrm d}2\log[2\pi e \mathrm d\mathrm T/(\mathrm d\Lambda_\ell+\log|\mathcal J|)]\).
At a projection dropping \(k\) coordinates \(z\), the chain rule and the
maximum-entropy Gaussian bound, applied conditionally on \((Y,S)\) and then
with Jensen's inequality, give
\begin{equation}
 \mathcal S(Y,S)\ge \mathcal S(\mathbf X,S)
       -\frac k2\log\frac{2\pi e \E|z-\bar z|^2}{k},
 \label{eq:rg-formation-entropy-projection}
\end{equation}
where \(\bar z\) can be the projected comparison mean, whose centered
covariance the moment proof of \Cref{lem:formation-prepared-continuation}
bounds by \(\mathrm C_L\mathrm T\varsigma_\epsilon^2H_\ell^{-1}\). Initial
entropies and all logarithms in these estimates are
\(O_L(\varsigma_\epsilon^2)\), because the covariance eigenvalues, \(\mathrm A\),
and \(\mathrm T\) are fixed powers of \(\epsilon\), the admitted initial density
ratios are bounded, and there are finitely many projections. Expanding
\(\mathrm T\KL(q\Vert\gamma_{\ell,\vartheta}\otimes\pi(\vartheta))\) and using
\cref{eq:rg-formation-actual-energy} proves
\cref{eq:rg-formation-activation-budget}; the bounded optimizer tilt adds at
most \(\mathrm TB_0\), and the source prior has a positive minimum on
\([1,21/10]\), so its logarithmic contribution does not depend on the masses
\(\lambda_j\).

For the time-integrated estimate put \(\mathcal V=y^\top H_\ell y/2\) and
\(\mathrm R_L=e^{B_0}2^{\dim M_0/2}\). The reference satisfies the first bound
below. Entropy duality gives
\(\vartheta\E_q\mathcal V\le2\mathfrak E_\ell+2\mathrm T\log\mathrm R_L\),
Jensen's inequality gives
\(\vartheta\E_{q^{\star,\ell}_\vartheta}\mathcal V\le2\mathrm T\log\mathrm R_L\),
and the other centered terms of the normalized full score are at most
\(2\mathrm C_{\mathrm{sc},L}\), which gives the second. Diffusion dissipates
\(\mathrm T^2\mathcal I_\ell\) and detailed balance makes the source
contribution nonpositive, which gives the third for a fixed
\(\mathrm C_{0,L}\):
\begin{align}
 \E_{q^{\star,\ell}_\vartheta}\exp\frac{\vartheta\mathcal V}{2\mathrm T}
 &\le \mathrm R_L,
 \label{eq:rg-formation-energy-exponential-moment}\\
 |\mathrm T\E_q\partial_\vartheta\log q^{\star,\ell}_\vartheta|
 &\le\frac{2\mathfrak E_\ell+4\mathrm T\log\mathrm R_L}{\vartheta}+2\mathrm T\mathrm C_{\mathrm{sc},L},
 \label{eq:rg-formation-entropy-parameter-budget}\\
 \mathfrak E_\ell'&\le -\mathrm T^2\mathcal I_\ell+
    \frac{2|\dot\vartheta|}{\vartheta}\bigl(\mathfrak E_\ell+\mathrm C_{0,L}\mathrm T\bigr).
 \label{eq:rg-formation-entropy-energy-inequality}
\end{align}
The integrating factor \(\exp[-2\int|\dot\vartheta|/\vartheta]\) has exponent
bounded over \(t_*\) by the speed bound, and integration proves
\cref{eq:rg-formation-full-fisher-budget}. At an activation the law and the
reference are pushed through the same channel at the same parameter, so the
chain rule makes \(\mathfrak E\) decrease along the finite lineage. All
entropy derivatives are first taken for positive truncated densities on
bounded domains, and the moment and entropy bounds remove the cutoffs.
\hfill\(\square\)

\lemmaheading{Conditional entropy and complete score error}{lem:formation-conditional-entropy}
There is a constant \(\mathrm c_{\mathrm{early},L}>0\) such that, for every
retention \(\ell-1\to\ell\), the conditional relative entropy \(\Delta_\ell\)
with respect to the optimizing hidden conditional satisfies
\begin{equation}
 \sup_{t\in I_\ell\text{ and later active times}}
 \Delta_\ell(t)
 \le
 \begin{cases}
 \mathrm C_L\epsilon^{14},&\ell<L,\\
 \mathrm C_L\epsilon^2,&\ell=L,
 \end{cases}
 \label{eq:rg-formation-conditional-final-bounds}
\end{equation}
where the early bound also applies, until the next activation, to the
corresponding unreplaced comparison process, and the terminal bound is
\Cref{lem:formation-terminal-joint-entropy}. At every such time the complete
fine/parent covector difference is at most
\(\mathrm C_L\mathrm T(\sqrt{\Delta_\ell}+\Delta_\ell)\). \status{ESTABLISHED}

\noindent\emph{Proof of the early bound and the covector estimate.}
For an early retention use the auxiliary continuous split of the first line
below. The true parent retains the loops, so the chain rule bounds its hidden
conditional defect \(\Delta_\ell\) by the auxiliary defect
\(\widehat\Delta_\ell\), which also integrates \(h\); the independent root
cancels, and conditional strong convexity in \(Z\) is at least
\(\lambda_\ell=\mathrm c_{\mathrm{early},L}\mathrm J_\ell\) after reducing the
constant if necessary. The mixed Hessian obeys the second line:
\begin{align}
 &Z=(\mathcal U_{\ell-1}\cap\mathcal U_\ell^\perp)\oplus\mathcal H,
 \qquad Y=\mathcal U_\ell,
 \qquad S\ \text{retained},
 \label{eq:rg-formation-auxiliary-split}\\
 &\|\nabla_Z\nabla_YV_{\ell-1,j}\|\le\Gamma_\ell=\mathrm C_L\mathrm J_{\ell+1},
 \qquad \Gamma_L=0.
 \label{eq:rg-formation-mixed-hessian-scale}
\end{align}
The nonlinear residuals lie in the first normal layer and the loop space, and
at \(\ell-1=0\) the nonlinear potential has no future normal dependence. For
\(\ell-1\ge1\), in coordinates \(r_E=B_Ew_E\) on the eliminated normal rows,
the Gaussian conditional precision
\(\Omega_E=\mathrm J_E+\mathcal C^\top\mathrm J_{\mathrm{fut}}\mathcal C\), with
the diagonal strengths \(\mathrm J_E\) and \(\mathrm J_{\mathrm{fut}}\) of the
eliminated and future rows, has \(\|\Omega_E-\mathrm J_E\|\le\mathrm C_L\mathrm J_\ell\)
and \(\|\Omega_E^{-1}\|\le\mathrm C_L/\mathrm J_{\ell-1}\), so the identity in the
first line below bounds the first block row of \(\Omega_E^{-1}\) by
\(\mathrm C_L/\mathrm J_1\). A future force column of level \(r\) has size
\(\mathrm C_L\mathrm J_r\), so the first conditional mean responds by
\(\mathrm C_L\mathrm J_r/\mathrm J_1\), with conditional covariance comparable to
\(\mathrm T/\mathrm J_1\). In the derivatives of the normalized Gaussian
integral defining \(\psi_{\ell-1,j}\), a first-layer or loop derivative of
\(\phi_j\) has size \(\mathrm C_L\sqrt{\mathrm J_1/\mathrm T}\) and a second
derivative size \(\mathrm C_L\mathrm J_1/\mathrm T\), each future derivative
brings that mean response, and the covariance term is bounded by products of
first-derivative bounds, which gives the second line:
\begin{align}
 \Omega_E^{-1}&=\mathrm J_E^{-1}
           -\mathrm J_E^{-1}(\Omega_E-\mathrm J_E)\Omega_E^{-1},
 \label{eq:rg-formation-first-row-resolvent}\\
 \mathrm T|\partial_h\partial_{w_r}\psi_{\ell-1,j}|&\le \mathrm C_L\mathrm J_r,\qquad
 \mathrm T|\partial_{w_i}\partial_{w_r}\psi_{\ell-1,j}|
       \le \mathrm C_L\frac{\mathrm J_i\mathrm J_r}{\mathrm J_1}.
 \label{eq:rg-formation-future-derivative-scales}
\end{align}
The Gaussian Schur cross block is at most \(\mathrm C_L\mathrm J_{\ell+1}\) by
the lower triangular residual form and its variational bound, and summing
finitely many blocks, generated conditional covariance included, proves
\cref{eq:rg-formation-mixed-hessian-scale}.

Write \(q=q_Y(y,j)q_{Z\mid Y}(z\mid y,j)\) and
\(q^\star=q^\star_Y(y,j)q^\star_{Z\mid Y}(z\mid y,j)\), and let
\(\mathfrak w_Y=\nabla_Y\log(q/q^\star)\),
\(\mathfrak v_Y=\nabla_Y\log(q_Y/q^\star_Y)\), and
\(\mathfrak e_Y=\E_{q_{Z\mid Y}}\nabla_Y\log q^\star_{Z\mid Y}\). Normalization
gives \(\E_{q^\star_{Z\mid Y}}\nabla_Y\log q^\star_{Z\mid Y}=0\) and
\(\E_{q_{Z\mid Y}}\mathfrak w_Y=\mathfrak v_Y-\mathfrak e_Y\), the marginal
flux is \(\mathrm Tq_Y(\mathfrak v_Y-\mathfrak e_Y)\), and subtracting the
marginal entropy derivative from the full one yields
\begin{align}
 \widehat\Delta_\ell'|_{\mathrm{cont}}
 ={}&-\mathrm T\mathcal I_Z
      -\mathrm T\E_{q_Y}\operatorname{Var}_{q_{Z\mid Y}}(\mathfrak w_Y)
      +\mathrm T\E_{q_Y}\bigl[\mathfrak e_Y\cdot\E_{q_{Z\mid Y}}\mathfrak w_Y\bigr]
 \notag\\
 &-\dot\vartheta \E_q\partial_\vartheta\log q^\star_{Z\mid Y}.
 \label{eq:rg-formation-conditional-balance}
\end{align}
The conditional logarithmic Sobolev inequality gives
\(\mathrm T\mathcal I_Z\ge2\lambda_\ell\widehat\Delta_\ell\). Applying
\cref{eq:rg-formation-herbst-proof} to each unit-vector component of
\(\nabla_Y\log q^\star_{Z\mid Y}\), whose Lipschitz constant in \(Z\) is at most
\(\Gamma_\ell/\mathrm T\), gives
\(|\mathfrak e_Y|\le\Gamma_\ell\sqrt{2\KL(q_{Z\mid Y}\Vert q^\star_{Z\mid Y})/(\mathrm T\lambda_\ell)}\),
so Cauchy--Schwarz and Young's inequality bound the mixed term by
\(\lambda_\ell\widehat\Delta_\ell/2+\mathrm T\Gamma_\ell^2\mathcal I_Y/\lambda_\ell^2\),
with \(\mathcal I_Y\) the \(Y\)-component of the full relative Fisher
information. The normalized conditional parameter score, sub-gamma by
\Cref{lem:formation-activation-entropy}, contributes at most
\(\lambda_\ell\widehat\Delta_\ell/2+\mathrm v_L|\dot\vartheta|^2/\lambda_\ell
+\mathrm b_L|\dot\vartheta|\widehat\Delta_\ell\).

For the source updates, at fixed \(Y=y\) each
\(\log(\gamma^{\mathrm{hid}}/q^\star_{Z\mid Y,j})\), with \(\gamma^{\mathrm{hid}}\)
the common label-independent hidden Gaussian, lies in \([-B_0,B_0]\), and the
conditional entropy relative to \(\gamma^{\mathrm{hid}}\) splits as
\(\KL(q_Z\Vert\gamma^{\mathrm{hid}})+\mathbb I(Z;S)\). A proposal leaves \(q_Z\)
unchanged, and the mutual information is convex in the conditional label
channel at fixed \(q_Z\) by joint convexity of relative entropy, so its
directional derivative toward the new channel is at most \(\log|\mathcal J|\);
the difference from the optimizing target, a linear expectation of
\(\log(\gamma^{\mathrm{hid}}/q^\star_{Z\mid Y,j})\), changes at rate at most
\(2B_0\). This divides by no source mass and extends to vanishing masses by
convex approximation. With \(\mathrm b_L|\dot\vartheta|\le\lambda_\ell/2\) for
small \(\epsilon\) and the negative variance term dropped,
\begin{align}
 \widehat\Delta_\ell'|_{\mathrm{jump}}&\le
 \mathrm C_{\mathcal J}:=\log|\mathcal J|+2B_0,
 \label{eq:rg-formation-conditional-jump-remainder}\\
 \widehat\Delta_\ell'
 &\le-\frac{\lambda_\ell}{2}\widehat\Delta_\ell
     +\frac{\mathrm T\Gamma_\ell^2}{\lambda_\ell^2}\mathcal I_Y
     +\frac{\mathrm v_L|\dot\vartheta|^2}{\lambda_\ell}
     +\mathrm C_{\mathcal J}.
 \label{eq:rg-formation-conditional-differential-bound}
\end{align}
The initial conditional entropy is at most
\(\mathfrak E_{\ell-1}(t^{\mathrm{act}}_{\ell-1})/\mathrm T\). Integrating with
\cref{eq:rg-formation-full-fisher-budget,eq:rg-formation-activation-budget}
and \(\mathrm A^2/\mathrm T=\epsilon^{-6}\) gives
\begin{align}
 \widehat\Delta_\ell(t)\le \mathrm C_L\bigg[
 &e^{-\mathrm c_{\mathrm{early},L}\mathrm J_\ell(t-t^{\mathrm{act}}_{\ell-1})}
                 (\epsilon^{-6}\mathrm J_\ell+\varsigma_\epsilon^2)
 \notag\\
 &+\epsilon^{-6}\frac{\mathrm J_{\ell+1}^2}{\mathrm J_\ell}
 +\varsigma_\epsilon^2\frac{\mathrm J_{\ell+1}^2}{\mathrm J_\ell^2}
 +\frac{\epsilon^2}{\mathrm J_\ell^2}+\frac{\mathrm C_{\mathcal J}}{\mathrm J_\ell}\bigg],
 \label{eq:rg-formation-early-entropy-bound}
\end{align}
with \(\mathrm c_{\mathrm{early},L}\) absorbing the factor \(1/2\). For
\(\ell<L\), \(\mathrm J_{\ell+1}^2/\mathrm J_\ell=\epsilon^{20}\) and
\(\mathrm J_\ell\ge\epsilon^{-20}\); on \(I_\ell\) the entrance time and
\(\mathrm c_{\mathrm{early},L}k_L/2>\mathrm Q_1+20\) bound the first term by
\(\mathrm C_L\epsilon^{14}\), and the other terms meet the same bound; at
\(\ell=L\) the last term would be \(\mathrm C_{\mathcal J}\).

For the true retention, \cref{eq:rg-formation-conditional-cost} with
\(\ell'=\ell-1\), differentiated at fixed recognition law and channel, gives
\(g_{\ell-1}-g_\ell=-\mathrm T\E_q\partial_\vartheta\log q^\star_{Z\mid Y,\mathrm{true}}\),
the common scalar and the source-prior derivative canceling. The normalized
conditional score has the same sub-gamma bound, so
\cref{eq:rg-formation-score-duality}, applied conditionally, with Jensen's
inequality gives the covector bound
\(|g_{\ell-1}-g_\ell|\le\mathrm C_L\mathrm T(\sqrt{\Delta_\ell}+\Delta_\ell)\).
\(\square\)

\lemmaheading{Joint relaxation at the weakest strength}{lem:formation-terminal-joint-entropy}
There is a constant \(\mathrm c_{\mathrm{joint},L}>0\), independent of
\(\epsilon\), such that every lineage satisfies
\begin{equation}
 \KL\bigl(q_t^\ell\Vert q^{\star,\ell}_{\vartheta_\ell(t)}\bigr)
 \le \mathrm C_L e^{-\mathrm c_{\mathrm{joint},L}t}
           [\epsilon^{-(\mathrm Q_1+6)}+\varsigma_\epsilon^2]
       +\mathrm C_L\epsilon^2.
 \label{eq:rg-formation-terminal-joint-bound}
\end{equation}
Consequently the terminal part of
\cref{eq:rg-formation-conditional-final-bounds} holds once
\(t\ge t_L=k_L\log(1/\epsilon)\). \status{ESTABLISHED}

\noindent\emph{Proof.}
Fix a level and a parameter, write \(q^\star=q^{\star,\ell}_\vartheta\) and
\(\mathfrak r_j=q_j/q^\star_j\). The bounded tilt gives
\(e^{-B_0}\le q^\star/(\gamma_{\ell,\vartheta}\otimes\pi)\le e^{B_0}\), and
\cref{eq:rg-formation-parent-generator} gives the conductance comparison in the
first line below. For a probability measure \(\mu\),
\(\operatorname{Ent}_\mu(\mathfrak r)=\inf_{r>0}\int[\mathfrak r\log(\mathfrak r/r)-\mathfrak r+r] d\mu\)
with a nonnegative integrand, so
\(\operatorname{Ent}_{q^\star}(\mathfrak r)\le e^{B_0}\operatorname{Ent}_{\gamma\otimes\pi}(\mathfrak r)\),
and the product chain rule and convexity give the second line, whose first
term the Gaussian entropy inequality bounds by
\(\mathrm T/(2\vartheta\mathrm h_L)\) times its continuous Fisher information:
\begin{align}
 q^\star_j(y)\beta_{jk}^\ell(y)
 &\ge e^{-2B_0}\gamma_{\ell,\vartheta}(y)\pi_j(\vartheta)\pi_k(\vartheta),
 \label{eq:rg-formation-conductance-comparison}\\
 \operatorname{Ent}_{\gamma\otimes\pi}(\mathfrak r)
 &\le\sum_j\pi_j\operatorname{Ent}_\gamma(\mathfrak r_j)
   +\E_\gamma\operatorname{Ent}_{\pi}(\mathfrak r(y)).
 \label{eq:rg-formation-product-entropy-bound}
\end{align}
For any positive finite vector \(\mathbf a\), Jensen's inequality for the
logarithm gives
\begin{equation}
 \frac12\sum_{j,k}\pi_j\pi_k
       (\mathbf a_j-\mathbf a_k)(\log\mathbf a_j-\log\mathbf a_k)
 =\sum_j\pi_j\mathbf a_j\log\mathbf a_j
       -\bar{\mathbf a}\sum_j\pi_j\log\mathbf a_j
 \ge \operatorname{Ent}_\pi(\mathbf a),
 \label{eq:rg-formation-categorical-entropy-bound}
\end{equation}
with \(\bar{\mathbf a}=\sum_j\pi_j\mathbf a_j\). Comparing the Fisher measures
and the conductances through their pointwise bounds, with the jump
dissipation of the first line below, gives the second:
\begin{align}
 \mathcal D_{\mathrm{jump}}
 &=\frac12\int\sum_{j,k}q^\star_j\beta_{jk}^\ell
           (\mathfrak r_j-\mathfrak r_k)(\log\mathfrak r_j-\log\mathfrak r_k) dy,
 \label{eq:rg-formation-jump-dissipation}\\
 \KL(q\Vert q^\star)&\le
 e^{2B_0}\frac{\mathrm T}{2\vartheta\mathrm h_L}\mathcal I_\ell
       +e^{3B_0}\mathcal D_{\mathrm{jump}}.
 \label{eq:rg-formation-joint-coercivity}
\end{align}
With \(\mathrm c_{0,L}=\min\{2\mathrm h_Le^{-2B_0},e^{-3B_0}\}>0\),
\(\mathrm T\mathcal I_\ell+\mathcal D_{\mathrm{jump}}\ge\mathrm c_{0,L}\KL\). The time-dependent normalized
full score satisfies \cref{eq:rg-formation-subgamma-proof}, hence, for small
\(\epsilon\),
\begin{equation}
 \KL'\le-\mathrm c_{0,L}\KL+
       \mathrm C_L\epsilon(\sqrt{\KL}+\KL)
 \le-\frac{\mathrm c_{0,L}}{2}\KL+\mathrm C_L\epsilon^2,
 \label{eq:rg-formation-joint-moving-reference}
\end{equation}
and \(\mathrm c_{\mathrm{joint},L}=\mathrm c_{0,L}/2\). Activations decrease the
relative entropy at the unchanged parameter, and each descendant then obeys
the same inequality with its own parameter speed. The initial relative entropy
is at most \(\mathrm C_L[\epsilon^{-(\mathrm Q_1+6)}+\varsigma_\epsilon^2]\), and
integration through finitely many decreases proves
\cref{eq:rg-formation-terminal-joint-bound}; the condition
\(\mathrm c_{\mathrm{joint},L}k_L>\mathrm Q_1+8\) makes its decaying term
\(O_L(\epsilon^2)\) at \(t_L\). Both the auxiliary and the true hidden
conditional entropies are bounded by the full joint relative entropy, which
proves the terminal assertion. \hfill\(\square\)

\section{Exact fast marginals and complete source histories}
\label{sec:rg-formation-history-proof}

\lemmaheading{Common fast marginal and environmental error}{lem:formation-fast-marginal}
In the flag coordinates \(w=(w_1,\ldots,w_L)\) of
\cref{eq:rg-formation-weighted-flag} and the loop coordinate \(h\), put
\(\Omega_*=\mathrm J_1\mathsf B_{11}^\top\mathsf B_{11}\) and
\begin{equation}
 \mu_{\vartheta,j}^*(dw_1,dh)
 =\frac{1}{Z_{\vartheta,j}^*}
   \exp\left\{-\frac{\vartheta}{2\mathrm T}
         (w_1^\top\Omega_*w_1+\mathrm J_1|h|^2)-\phi_j(w_1,h)\right\}dw_1 dh,
 \label{eq:rg-formation-fast-law}
\end{equation}
with loop marginal \(\mu_{\vartheta,j}^{*,h}\), first-normal conditional
\(\mu_{\vartheta,j}^{*,1}(dw_1\mid h)\), and, with environment scales
\(\varrho_\epsilon\), \(\varrho'_\epsilon\),
\begin{align}
 \beta_{jk}^{*,h}(h,\vartheta)
 &=\int\beta_k(w_1,h,\vartheta)
                   \mu_{\vartheta,j}^{*,1}(dw_1\mid h),
 \label{eq:rg-formation-fast-loop-row}\\
 \bar\beta_{jk}(\vartheta)
 &=\int\beta_k(w_1,h,\vartheta) \mu_{\vartheta,j}^*(dw_1,dh)
   =\int\beta_{jk}^{*,h}(h,\vartheta)
                    \mu_{\vartheta,j}^{*,h}(dh),
 \label{eq:rg-formation-common-averaged-row}\\
 \varrho_\epsilon&=\frac{\mathrm J_2}{\mathrm J_1},\qquad
 \varrho'_\epsilon=\frac{\mathrm J_2}{\sqrt{\mathrm J_1\mathrm T}},\qquad
 \mathrm e_\epsilon(y)=\varrho_\epsilon+\varrho'_\epsilon|y|.
 \label{eq:rg-formation-environment-scales}
\end{align}
Then \cref{eq:rg-formation-fast-law} is the exact \((w_1,h)\)-marginal of the
optimizer conditional on \(S=j\); averaging the inherited row
\(\beta^\ell_{jk}\) against the optimizer conditional on \(j\) gives the same
\(\bar\beta_{jk}(\vartheta)\) for every \(\ell\); and the exact parent rows and
loop forces satisfy
\cref{eq:rg-formation-parent-row-environment,eq:rg-formation-parent-force-environment,eq:rg-formation-row-fast-lipschitz}
below. \status{ESTABLISHED}

\noindent\emph{Proof.}
In the flag coordinates
\begin{equation}
 w^\top H_{\mathcal U}w
 =\mathrm J_1|\mathsf B_{11}w_1|^2
       +\sum_{\ell=2}^L\mathrm J_\ell
             \left|\sum_{r\le\ell}\mathsf B_{\ell r}w_r\right|^2,
 \label{eq:rg-formation-triangular-fast-quadratic}
\end{equation}
and \(\phi_j\) depends only on \((w_1,h)\). For fixed \(w_1,h\), change the
later \(w_\ell\) to the residuals \(r_\ell=\sum_{r\le\ell}\mathsf B_{\ell r}w_r\);
the domain stays the whole Euclidean product and the Jacobian is the constant
\(\prod_{\ell\ge2}|\det\mathsf B_{\ell\ell}|\), so integrating those Gaussian
residuals contributes a scalar independent of \(w_1,h,j\), which proves the
first claim, the prior factor \(\pi_j(\vartheta)\) canceling on conditioning.
The change of variables and the tower property give the second.

For an intermediate \(\ell\ge1\), keep \(y=(w_{\ell+1},\ldots,w_L)\) fixed and
integrate \(w_2,\ldots,w_\ell\). The unperturbed conditional first-normal law
\(\gamma^1_{\ell,\vartheta,y}\) and its divergence from
\(\gamma^{1,*}_\vartheta=\mathcal N(0,(\mathrm T/\vartheta)\Omega_*^{-1})\) are
\begin{align}
 &\gamma^1_{\ell,\vartheta,y}
 =\mathcal N\left(\mu_\ell(y),\frac{\mathrm T}{\vartheta}(\Omega^{(1)}_\ell)^{-1}\right),\qquad
 \Omega_*\preceq\Omega^{(1)}_\ell\preceq\Omega_*+\mathrm C_L\mathrm J_2I,\qquad
 |\mu_\ell(y)|\le \mathrm C_L\frac{\mathrm J_2}{\mathrm J_1}|y|,
 \label{eq:rg-formation-first-normal-conditional}\\
 &\frac12\left[
  \frac{\vartheta}{\mathrm T}\mu_\ell^\top\Omega_*\mu_\ell
  +\Tr\bigl(\Omega_*(\Omega^{(1)}_\ell)^{-1}\bigr)-\dim w_1
            -\log\det\bigl(\Omega_*(\Omega^{(1)}_\ell)^{-1}\bigr)\right]
 \le \mathrm C_L[\varrho_\epsilon^2+\varrho_\epsilon'^2|y|^2].
 \label{eq:rg-formation-first-normal-gaussian-kl}
\end{align}
For the lower precision bound, minimize the nonnegative later residual squares
at \(y=0\); for the upper, set the eliminated later coordinates to zero and use
\(\mathrm J_\ell\le\mathrm J_2\) for \(\ell\ge2\); at \(\ell=L\),
\(\Omega^{(1)}_L=\Omega_*\) and \(\mu_L=0\). The mean bound is the
first-block-row estimate \cref{eq:rg-formation-first-row-resolvent}, each
future coordinate of level \(\ell\) contributing at most
\(\mathrm C_L\mathrm J_\ell/\mathrm J_1\). The divergence bound holds because
the whitened covariance eigenvalues lie between
\((1+\mathrm C_L\varrho_\epsilon)^{-1}\) and one, where
\(t-1-\log t=O((t-1)^2)\), and Pinsker's inequality turns it into a total
variation at most \(\mathrm C_L\mathrm e_\epsilon(y)\).

Tilting two probability measures by one weight \(w\in[e^{-B_0},1]\) multiplies
their total-variation distance by at most \(2e^{B_0}\), since each term of
\(w\mu/\mu(w)-w\nu/\nu(w)=w(\mu-\nu)/\mu(w)+w\nu[1/\mu(w)-1/\nu(w)]\) has
total variation at most \(e^{B_0}\|\mu-\nu\|_{\mathrm{TV}}\). Applied to
\(w=e^{-\phi_j(w_1,h)}\) uniformly in \(h\), this gives the first bound below;
the second follows through the common Gaussian force
\(\vartheta\mathrm J_1h\) and \(\mathrm T\E\nabla_h\phi_j\), and the third
because \(\gamma^{1,*}_\vartheta\) does not depend on \(h\), so differentiation
under its normalized tilt gives
\[
\nabla_h\beta^{*,h}_{jk}
=\E_{\mu^{*,1}_{\vartheta,j}(\cdot\mid h)}\nabla_h\beta_k
-\operatorname{Cov}_{\mu^{*,1}_{\vartheta,j}(\cdot\mid h)}
 (\beta_k,\nabla_h\phi_j),
\]
and both terms have the asserted Lipschitz bound: the standardized source
potential has bounded first derivatives and \(\sum_k\beta_k=1\). Thus
\begin{align}
 \|\beta_{j\cdot}^\ell(y,h,\vartheta)
          -\beta_{j\cdot}^{*,h}(h,\vartheta)\|_{\mathrm{TV}}
 &\le \mathrm C_L\mathrm e_\epsilon(y),
 \label{eq:rg-formation-parent-row-environment}\\
 |\nabla_hV_{\ell,j}(y,h,\vartheta)
           -\nabla_hV_j^{*,h}(h,\vartheta)|
 &\le \mathrm C_L\sqrt{\mathrm T\mathrm J_1} \mathrm e_\epsilon(y),
 \label{eq:rg-formation-parent-force-environment}\\
 \operatorname{Lip}_{(w_1,h)}\beta_{j\cdot}
       +\operatorname{Lip}_h\beta_{j\cdot}^{*,h}
 &\le \mathrm C_L\sqrt{\mathrm J_1/\mathrm T}.
 \label{eq:rg-formation-row-fast-lipschitz}
\end{align}
\hfill\(\square\)

\lemmaheading{Fast mixing with a learned coefficient}{lem:formation-fast-mixing}
There is a constant \(\mathrm c_{\mathrm{fast},L}>0\) with the following property.
Let the auxiliary process have, on a proposal-free interval, fixed assignment
\(j\) and either the full fast potential of \cref{eq:rg-formation-fast-law} or
its exact loop marginal, and a deterministic coefficient with
\(1\le\vartheta\le21/10\) and \(|\dot\vartheta|\le\mathrm C_L\epsilon\). For a
point start \(\xi\) at time \(t'\) and \(t-t'\ge(\mathrm c_{\mathrm{fast},L}\mathrm J_1)^{-1}\),
\begin{equation}
 \|P^*_{t',t}(\xi,\cdot)-\mu^*_{\vartheta(t),j}\|_{\mathrm{TV}}
 \le \mathrm C_L\bigl(1+\sqrt{\mathrm J_1/\mathrm T}|\xi|\bigr)
              e^{-\mathrm c_{\mathrm{fast},L}\mathrm J_1(t-t')}
       +\mathrm C_L\epsilon/\mathrm J_1,
 \label{eq:rg-formation-fast-mixing-bound}
\end{equation}
where \(\mu^*\) means its loop marginal in the loop case. Also
\(\bar\beta_{jk}(\vartheta)\) is uniformly Lipschitz in \(\vartheta\).
\status{ESTABLISHED}

\noindent\emph{Proof.}
The relative Hessian bounds and the normalized marginal covariance calculation
give \(\mathrm c_L\mathrm J_1I\preceq\nabla^2V^*\preceq\mathrm C_L\mathrm J_1I\)
for both potentials, which are even, minimize at zero, and have equilibrium
second moments at most \(\mathrm C_L\mathrm T/\mathrm J_1\). For two point
starts \(\xi,\xi'\), reflection coupling before meeting gives their distance a
one-dimensional noise of size \(2\sqrt{2\mathrm T}\) and a drift at most
\(-\lambda\) times the distance, \(\lambda=\mathrm c_L\mathrm J_1\); multiplying
by \(e^{\lambda(t-t')}\) gives a time-changed Brownian motion of variance
\(4\mathrm T(e^{2\lambda(t-t')}-1)/\lambda\), and the reflection principle
bounds the probability of not meeting by
\begin{equation}
 \mathrm C|\xi-\xi'|\sqrt{\lambda/\mathrm T}
                 \bigl(e^{2\lambda(t-t')}-1\bigr)^{-1/2}.
 \label{eq:rg-formation-reflection-bound}
\end{equation}
The coupling reflects one Brownian motion orthogonally and predictably, so
both marginals have the required law; the drifts are globally Lipschitz for
each instance, so the copies are constructed away from the diagonal and joined
at meeting, the comparison and hitting bound following by stopping before the
diagonal and letting the stopping radius tend to zero. Averaging \(\xi'\)
against the equilibrium at time \(t'\) gives the first term of
\cref{eq:rg-formation-fast-mixing-bound}.

For the second, compare that evolving equilibrium start with the
instantaneous equilibrium at time \(t\). After standardizing \((w_1,h)\) by
\(\sqrt{\mathrm J_1/\mathrm T}\), the normalized parameter score is a centered
quadratic with uniform exponential moments, and the loop-marginal score, its
conditional expectation, has the same bound by Jensen's inequality. The
moving-reference entropy calculation with the continuous entropy inequality at
curvature \(\mathrm c_L\mathrm J_1\) gives, for the relative entropy
\(\mathrm D_{\mathrm{lag}}\) of the evolved law to the current equilibrium,
\begin{equation}
 \mathrm D_{\mathrm{lag}}'\le-\mathrm c_L\mathrm J_1\mathrm D_{\mathrm{lag}}
 +\mathrm C_L\epsilon^2/\mathrm J_1,\qquad \mathrm D_{\mathrm{lag}}(t')=0,
 \label{eq:rg-formation-fast-equilibrium-lag}
\end{equation}
the parameter term linear in \(\mathrm D_{\mathrm{lag}}\) being absorbed for
small \(\epsilon\); hence \(\mathrm D_{\mathrm{lag}}\le\mathrm C_L\epsilon^2/\mathrm J_1^2\),
and Pinsker's inequality gives the second term. Finally, differentiating
\(\bar\beta_{jk}=\E_{\mu_{\vartheta,j}^*}\beta_k\) gives the explicit row
derivative, whose absolute values sum to at most eight, and a covariance with
the centered scale score, whose standardized first moment is bounded, so
\(\sum_k|\partial_\vartheta\bar\beta_{jk}|\le\mathrm C_L\). \hfill\(\square\)

\lemmaheading{Coupling of complete source histories}{lem:formation-history-coupling}
The source-assignment histories of a process and of an activated exact parent
can be coupled, together with their common initial marked sample, so that the
probability of any disagreement through \(t_*\) is at most
\(\mathrm C_L\epsilon^6\). This holds through the finite activation lineage,
for the full finite-state path with arbitrary label correlations and the
rate-one proposal clock. \status{ESTABLISHED}

\noindent\emph{Proof.}
First bound a proposal row on one proposal-free interval. For a microscopic
process let \(\xi=(w_1,h)\), and couple its auxiliary fast diffusion to the
same initial fast point, Brownian increments, assignment, and parameter path.
The omitted Gaussian environmental force
\(\vartheta[(H_{11}-\Omega_*)w_1+H_{1,>1}w_{>1}]\) has norm at most
\(\mathrm C_L\mathrm J_2|w|\), and strong contraction from the most recent
proposal or activation time \(t_{\mathrm{prev}}\) gives the first bound below.
For a parent at level \(\ell\ge1\), coupling its retained loop to the
auxiliary loop process, \cref{eq:rg-formation-parent-force-environment} gives
the second:
\begin{align}
 |\xi_t-\xi_t^*|
 &\le \mathrm C_L\mathrm J_2\int_{t_{\mathrm{prev}}}^t
             e^{-\mathrm c_L\mathrm J_1(t-\sigma)}|w_\sigma| d\sigma,
 \label{eq:rg-formation-fast-environment-coupling}\\
 |h_t-h_t^*|
 &\le \mathrm C_L\sqrt{\mathrm T\mathrm J_1}
       \int_{t_{\mathrm{prev}}}^t e^{-\mathrm c_L\mathrm J_1(t-\sigma)}\mathrm e_\epsilon(y_\sigma) d\sigma,
 \label{eq:rg-formation-loop-environment-coupling}
\end{align}
with the adapted environment of the process. The row Lipschitz constant and
the direct parent row error \cref{eq:rg-formation-parent-row-environment}
convert these into proposal-row errors.

At random proposal times, condition on the whole independent future Poisson
grid, but not on its marks or on any
future continuous path, so that the Brownian law is unchanged. Evenness and
the uniform lower curvature on \(M_\ell\) give the first bound below on every
sheet and level; proposals preserve \(y\), projections contract its norm, and
the preparation has second moment at most \(\mathrm C_L(\mathrm A^2+\mathrm T)\),
which gives the second:
\begin{align}
 (\mathrm T\Delta-\nabla V_{\ell,j}\cdot\nabla)|y|^2
 &\le2\mathrm T\dim M_0-2\mathrm c_L|y|^2,
 \label{eq:rg-formation-grid-moment-generator}\\
 \sup_{t\le t_*}\E\bigl[|\mathbf X_{M_\ell,t}|^2\bigm|\text{Poisson grid}\bigr]
 &\le \mathrm C_L(\mathrm A^2+\mathrm T).
 \label{eq:rg-formation-grid-moment}
\end{align}
If the initial sample is fixed as well, the right side first contains its
squared norm, and integrating that sample returns the bound; the future grid
after a deterministic activation is independent of the past, so the argument
applies at each activation and at grid-selected times. Averaging the row
errors of
\cref{eq:rg-formation-fast-environment-coupling,eq:rg-formation-loop-environment-coupling}
with \cref{eq:rg-formation-grid-moment}, where the exponential convolution has
integral at most \(\mathrm C_L/\mathrm J_1\), over the distribution of every
past prefix, and dividing by the probability of no rare prefix, gives the
common expected per-proposal bound
\begin{equation}
 \eta_\epsilon
 =\mathrm C_L\left[\frac{\mathrm J_2}{\mathrm J_1}
       +\frac{\mathrm J_2}{\sqrt{\mathrm J_1}}
                       \left(1+\frac{\mathrm A}{\sqrt{\mathrm T}}\right)\right]
 =O_L(\epsilon^7),
 \label{eq:rg-formation-averaged-environment-error}
\end{equation}
since \(\mathrm Q_1=2\mathrm Q_2+20\) and \(\mathrm A/\sqrt{\mathrm T}=\epsilon^{-3}\).

Define one comparison source process on the same proposal grid, started at
the initial assignment and using the row \(\bar\beta_{j\cdot}(\vartheta_\circ(t))\),
\(\vartheta_\circ(t)=2-e^{-\epsilon t}\), a single comparison law common to both
processes while each process keeps its own law and coefficient. At a proposal
time \(t\), condition on the common starting sample, the full future grid, and
all previous marks, self marks included; the assignment and level are fixed
on \((t_{\mathrm{prev}},t)\), and no future continuous trajectory is conditioned
on. Starting the auxiliary diffusion from the fast point at
\(t_{\mathrm{prev}}\), the environmental coupling, then
\cref{eq:rg-formation-fast-mixing-bound}, then the Lipschitz comparison of
\(\vartheta(t)\) with \(\vartheta_\circ(t)\) bound the averaged conditional row
error, for gaps at least \(t_{\mathrm{gap}}\), by the second line below:
\begin{align}
 t_{\mathrm{gap}}&=\frac{k_L}{\mathrm J_1}\log(1/\epsilon),
 \label{eq:rg-formation-history-gap}\\
 &\mathrm C_L\left[
 \bigl(1+\mathrm A\sqrt{\mathrm J_1/\mathrm T}\bigr)e^{-\mathrm c_{\mathrm{fast},L}\mathrm J_1t_{\mathrm{gap}}}
       +\epsilon/\mathrm J_1+\eta_\epsilon+r_\epsilon\right];
 \label{eq:rg-formation-long-gap-row-error}
\end{align}
the moment of the point-start mixing factor, which includes
\(\sqrt{\mathrm J_1/\mathrm T}(\mathrm A+\sqrt{\mathrm T})\), is absorbed because
\(\mathrm A/\sqrt{\mathrm T}\ge1\), and averaging over the past marks returns
the grid-conditioned moments exactly. Short gaps use the trivial bound one.
For a rate-one Poisson grid the expected number of ordered pairs of proposal
times less than \(t_{\mathrm{gap}}\) apart is at most \(t_*t_{\mathrm{gap}}\),
by integrating the two-time intensity; proposals within \(t_{\mathrm{gap}}\)
after the initial time or one of the finitely many activations add at most
\((L+1)t_{\mathrm{gap}}\), and the expected total count is \(t_*\).

To pass from conditional row errors to the whole history, condition on a grid
with \(\mathrm n\) proposals, write the two joint mark laws as products of
their conditional kernels \(\mathsf K_m\) and \(\overline{\mathsf K}_m\), and
replace kernels in order, keeping the earlier prefix distributed according to
the process and integrating the comparison tail, whose mass is one. With
\(\mathsf j_m\in\mathcal J\) the \(m\)th proposal mark,
\begin{equation}
 \bigl\|\operatorname{Law}(\mathsf j_1,\ldots,\mathsf j_{\mathrm n})
          -\operatorname{Law}(\overline{\mathsf j}_1,\ldots,\overline{\mathsf j}_{\mathrm n})\bigr\|_{\mathrm{TV}}
 \le \sum_{m=1}^{\mathrm n}
       \E_{\text{prefix}}
       \|\mathsf K_m(\cdot\mid\mathsf j_{<m})
                 -\overline{\mathsf K}_m(\cdot\mid\mathsf j_{<m})\|_{\mathrm{TV}},
 \label{eq:rg-formation-history-telescoping}
\end{equation}
and integrating over the grid with the preceding bounds gives
\begin{align}
 \mathrm{Err}_{\mathrm{hist}}\le \mathrm C_L\big[
 &(t_*+L+1)t_{\mathrm{gap}}
 +t_*\bigl(1+\mathrm A\sqrt{\mathrm J_1/\mathrm T}\bigr)
                      e^{-\mathrm c_{\mathrm{fast},L}\mathrm J_1t_{\mathrm{gap}}}
 \notag\\
 &+t_*\epsilon/\mathrm J_1
       +t_*\eta_\epsilon+t_*r_\epsilon\big].
 \label{eq:rg-formation-full-history-error}
\end{align}
With \(\mathrm c_{\mathrm{fast},L}k_L>\mathrm Q_1/2+12\), the five contributions
are bounded by constants times \(\epsilon^{19}\varsigma_\epsilon\),
\(\epsilon^8\), \(\epsilon^{\mathrm Q_1}\), \(\epsilon^6\), and
\(\epsilon^{2\mathrm Q_1+19}+\epsilon^{3\mathrm Q_1+25}\varsigma_\epsilon^2\), so
\(\mathrm{Err}_{\mathrm{hist}}\le\mathrm C_L\epsilon^6\).

Fix the common initial sample and the grid, couple each mark sequence
maximally with the comparison sequence, and glue the two couplings along it;
the glued sequences disagree with probability at most
\(\mathrm C_L\epsilon^6\). The path spaces, continuous paths in a Euclidean
space and right-continuous paths in \(\mathcal J\), are standard Borel, so each
continuous trajectory can be drawn from its regular conditional law given the
initial sample, grid, and marks, which preserves both full-process marginals
and the parent projection of the initial sample. Once self-proposals are
discarded, each source path is a function of grid and marks, which proves the
first bound of \cref{eq:rg-formation-two-path-bounds}; finitely many
activations change only \(\mathrm C_L\). \hfill\(\square\)

\section{Coherence, source movement, and continuous paths}
\label{sec:rg-formation-dynamics-transfer}

\lemmaheading{Successive coherence with the original conditional marks}{lem:formation-coherence}
Under the preparation and regime of
\cref{eq:rg-formation-preparation-margin,eq:rg-formation-preparation-transport,eq:rg-formation-weakest-one-regime},
there are a finite \(k_{\mathrm{geom},L}\) and a tolerance bound
\(\eta_L^{\max}>0\), a fixed fraction of \(\mathrm m_L\), both depending only on
\(L\) and the fixed data, with the following property. For every
\(k_L>k_{\mathrm{geom},L}\) and every \(0<\eta_L\le\eta_L^{\max}\) there are
\(\epsilon_L,\mathrm c_L,\mathrm C_L>0\) such that, for \(0<\epsilon<\epsilon_L\),
conclusion (b) of \Cref{thm:finite-depth-adaptive-formation} holds for the fine
process and its successive parent descriptions on the windows
\(I_\ell=[t_\ell,2t_\ell]\), \(t_\ell=k_L\log(1/\epsilon)/\mathrm J_\ell\), with
the probability bound \cref{eq:rg-formation-coherence-probability}. On
\(I_\ell\) the discrepancies \(\mathfrak d_i\) and the residuals of level at
most \(\ell\) are below \(\eta_L\mathrm A\) and future bridges stay separated
by a fixed positive multiple of \(\mathrm A\), in the conditional
root-mean-square original residuals of
\cref{eq:rg-formation-inherited-edge-statistic} for eliminated coordinates,
separately for the physical posterior and the restricted optimizer.
\status{ESTABLISHED}

\noindent\emph{Proof.}
On an active level \(\ell'<\ell\), write the non-root Gaussian comparison of
\Cref{lem:formation-prepared-continuation} through a current normal component
\(z\in\mathcal U_{\ell'}\cap\mathcal U_\ell^\perp\), a future component
\(y\in\mathcal U_\ell\), and the separate loop component. The weighted flag
gives, for the current and future blocks of \(H_{\ell'}\),
\begin{equation}
 H_{zz}\succeq \mathrm c_L\mathrm J_\ell I,
 \qquad \|H_{zy}\|\le \mathrm C_L\mathrm J_{\ell+1},
 \qquad \|H_{yy}\|\le \mathrm C_L\mathrm J_{\ell+1},
 \qquad \mathrm J_{L+1}=0.
 \label{eq:rg-formation-current-future-blocks}
\end{equation}
Indeed rows of levels at most \(\ell\) annihilate \(\mathcal U_\ell\), so the
original cross and future blocks contain only strengths at most
\(\mathrm J_{\ell+1}\); in an earlier elimination through \(\ell''\), the
eliminated-current block is \(O_L(\mathrm J_{\ell''+1})\), the inverse hidden
block \(O_L(\mathrm J_{\ell''}^{-1})\), and the eliminated-future block
\(O_L(\mathrm J_{\ell+1})\), so the Schur correction is still
\(O_L(\mathrm J_{\ell+1})\); at level zero there is no correction, and
loop/normal cross blocks vanish in the Gaussian core.

The deterministic reference mean \(\mu^{\mathrm G}\) has norm at most
\(\mathrm C_L\mathrm A\) and zero loop component. Integrating
\(\dot\mu_z^{\mathrm G}=-\vartheta(t)(H_{zz}\mu_z^{\mathrm G}+H_{zy}\mu_y^{\mathrm G})\)
with \(1\le\vartheta\le21/10\) gives the first bound below; integrating the
full future row, of norm \(O_L(\mathrm J_{\ell+1})\), from the common center,
earlier projections preserving \(P_\ell\mu^{\mathrm G}\), gives the second for
a sequential process begun at \(t^{\mathrm{act}}_{\ell-1}\):
\begin{align}
 |\mu_z^{\mathrm G}(t)|
 &\le \mathrm C_L\mathrm A\left[e^{-\mathrm c_L\mathrm J_\ell(t-t_0)}+\frac{\mathrm J_{\ell+1}}{\mathrm J_\ell}\right],
 \label{eq:rg-formation-current-mean}\\
 |\mu^{\mathrm G}(t)-\mathrm AP_\ell\mathbf v_0|
 &\le \mathrm C_L\mathrm A\left[
 e^{-\mathrm c_L\mathrm J_\ell(t-t^{\mathrm{act}}_{\ell-1})}
 +\frac{\mathrm J_{\ell+1}}{\mathrm J_\ell}+\mathrm J_{\ell+1}t\right].
 \label{eq:rg-formation-window-mean}
\end{align}
The future terms are absent when \(\ell=L\), and a fine process uses \(t_0=0\).

For the fluctuations, on each active non-root space
\(\mathbf Z=\mathbf X^{\mathrm G}-\mu^{\mathrm G}\) obeys
\(d\mathbf Z=-\vartheta(t)H_{\ell'}\mathbf Zdt+\sqrt{2\mathrm T} dW\), with
\(\vartheta H_{\ell'}\succeq\mathrm h_LI\) and initial covariance at most
\(\mathrm C_L\mathrm TH_{\ell'}^{-1}\). Choose \(\mathrm a_*>0\) below
\(\mathrm h_L\) and small enough for the initial Gaussian exponential moment.
For \(\mathcal Y(\mathbf z)=e^{\upsilon}\) with
\(\upsilon=\mathrm a_*|\mathbf z|^2/(4\mathrm T)\), direct differentiation gives
the first line below; the initial expectation of \(\mathcal Y\) is bounded,
orthogonal retention decreases \(\mathcal Y\), and It\^o's formula up to the
first exit from a ball, localization, and the nonnegative stopped value of
\(\mathcal Y\) yield the second for every \(\varrho>0\), higher stiffnesses only
strengthening the negative term:
\begin{align}
 \mathscr L_t\mathcal Y
 &=\left[\frac{\mathrm a_*\mathrm d}{2}
 +\frac{\mathrm a_*^2|\mathbf z|^2}{4\mathrm T}
 -\frac{\mathrm a_*\vartheta\mathbf z^\top H_{\ell'}\mathbf z}{2\mathrm T}\right]\mathcal Y
 \le \mathrm a_*(\mathrm d/2-\upsilon)e^\upsilon\le \mathrm C_L,
 \label{eq:rg-formation-gaussian-exponential-generator}\\
 \Pr\Bigl\{\sup_{t\le t_*}|\mathbf X^{\mathrm G}_t-\mu^{\mathrm G}(t)|&>\varrho\mathrm A\Bigr\}
 \le \mathrm C_L(1+t_*)
       \exp\left[-\mathrm c_L\varrho^2\mathrm A^2/\mathrm T\right].
 \label{eq:rg-formation-gaussian-maximum}
\end{align}

Let \(\bar r_L\) bound the norms of the finitely many residual and discrepancy
covectors, let \(\eta_L^{\max}\) be a fixed fraction of \(\mathrm m_L\), fix
\(\eta_L\in(0,\eta_L^{\max}]\), and choose a fluctuation radius
\(\varrho\mathrm A\) with \(\bar r_L\varrho\mathrm A\) smaller than the reserved
tolerance margin, \(\mathrm c_L\) in the probability depending on \(\varrho\).
Take for example \(k_{\mathrm{geom},L}=2/\mathrm c_L\) for the restoring constant
of \cref{eq:rg-formation-current-mean}; increasing this finite constant is
allowed. For small \(\epsilon\), \(t^{\mathrm{act}}_{\ell-1}<t_\ell/2\), so the
exponential term of \cref{eq:rg-formation-window-mean} tends to zero uniformly
on \(I_\ell\), and
\begin{equation}
 \frac{\mathrm J_{\ell+1}}{\mathrm J_\ell}\longrightarrow0,
 \qquad \mathrm J_{\ell+1}t_\ell\longrightarrow0\quad(\ell<L),
 \qquad \frac{t^{\mathrm{act}}_L}{t_*}\longrightarrow0.
 \label{eq:rg-formation-geometric-entrance}
\end{equation}
The nonlinear correction of \cref{eq:rg-formation-path-force-bound} is at most
\(\mathrm C_L\mathrm A\epsilon^{5/2}\), including the admitted initial transport
and finitely many inherited errors; it enters through the source force only via
the uniform bound \(B_*\sqrt{\mathrm T}\), so it holds for every adapted
assignment path. Since \(\mathrm A^2/\mathrm T=\epsilon^{-6}\), the strict
prepared residual margins and \cref{eq:rg-formation-gaussian-maximum} give the
asserted probability and separation on every window. The stable current-block
equation keeps bounding formed residuals as their future forcing decreases,
and after the last normal block there is no future forcing, so the final
regime persists through \(t_*\) on the same event, with every loop retained.

For an eliminated edge and a cumulative elimination through \(\ell'\ge1\),
write a fine non-root state as \(\mathbf z+y\) with
\(\mathbf z\in W_{\ell'}^\perp\) and \(y\in M_{\ell'}\). The optimizing and
physical Gaussian precisions are \(\vartheta H\) and
\(\vartheta H-\mathrm TR_{\mathrm{ev}}\). Every row of \(\mathfrak D\) lies in the
first normal layer, so \(R_{\mathrm{ev}}\) annihilates \(\mathcal U_1\),
\(\mathcal H\), and \(\mathcal K\), and with \(\mathrm TR_{\mathrm{ev}}\preceq H/2\)
both laws have hidden precision at least \(\mathrm c_L\mathrm J_{\ell'}I\),
hidden/future block \(O_L(\mathrm J_{\ell'+1})\), and zero hidden/loop block.
The Gaussian conditional mean therefore has norm at most
\(\mathrm C_L(\mathrm J_{\ell'+1}/\mathrm J_{\ell'})|y_{\mathcal U_{\ell'}}|\) and
covariance trace at most \(\mathrm C_L\mathrm T/\mathrm J_{\ell'}\), and, given
the full assignment, each exact hidden density is a normalized tilt of its
Gaussian conditional with log oscillation at most \(B_0\), hence density ratio
at most \(e^{B_0}\). For
\(\mu_\vartheta\in\{\boldsymbol\Pi_{\vartheta,\mathsf o},q^\star_\vartheta\}\) and
every retained state with \(|y_{\mathcal U_{\ell'}}|\le\mathrm C_L\mathrm A\), this
gives the first bound below, and for every original residual covector the
triangle inequality and its reverse in conditional \(L^2\) give the second:
\begin{align}
 \E_{\mu_\vartheta}
   [|\mathbf z|^2\mid P_{\ell'}\mathbf X=y,S=j]
 &\le \mathrm C_Le^{B_0}\left[
       \mathrm A^2\left(\frac{\mathrm J_{\ell'+1}}{\mathrm J_{\ell'}}\right)^2
       +\frac{\mathrm T}{\mathrm J_{\ell'}}\right],
 \label{eq:rg-formation-hidden-original-mark-moment}\\
 \left|
 \sqrt{\E_{\mu_\vartheta}[r_e(\mathbf X)^2\mid P_{\ell'}\mathbf X=y,S=j]}
 -|r_e(y)|\right|
 &\le\|\mathbf r_e\|
      \sqrt{\E_{\mu_\vartheta}[|\mathbf z|^2\mid P_{\ell'}\mathbf X=y,S=j]}.
 \label{eq:rg-formation-original-mark-rms}
\end{align}
On the event already constructed the right side of
\cref{eq:rg-formation-original-mark-rms} is \(o(\mathrm A)\), uniformly in the
assignment and separately for the two laws, and the same holds for
\(\mathfrak d_i\); this transfers the small formed and the separated future
residuals to \cref{eq:rg-formation-inherited-edge-statistic}, with no erased
mark replaced by its value at \(P_{\ell'}\mathbf x\). The bridge levels join
successive pairs of already formed components, giving the stated counts.
\hfill\(\square\)

\lemmaheading{Source envelope and changing hidden kernels}{lem:formation-source-envelope}
Every process or finite replacement lineage of
\Cref{lem:formation-prepared-continuation} satisfies
\cref{eq:rg-formation-source-movement} for the distinguished binary source,
from any of its activation times. Between \(\vartheta=1\) and its final
parameter, the first hidden physical conditional and the first hidden
optimizing conditional each change by relative entropy greater than
\(1/200\), uniformly at the same retained input and assignment.
\status{ESTABLISHED}

\noindent\emph{Proof.}
The distinguished receiver's posterior row \(\beta_{i_0,\pm}\), a factor of
the product row \cref{eq:rg-formation-full-source-row}, has the log odds of the
first line below. Since \(0\le\chi\le1\), it lies between
\(\operatorname{sig}(4(\vartheta-3/2)\mp\delta)\), and summing the inherited
row over the assignments with distinguished label \(+\) and taking the
expectation preserves these bounds, so the generator applied to the indicator
of that label gives the second line, with \(\lambda(t)=\Pr\{S_{i_0}(t)=+\}\):
\begin{align}
 \log\frac{\beta_{i_0,+}(\mathbf x,\vartheta)}
               {\beta_{i_0,-}(\mathbf x,\vartheta)}
 &=4(\vartheta-3/2)
  -\delta\left[
 \chi\bigl(\sqrt{\mathrm J_1/\mathrm T} r_{i_0,+}(\mathbf x)\bigr)
 -\chi\bigl(\sqrt{\mathrm J_1/\mathrm T} r_{i_0,-}(\mathbf x)\bigr)\right],
 \label{eq:rg-formation-source-log-odds}\\
 \lambda'&=\bar\beta_\lambda-\lambda,
 \qquad
 \operatorname{sig}(4(\vartheta-3/2)-\delta)
 \le\bar\beta_\lambda(t)\le
 \operatorname{sig}(4(\vartheta-3/2)+\delta).
 \label{eq:rg-formation-source-forcing-envelope}
\end{align}
Neither diffusion nor a marked projection changes this indicator, so variation
of constants remains valid across every activation, even when the inherited row
depends on the current assignment and continuous state.

By \Cref{lem:formation-parameter-comparison}, uniformly before
\(t^{\mathrm{act}}_L=2k_L\log(1/\epsilon)\),
\(\vartheta\le1+\mathrm C_L\epsilon\log(1/\epsilon)+\mathrm C_Lr_\epsilon\). A
solution of \cref{eq:rg-formation-source-forcing-envelope} cannot exceed the
larger of its starting value and the supremum forcing, so the preparation
tolerance and \(\|\operatorname{sig}'\|_\infty\le1/4\) give the first bound
below. For small \(\epsilon\),
every activation precedes \(t_*-10\) and \(\vartheta\ge7/4\) throughout
\([t_*-10,t_*]\), so nonnegativity of the earlier contribution in variation of
constants yields the second; for \(\delta<10^{-4}\), the elementary bounds
\(e^{1-\delta}>249/100\), \(e^2>7\), and \(e^{10}>2^{10}\) give the third:
\begin{align}
 &\lambda(t^{\mathrm{act}}_\ell)
 \le\operatorname{sig}(-2)+10^{-2}+\delta/4
       +\mathrm C_L[\epsilon\log(1/\epsilon)+r_\epsilon],
 \qquad 0\le \ell\le L,
 \label{eq:rg-formation-source-at-activation}\\
 &\lambda(t_*)\ge(1-e^{-10})\operatorname{sig}(1-\delta),
 \label{eq:rg-formation-source-final-mass}\\
 &(1-e^{-10})\operatorname{sig}(1-\delta)-\operatorname{sig}(-2)
       -10^{-2}-\delta/4
 >\frac{1023}{1024} \frac{249}{349}
       -\frac18-\frac1{100}-\frac1{40000}
 >\frac{57}{100}.
 \label{eq:rg-formation-source-strict-margin}
\end{align}
The remaining error in \cref{eq:rg-formation-source-at-activation} tends to
zero, hence \(\lambda(t_*)-\lambda(t^{\mathrm{act}}_\ell)>1/2\) for every
activation, which for a binary row is its total-variation change.

For the hidden kernels, fix the retained input and the assignment, and first
omit the bounded source factor. On the first hidden normal space let
\(\Omega_1,\Omega_{\mathrm f}\) be the initial and final Gaussian precisions, with
covariances \(\mathrm T\Omega_1^{-1},\mathrm T\Omega_{\mathrm f}^{-1}\). For the
optimizer the precision scales with \(\vartheta\); for the physical law, on the
non-root space,
\begin{equation}
 \Omega_1=H-\mathrm TR_{\mathrm{ev}}\succ0,
 \qquad
 \Omega_{\mathrm f}=\Omega_1+(\vartheta_{\mathrm f}-1)H
       \succeq\vartheta_{\mathrm f}\Omega_1,
 \qquad \vartheta_{\mathrm f}\ge7/4.
 \label{eq:rg-formation-physical-precision-growth}
\end{equation}
Restriction to a conditional block preserves the order, and so does an earlier
Gaussian marginal, whose Schur quadratic is the minimum over the discarded
coordinates, so in either case \(\Omega_{\mathrm f}\succeq(5/4)\Omega_1\). For
the Gaussian laws
\(\gamma_1,\gamma_{\mathrm f}\) on a hidden space of dimension
\(\mathrm d_{\mathrm h}\ge1\), whitening \(\Omega_1\) and diagonalizing the
conjugated final precision, with eigenvalues \(\varrho_i\ge5/4\), and completing
the square in the Hellinger affinity give the first line below; each factor
decreases for \(\varrho_i\ge1\), so the affinity is at most
\([2\sqrt{5/4}/(1+5/4)]^{\mathrm d_{\mathrm h}/2}\). Each exact conditional
\(\nu_1,\nu_{\mathrm f}\) is a normalized source tilt of its Gaussian
conditional bounded above by \(e^{B_0}\), also after exact hidden integration,
so \(\mathcal A(\nu_1,\nu_{\mathrm f})\le e^{B_0}\mathcal A(\gamma_1,\gamma_{\mathrm f})\),
and Jensen's inequality under \(\nu_1\) gives
\(-\KL(\nu_1\Vert\nu_{\mathrm f})/2\le\log\mathcal A(\nu_1,\nu_{\mathrm f})\).
With \(\log(1+1/80)\ge1/81\) this gives the second line, separately for the
physical and optimizing conditionals and uniformly in the retained input and
assignment:
\begin{align}
 \mathcal A(\gamma_1,\gamma_{\mathrm f})
 &=\prod_{i=1}^{\mathrm d_{\mathrm h}}
       \left(\frac{2\sqrt{\varrho_i}}{1+\varrho_i}\right)^{1/2}
       e^{-\mathrm M},
 \qquad \mathrm M\ge0,
 \label{eq:rg-formation-gaussian-affinity}\\
 \KL(\nu_1\Vert\nu_{\mathrm f})
 &\ge\frac{\mathrm d_{\mathrm h}}{2}\log\frac{81}{80}-2B_0
 \ge\frac1{162}-\frac1{5000}>\frac1{200}.
 \label{eq:rg-formation-hidden-kernel-change}
\end{align}
\hfill\(\square\)

\lemmaheading{Concentration of continuous paths on a common deterministic path}{lem:formation-continuous-path}
Let a source process and its exact parent start from a common retained sample
and parameter at \(t^{\mathrm{act}}_\ell\). There is a deterministic path
\(\bar{\mathbf x}_\ell:[t^{\mathrm{act}}_\ell,t_*]\to W_\ell\), namely the
parent's Gaussian reference mean on the non-root coordinates and the constant
equilibrium root mean, such that
\begin{equation}
 \Pr\Bigl\{\sup_{t^{\mathrm{act}}_\ell\le t\le t_*}|P_\ell\mathbf X_t-\bar{\mathbf x}_\ell(t)|
   >\mathrm C_L\mathrm A\epsilon^2\Bigr\}
 +\Pr\Bigl\{\sup_{t^{\mathrm{act}}_\ell\le t\le t_*}|\overline{\mathbf X}_t-\bar{\mathbf x}_\ell(t)|
   >\mathrm C_L\mathrm A\epsilon^2\Bigr\}
 \le \mathrm C_L(1+t_*)e^{-\mathrm c_L/\epsilon^2}.
 \label{eq:rg-formation-common-path-proof}
\end{equation}
Consequently the displacement bound of \cref{eq:rg-formation-two-path-bounds}
holds under every coupling of the two processes, in particular under the
coupling of \Cref{lem:formation-history-coupling}. A finite number of
successive replacements changes only the constants. \status{ESTABLISHED}

\noindent\emph{Proof.}
First compare Gaussian means at one deterministic parameter path. In the split
of \cref{eq:rg-formation-current-future-blocks} the inherited normal precision
is \(H_\ell=H_{yy}-H_{zy}^\top H_{zz}^{-1}H_{zy}\). With
\(\mathfrak w=\mu_z^{\mathrm G}+H_{zz}^{-1}H_{zy}\mu_y^{\mathrm G}\), the future
mean of the source process satisfies the first line below, the bound following
from \cref{eq:rg-formation-current-mean} and
\(\|H_{zz}^{-1}H_{zy}\|\le\mathrm C_L\mathrm J_{\ell+1}/\mathrm J_\ell\), and
integrating the contractive future propagator from a matching retained mean
gives the second:
\begin{align}
 \dot\mu_y^{\mathrm G}&=-\vartheta H_\ell\mu_y^{\mathrm G}-\vartheta H_{zy}^\top\mathfrak w,
 \qquad
 |\mathfrak w(t)|\le \mathrm C_L\mathrm A\left[
 e^{-\mathrm c_L\mathrm J_\ell(t-t^{\mathrm{act}}_\ell)}+\frac{\mathrm J_{\ell+1}}{\mathrm J_\ell}\right],
 \label{eq:rg-formation-gaussian-mean-defect}\\
 \sup_{t^{\mathrm{act}}_\ell\le t\le t_*}
 |\mu_y^{\mathrm G}(t)-\bar\mu^{\mathrm G}(t)|
 &\le \mathrm C_L\mathrm A\left[
 \frac{\mathrm J_{\ell+1}}{\mathrm J_\ell}
 +t_*\frac{\mathrm J_{\ell+1}^2}{\mathrm J_\ell}\right].
 \label{eq:rg-formation-gaussian-mean-reduction}
\end{align}
There is no future normal contribution when \(\ell=L\), and the loop Gaussian
core is unchanged by the reduction and has zero mean. If the independently
learned parameters differ, with
\(\mathrm e_\vartheta=\sup|\vartheta-\bar\vartheta|\le\mathrm C_Lr_\epsilon\) by
\Cref{lem:formation-parameter-comparison}, Young's inequality in the \(H_\ell\)
quadratic form gives, for homogeneous parent means \(\mu,\bar\mu\) with
difference \(\mathfrak e\),
\begin{equation}
 \frac d{dt}|\mathfrak e|^2
 \le-\vartheta \mathfrak e^\top H_\ell\mathfrak e
       +\mathrm C_L\mathrm e_\vartheta^2 \bar\vartheta
                         \bar\mu^\top H_\ell\bar\mu,
 \label{eq:rg-formation-mean-parameter-energy}
\end{equation}
and since \(\frac d{dt}|\bar\mu|^2=-2\bar\vartheta\bar\mu^\top H_\ell\bar\mu\),
the time integral of the final energy is at most \(\mathrm C_L\mathrm A^2\), so
the extra displacement is at most \(\mathrm C_L\mathrm A \mathrm e_\vartheta\).

Attach to each process its own Gaussian comparison from
\Cref{lem:formation-prepared-continuation}; the discrepancy between the
parent's projected starting reference and its full retained starting sample is
included in \cref{eq:rg-formation-path-force-bound}. On the non-root
coordinates let \(\bar{\mathbf x}_\ell=\bar\mu^{\mathrm G}\) be the parent's
reference mean. For the parent,
\cref{eq:rg-formation-path-force-bound,eq:rg-formation-gaussian-maximum} with
\(\varrho=\epsilon^2\) give
\(\sup_t|\overline{\mathbf X}_t-\bar{\mathbf x}_\ell(t)|\le\mathrm C_L\mathrm A(\epsilon^2+\epsilon^{5/2})\)
outside a set of probability \(\mathrm C_L(1+t_*)e^{-\mathrm c_L/\epsilon^2}\).
For the source process the same two estimates for its own reference, then
\cref{eq:rg-formation-gaussian-mean-reduction} and the parameter comparison,
give, outside a set of the same form,
\begin{equation}
 \sup_{t^{\mathrm{act}}_\ell\le t\le t_*}|P_\ell\mathbf X_t-\bar{\mathbf x}_\ell(t)|
 \le \mathrm C_L\mathrm A\left[
 \epsilon^2+\epsilon^{5/2}+r_\epsilon
 +\frac{\mathrm J_{\ell+1}}{\mathrm J_\ell}
 +t_*\frac{\mathrm J_{\ell+1}^2}{\mathrm J_\ell}\right].
 \label{eq:rg-formation-continuous-error-terms}
\end{equation}
For \(\ell<L\) the last term equals \(2\epsilon^{19}\),
\(\mathrm J_{\ell+1}/\mathrm J_\ell=\epsilon^{\mathrm Q_{\ell+1}+20}\), and
\(r_\epsilon=o(\epsilon^2)\); both future terms vanish at \(\ell=L\). The root
coordinates of each process form the same autonomous Ornstein--Uhlenbeck
process in law, with stiffness \(1+\kappa\ge1\), started at its equilibrium
and independent of the source history, and the argument of
\cref{eq:rg-formation-gaussian-maximum} keeps each within \(\epsilon^2\mathrm A\)
of the constant mean \(\kappa z_\star/(1+\kappa)\) outside a set of the same
form. Taking that constant as the root component of \(\bar{\mathbf x}_\ell\)
proves \cref{eq:rg-formation-common-path-proof}. Each of its events concerns
the law of one process alone, so for any coupling the displacement event of
\cref{eq:rg-formation-two-path-bounds}, after doubling \(\mathrm C_L\), lies in
their union; a finite sequence of replacements uses a finite triangle
inequality and union bound. \hfill\(\square\)

\section{The Gaussian core across an early interface}
\label{sec:rg-formation-core-proof}

\noindent\emph{Proof of parts (b) and (c) of \Cref{prop:rg-formation-gaussian-core}.}
Fix \(1\le\ell\le L-1\). In this section \(\mathrm C\) and \(\mathrm c\)
denote positive constants depending only on \(L\) and the carrier, never on
the strengths, \(\mathrm A\), \(\mathrm T\), \(t\), or \(\epsilon\); put
\(\mathrm J_{L+1}=0\) and
\(\bar\varrho_\ell=\varrho_{\ell-1}+\varrho_\ell+\varrho_{\ell+1}\).

\emph{Flag layers and block scales.} In the flag coordinates
\(w=(w_1,\ldots,w_L)\) of \cref{eq:rg-formation-weighted-flag},
\(\mathbf x_{\mathcal U}=\sum_s\mathsf Q_sw_s\), the quadratic form
\(\mathbf x_{\mathcal U}^\top H\mathbf x_{\mathcal U}=\sum_s\mathrm J_s|\sum_{r\le s}\mathsf B_{sr}w_r|^2\)
has the blocks
\((H_{\mathcal U})_{rr'}=\sum_{s\ge\max(r,r')}\mathrm J_s\mathsf B_{sr}^\top\mathsf B_{sr'}\),
each of norm at most \(\mathrm C\mathrm J_{\max(r,r')}\); the loop block
\(\mathrm J_1I_{\mathcal H}\) is orthogonal to every layer and plays no part.
On the active level \(\ell-1\) the normal part of \(H_{\ell-1}\) is the Schur
complement of \(H_{\mathcal U}\) onto the layers \(r\ge\ell\), obtained by
eliminating the layers \(\mathcal E=\{1,\ldots,\ell-1\}\), with the discarded
block \(z=w_\ell\) and the retained block \(y=(w_{\ell+1},\ldots,w_L)\). The
lower bound of \cref{eq:rg-formation-weighted-flag} on \(\mathcal E\) gives
\(H_{\mathcal E\mathcal E}\succeq\mathrm c\mathrm J_{\ell-1}I\), so the Schur
correction to a block \((r,r')\) with \(r,r'\ge\ell\) has norm at most
\(\mathrm C\mathrm J_r\mathrm J_{r'}/\mathrm J_{\ell-1}\le\mathrm C\mathrm J_{\max(r,r')}\varrho_{\ell-1}\).
Since in each block the summand with the largest admissible strength leads
and the remainder carries the next strength,
\begin{equation}
\begin{aligned}
H_{zz}&=\mathrm J_\ell\bigl[\mathsf B_{\ell\ell}^\top\mathsf B_{\ell\ell}+O(\varrho_\ell+\varrho_{\ell-1})\bigr],\\
H_{z,\ell+1}&=\mathrm J_{\ell+1}\bigl[\mathsf B_{\ell+1,\ell}^\top\mathsf B_{\ell+1,\ell+1}+O(\varrho_{\ell+1}+\varrho_{\ell-1})\bigr],\\
H_{\ell+1,\ell+1}&=\mathrm J_{\ell+1}\bigl[\mathsf B_{\ell+1,\ell+1}^\top\mathsf B_{\ell+1,\ell+1}+O(\varrho_{\ell+1}+\varrho_{\ell-1})\bigr],
\end{aligned}
\label{eq:rg-formation-core-block-scales}
\end{equation}
with the constants of the remainders depending only on \(L\) and the carrier,
together with \(\|H_{z,r}\|\le\mathrm C\mathrm J_r\) and
\(\|H_{\ell+1,r}\|\le\mathrm C\mathrm J_{\ell+1}(\varrho_{\ell+1}+\varrho_{\ell-1})\)
for \(r\ge\ell+2\), and \(\|H_{r,\cdot}\|\le\mathrm C\mathrm J_r\) for every
row block \(r\ge\ell+1\), in particular
\(\|H_{zy}\|\le\mathrm C\mathrm J_{\ell+1}\). Hence, for
\(\bar\varrho_\ell\le\mathrm c_0\) with \(\mathrm c_0\) small,
\(\lambda_{\min}(H_{zz})\ge\tfrac34\mathfrak j_\ell\mathrm J_\ell\),
\(\|H_{zz}^{-1}\|\le\mathrm C/\mathrm J_\ell\), and
\(\|H_{zz}^{-1}H_{zy}\|\le\mathrm C\varrho_\ell\). The same computation at
every earlier active level \(p\le\ell-1\) bounds the rows of layer
\(\ell+1\) of \(H_p\) by \(\mathrm C\mathrm J_{\ell+1}\).

\emph{Full row rank.} Write \(\mathbf r_e\) for the covector of \(r_e\).
Since \(\mathcal U_\ell=\mathcal U_{\ell-1}\cap\bigcap_e\ker\mathbf r_e\) over
the rows \(e\) of \(B_\ell\), the layer spanned by \(\mathsf Q_\ell\) is the
span of the projections \(P_{\mathcal U_{\ell-1}}\mathbf r_e\), and
\(\mathsf B_{\ell+1,\ell}\) has full row rank exactly when the matrix
\(\mathrm M_{fe}=\mathbf r_f\cdot P_{\mathcal U_{\ell-1}}\mathbf r_e\), with
\(f\) over the rows of \(B_{\ell+1}\) and \(e\) over those of \(B_\ell\), has
full row rank, because
\(\mathrm M_{fe}=(P_{\mathcal U_{\ell-1}}\mathbf r_f)\cdot(P_{\mathcal U_{\ell-1}}\mathbf r_e)\),
the component of \(P_{\mathcal U_{\ell-1}}\mathbf r_f\) in \(\mathcal U_\ell\)
pairs to zero with the layer, and the second factors span the layer. For
\(\ell=1\), \(\mathcal U_0=\mathcal U\) contains every row of \(B\), so
\(\mathrm M_{fe}=\mathbf r_f\cdot\mathbf r_e\); in the column of the
discrepancy row of agent \(i\) the entry of a bridge \(f\) is
\([i=\operatorname{head}(f)]-[i=\operatorname{tail}(f)]\), and the level-two
bridges, those of odd index, join disjoint pairs of triangles and have
pairwise disjoint endpoint sets, so the columns of the discrepancy rows of
their heads form an identity submatrix. For \(\ell\ge2\) let \(S\) be the
set of bridges of level at least \(\ell\), let \(\boldsymbol\partial_f\in\R^V\)
be the vector with entry \(1\) at the head of the bridge \(f\), \(-1\) at its
tail, and \(0\) elsewhere, let \(\partial_S\) have the rows
\(\boldsymbol\partial_f^\top\), \(f\in S\), and let \(\partial\) be the
oriented incidence operator of the whole graph, so that
\(\partial^\top\partial\) is its Laplacian. For a vector \(\mathsf b\) on the
bridges put \(\mathbf x_{\mathrm{br}}(\mathsf b)=(0,0,-\mathsf b)\), the
configuration with only the bridge links, and
\(\mathbf x(\mathsf b)=P_{\mathcal K^\perp}\mathbf x_{\mathrm{br}}(\mathsf b)\).
A configuration of \(\mathcal U_1\) has \(\bar k=\bar m=:t\) and, in each
triangle, three equal residuals which sum to \(-\mathfrak h_q=0\) by
orthogonality to \(\mathcal H\), so its internal links are \(\partial t\) and
its bridge links are \((\partial t)_f-\mathsf b_f\) with \(\mathsf b_f=r_f\);
it is therefore \((t,t,\partial t)+\mathbf x_{\mathrm{br}}(\mathsf b)\) and,
being orthogonal to \(\mathcal K\), equals \(\mathbf x(\mathsf b)\).
Conversely \(\mathbf x_{\mathrm{br}}(\mathsf b)\in W_1\cap\mathcal H^\perp\)
and \(\mathcal K\subseteq W_1\cap\mathcal H^\perp\), so
\(\mathbf x(\mathsf b)\in\mathcal U_1\), and
\(r_f(\mathbf x(\mathsf b))=r_f(\mathbf x_{\mathrm{br}}(\mathsf b))=\mathsf b_f\)
because residuals annihilate \(\mathcal K\). Hence
\(\mathcal U_{\ell-1}=\{\mathbf x(\mathsf b):\operatorname{supp}\mathsf b\subseteq S\}\)
with \(\mathsf b\mapsto\mathbf x(\mathsf b)\) injective, and from
\(\langle(t,t,\partial t),(t',t',\partial t')\rangle=t^\top(2I+\partial^\top\partial)t'\)
and
\(\langle(t,t,\partial t),\mathbf x_{\mathrm{br}}(\mathsf b)\rangle=-t^\top\partial_S^\top\mathsf b\)
the Gram matrix is
\(\langle\mathbf x(\mathsf b),\mathbf x(\mathsf b')\rangle=\mathsf b^\top\mathrm G\mathsf b'\)
with \(\mathrm G=I-\partial_S(2I+\partial^\top\partial)^{-1}\partial_S^\top\)
on \(\R^S\). The functional \(\mathsf b\mapsto\mathsf b_f\), \(f\in S\), has
the Riesz representative \(\mathbf x(\mathsf b^{(f)})\) with
\(\mathsf b^{(f)}\) the \(f\)th column of \(\mathrm G^{-1}\), so
\(P_{\mathcal U_{\ell-1}}\mathbf r_f=\mathbf x(\mathsf b^{(f)})\) and
\(\mathrm M_{fe}=(\mathrm G^{-1})_{fe}\). Since
\(2I+\partial^\top\partial-\partial_S^\top\partial_S=2I+\partial_{\setminus S}^\top\partial_{\setminus S}\),
with \(\partial_{\setminus S}\) the incidence operator of the graph with the
bridges of \(S\) removed, is positive definite, the Woodbury identity gives
\(\mathrm M_{fe}=\delta_{fe}+\boldsymbol\partial_f^\top(2I+\partial_{\setminus S}^\top\partial_{\setminus S})^{-1}\boldsymbol\partial_e\).
The graph with the bridges of \(S\) removed is the disjoint union of \(2^{L-\ell+1}\)
connected blocks of \(2^{\ell-2}\) consecutive triangles, and the inverse is
block diagonal with entrywise strictly positive blocks: on a block \(B\) with
Laplacian \(\partial_B^\top\partial_B\) and maximal degree \(\mathrm D_B\),
the matrix \((\mathrm D_B+2)I-(2I+\partial_B^\top\partial_B)\) is entrywise
nonnegative, irreducible, and of spectral radius at most
\(\mathrm D_B<\mathrm D_B+2\), since the Laplacian's eigenvalues lie in
\([0,2\mathrm D_B]\), so the Neumann series of
\((2I+\partial_B^\top\partial_B)^{-1}\) in it converges and is entrywise
positive. The bridges of \(S\) are those of index \(j=\mathrm w2^{\ell-2}\),
\(1\le\mathrm w\le2^{L-\ell+1}-1\), of level \(\ell+v_2(\mathrm w)\), the
bridge at position \(\mathrm w\) joining block \(\mathrm w\) to block
\(\mathrm w+1\); level \(\ell\) means \(\mathrm w\) odd and level
\(\ell+1\) means \(\mathrm w\equiv2\pmod4\). For \(f\ne e\) in \(S\), with
positions \(\mathrm w_f,\mathrm w_e\), the entry
\(\boldsymbol\partial_f^\top(2I+\partial_{\setminus S}^\top\partial_{\setminus S})^{-1}\boldsymbol\partial_e\)
is a sum over the blocks both bridges touch, so it vanishes unless
\(|\mathrm w_f-\mathrm w_e|=1\), and then it is plus or minus one entry of
the positive inverse on the shared block, hence nonzero. Pair each
level-\((\ell+1)\) bridge \(f\) with the level-\(\ell\) bridge \(e(f)\) at
position \(\mathrm w_f+1\), which exists because \(4\) divides
\(2^{L-\ell+1}\); then \(\mathrm M_{f,e(f')}\ne0\) exactly when
\(|\mathrm w_f-\mathrm w_{f'}-1|=1\), that is
\(\mathrm w_f-\mathrm w_{f'}\in\{0,2\}\), and as
\(\mathrm w_f-\mathrm w_{f'}\in4\mathbb Z\) exactly when \(f=f'\). The
square submatrix \((\mathrm M_{f,e(f')})_{f,f'}\) is diagonal with nonzero
entries, so the rows of \(\mathrm M\) of level \(\ell+1\) are independent and
\(\mathsf B_{\ell+1,\ell}\) has full row rank; as
\(\mathsf B_{\ell+1,\ell+1}\mathsf B_{\ell+1,\ell+1}^\top\) is invertible and
\(\mathsf B_{\ell+1,\ell}^\top\) is injective, \(\mathfrak Y_\ell\) is
injective. Nothing used the choice of bridge endpoints.

\emph{Slaving on the active interval.} Let level \(\ell-1\) be active on
\([t_0,t^{\mathrm{act}}_\ell]\) with \(t_0=t^{\mathrm{act}}_{\ell-1}\), with
mean \(\bar{\mathbf x}(t)\) in the layers \(r\ge\ell\),
\(\dot{\bar{\mathbf x}}=-H_{\ell-1}\bar{\mathbf x}\), and
\(|\bar{\mathbf x}(t_0)|\le\mathrm C\mathrm A\), which holds because
\(|\bar{\mathbf x}(0)|=\mathrm A|\mathbf v_0|\), the flow contracts
\(|\bar{\mathbf x}|\), and the projections do not increase it. Write
\(\mathbf g=H_{\ell-1}\bar{\mathbf x}\), so that \(\mathfrak v=\mathbf g_z\)
and \(\mathrm T\Delta_{\ell-1:\ell}=\tfrac12\mathbf g_z^\top H_{zz}^{-1}\mathbf g_z\)
by \cref{eq:rg-formation-gaussian-core-loss}, and
\(\mathsf b_{\ell+1}(t)=B_{\ell+1}P_\ell\bar{\mathbf x}(t)=\mathsf B_{\ell+1,\ell+1}\bar{\mathbf x}_{\ell+1}(t)\),
since \(P_\ell\) removes the layer \(\ell\) and \(\mathsf B_{\ell+1,r}=0\) for
\(r>\ell+1\). The row bound gives
\(|\dot{\bar{\mathbf x}}_{\ell+1}|=|\mathbf g_{\ell+1}|\le\mathrm C\mathrm J_{\ell+1}\mathrm A\),
and the same bound at the earlier levels, whose projections \(P_p\),
\(p\le\ell\), leave the layer \(\ell+1\) unchanged, gives the second line of
\cref{eq:rg-formation-core-two-sided}. Since
\(\dot{\mathbf g}=-H_{\ell-1}\mathbf g\), the variable
\(\mathbf u=\mathbf g_z+H_{zz}^{-1}H_{zy}\mathbf g_y\) satisfies
\(\dot{\mathbf u}=-H_{zz}\mathbf u-H_{zz}^{-1}H_{zy}(H_{\ell-1}\mathbf g)_y\),
where \(|(H_{\ell-1}\mathbf g)_y|\le\mathrm C\mathrm J_{\ell+1}(|\mathbf g_z|+|\mathbf g_y|)\),
\(|\mathbf g_y|\le\mathrm C\mathrm J_{\ell+1}\mathrm A\), and
\(|\mathbf g_z|\le|\mathbf u|+\mathrm C\varrho_\ell\mathrm J_{\ell+1}\mathrm A\).
Variation of constants with \(\|e^{-H_{zz}s}\|\le e^{-\lambda s}\),
\(\lambda=\lambda_{\min}(H_{zz})\), gives
\begin{equation}
|\mathbf u(t)|\le e^{-\lambda(t-t_0)}|\mathbf u(t_0)|
+\int_{t_0}^te^{-\lambda(t-s)}\bigl(\mathrm C\varrho_\ell^2\mathrm J_\ell|\mathbf u(s)|
+\mathrm C\varrho_\ell\mathrm J_{\ell+1}^2\mathrm A\bigr)ds,
\label{eq:rg-formation-core-volterra}
\end{equation}
and iterating the positive Volterra operator of the right side compares
\(|\mathbf u|\) with the solution of the scalar equation with decay rate
\(\lambda-\mathrm C\varrho_\ell^2\mathrm J_\ell\), constant forcing
\(\mathrm C\varrho_\ell\mathrm J_{\ell+1}^2\mathrm A\), and initial value
\(|\mathbf u(t_0)|\), so that
\(|\mathbf u(t)|\le e^{-(\lambda-\mathrm C\varrho_\ell^2\mathrm J_\ell)(t-t_0)}|\mathbf u(t_0)|
+\mathrm C\varrho_\ell\mathrm J_{\ell+1}^2\mathrm A/(\lambda-\mathrm C\varrho_\ell^2\mathrm J_\ell)\).
For \(\bar\varrho_\ell\le\mathrm c_0\),
\(\lambda-\mathrm C\varrho_\ell^2\mathrm J_\ell\ge\tfrac12\mathfrak j_\ell\mathrm J_\ell\),
and \(|\mathbf u(t_0)|\le\mathrm C\mathrm J_\ell\mathrm A\), so
\(|\mathbf u(t)|\le\mathrm C(\mathrm J_{\ell+1}^2\mathrm A/\mathrm J_\ell)[\varrho_\ell+\mathfrak x_\ell(t)]\).
For the forcing,
\(H_{zy}\mathbf g_y=H_{z,\ell+1}\mathbf g_{\ell+1}+\sum_{r\ge\ell+2}H_{z,r}\mathbf g_r\),
the sum being at most
\(\mathrm C\mathrm J_{\ell+2}^2\mathrm A\le\mathrm C\mathrm J_{\ell+1}^2\varrho_{\ell+1}^2\mathrm A\);
from \(\mathbf g_z=H_{zz}\bar{\mathbf x}_z+H_{zy}\bar{\mathbf x}_y\),
\(|\bar{\mathbf x}_z|\le\mathrm C\mathrm A(\varrho_\ell+\mathfrak x_\ell)\),
and \cref{eq:rg-formation-core-block-scales} then gives
\(\mathbf g_{\ell+1}=\mathrm J_{\ell+1}[\mathsf B_{\ell+1,\ell+1}^\top\mathsf b_{\ell+1}+O(\mathrm A(\bar\varrho_\ell+\mathfrak x_\ell))]\)
and
\(H_{zy}\mathbf g_y=\mathrm J_{\ell+1}^2[\mathfrak Y_\ell\mathsf b_{\ell+1}+O(\mathrm A(\bar\varrho_\ell+\mathfrak x_\ell))]\).
Writing \(\mathbf g_z=-H_{zz}^{-1}H_{zy}\mathbf g_y+\mathbf u\) with
\(H_{zz}^{-1}=\mathrm J_\ell^{-1}[(\mathsf B_{\ell\ell}^\top\mathsf B_{\ell\ell})^{-1}+O(\bar\varrho_\ell)]\)
and inserting in \(\tfrac12\mathbf g_z^\top H_{zz}^{-1}\mathbf g_z\), with
\(|\mathsf b_{\ell+1}|\le\mathrm C\mathrm A\), gives
\(\mathrm T\Delta_{\ell-1:\ell}=(\mathrm A^2\mathrm J_{\ell+1}^4/\mathrm J_\ell^3)
[\tfrac12|(\mathsf B_{\ell\ell}^\top\mathsf B_{\ell\ell})^{-3/2}\mathfrak Y_\ell\mathsf b_{\ell+1}/\mathrm A|^2
+O(\bar\varrho_\ell+\mathfrak x_\ell+\mathfrak x_\ell^2)]\), which is the
first line of \cref{eq:rg-formation-core-two-sided} by the definition of
\(\mathfrak y_\ell\).

\emph{Windows.} On \(I_\ell\),
\(\mathrm J_\ell(t-t_0)\ge\mathrm J_\ell(t_\ell-2t_{\ell-1})=k_L\log(1/\epsilon)(1-2\varrho_{\ell-1})\),
so \(\mathfrak x_\ell\to0\) uniformly on \(I_\ell\) under the stated
condition, and \(\mathrm J_{\ell+1}t\le2k_L\log(1/\epsilon)\varrho_\ell\to0\)
gives \(\mathsf b_{\ell+1}(t)/\mathrm A=B_{\ell+1}P_\ell\mathbf v_0+o(1)\), so
the quadratic \(\mathfrak y_\ell\) satisfies
\(\mathfrak y_\ell(\bar{\mathbf x}(t)/\mathrm A)=\mathfrak y_\ell(\mathbf v_0)+o(1)\).
For the lower bound,
\(|(\mathsf B_{\ell\ell}^\top\mathsf B_{\ell\ell})^{-3/2}\mathfrak Y_\ell\mathsf b_{\ell+1}|
\ge\|\mathsf B_{\ell\ell}\|^{-3}\sigma_{\min}(\mathfrak Y_\ell)|\mathsf b_{\ell+1}|\),
and \(|B_{\ell+1}P_\ell\mathbf v_0|^2=\sum_er_e(P_\ell\mathbf v_0)^2\) over
the \(N_\triangle/2^\ell\) bridges \(e\) of level \(\ell+1\) is at least
\(N_\triangle\mathrm m_L^2/2^{\ell+2}\) by the margin of
\Cref{lem:formation-preparation-margin}, which gives
\cref{eq:rg-formation-core-window}.

\emph{Power laws.} With \(\mathrm J_s=\epsilon^{-\mathrm Q_s}\),
\(\varrho_s=\epsilon^{\mathrm Q_s-\mathrm Q_{s+1}}\to0\),
\(k_L\log(1/\epsilon)\varrho_\ell\to0\), and on \(I_\ell\)
\[
\mathfrak x_\ell\le\epsilon^{\mathfrak j_\ell k_L(1-2\varrho_{\ell-1})/2-2(\mathrm Q_\ell-\mathrm Q_{\ell+1})},
\]
which tends to zero exactly under the stated threshold on \(k_L\); the
window statement then places \(\Delta_{\ell-1:\ell}\) between
\(\tfrac12\mathfrak y_\ell(\mathbf v_0)\) and \(2\mathfrak y_\ell(\mathbf v_0)\)
times \((\mathrm A^2/\mathrm T)\mathrm J_{\ell+1}^4/\mathrm J_\ell^3
=\epsilon^{3\mathrm Q_\ell-4\mathrm Q_{\ell+1}-\mathrm Q_{\mathrm T}}\) on
\(I_\ell\) for small \(\epsilon\). It is \(O(\epsilon^{\mathrm Q_\Delta})\)
on every early window exactly when
\(3\mathrm Q_\ell-4\mathrm Q_{\ell+1}-\mathrm Q_{\mathrm T}\ge\mathrm Q_\Delta\)
for every \(\ell\le L-1\), that is
\(\mathrm Q_\ell+\mathrm Q_{\mathrm T}+\mathrm Q_\Delta\ge\tfrac43(\mathrm Q_{\ell+1}+\mathrm Q_{\mathrm T}+\mathrm Q_\Delta)\),
and iterating from \(\mathrm Q_L=0\) gives
\(\mathrm Q_1+\mathrm Q_{\mathrm T}+\mathrm Q_\Delta\ge(4/3)^{L-1}(\mathrm Q_{\mathrm T}+\mathrm Q_\Delta)\),
with equality along the recursion; vanishing is the case
\(\mathrm Q_\Delta=0\) with strict inequality. \(\square\)

\section{The endogenous four-cycle}
\label{sec:rg-formation-cycle-proof}

\noindent\emph{Proof of \Cref{prop:rg-formation-endogenous-cycle}.}
\emph{Normalization, free energy, and descent.} Every factor of the joint is a
normalized kernel, so the joint has mass one, \(E=\mathbf 1\) has positive
probability, and every candidate source has positive prior and posterior mass
at every finite configuration. The reference density is
\((2\pi)^{-6}e^{-|\mathbf x|^2/2}\), the four likelihood factors contribute
\(4\mathrm T\log2\), and \cref{eq:rg-formation-evaluator-optimum} with
\(\sigma^2_{\mathrm{ev}}=1/2\) adds
\(\tfrac{\mathrm T}2\mathbf x^\top R_{\mathrm{ev}}\mathbf x\), so the scaled
negative ELBO is
\(\mathrm T\E_q\log(q/\pi_S)+\E_qU_S+\mathrm B_{\mathrm T}
=\mathrm T\KL(q\Vert q^\star)+\mathrm B_{\mathrm T}-\mathrm T\log Z_{\mathrm{opt}}\),
and the evidence is \(e^{-\mathrm B_{\mathrm T}/\mathrm T}Z_{\mathrm{phys}}\),
the physical energy omitting \(R_{\mathrm{ev}}\). Every \(U_S\) is a positive
definite quadratic form with Hessian at least \(\mathrm TI\), whose gradient
contains both the sending and the receiving appearances of each agent.
Between proposals the process is a linear diffusion with one of sixteen
matrices, and proposals are finitely many on bounded intervals, so it exists,
is pathwise unique with common random inputs, and does not explode. Since
\(q^\star(\mathbf x,S)\beta_{S'}(\mathbf x)
=q^\star_{\mathbf X}(\mathbf x)\beta_S\beta_{S'}\) is symmetric in \(S,S'\), the
dissipation identity \cref{eq:rg-formation-dissipation} at level zero, without
a parameter term, shows that \(\Fenergy_{\mathrm T}\) does not increase. A cycle
automorphism, with link signs transformed under reflections, preserves the
reference, the kernels, the rows, the kinetics, and the preparation class, so
the law at every time is invariant.

\emph{Matching sector.} Fix \(S^{(0)}\), with selected edges \(e=e_0\) and
\(f=e_2\); \(S^{(1)}\) follows by rotation. For each edge the covectors
\(\mathbf a^b_g,\mathbf a^m_g\) of \(r^b_g,r^m_g\) have Gram matrix
\(\begin{pmatrix}3&1\\1&3\end{pmatrix}\), and the two selected edges have
disjoint supports. The interaction Hessian of \(V_{S^{(0)}}\) is
\(H_M=2\sum_{g\in\{e,f\}}(\mathbf a^b_g\mathbf a^{b\top}_g+\mathbf a^m_g\mathbf a^{m\top}_g)\),
and the unit covectors
\begin{equation}
 n^+_g=\frac{\mathbf a^b_g+\mathbf a^m_g}{2\sqrt2},\qquad
 n^-_g=\frac{\mathbf a^b_g-\mathbf a^m_g}{2},\qquad g\in\{e,f\},
 \label{eq:rg-formation-cycle-normals}
\end{equation}
are orthonormal eigenvectors with eigenvalues \(8\) and \(4\); the kernel of
\(H_M\) has dimension eight. With \(\sigma^2_{\mathrm{ev}}=1/2\),
\(R_{\mathrm{ev}}n^+_g=0\), \(R_{\mathrm{ev}}n^-_g=4n^-_g\), and
\(R_{\mathrm{ev}}\) leaves the kernel invariant, so the held-matching drift
matrix \(K_M=H_M+\mathrm T(I+R_{\mathrm{ev}})\) has normal rates
\(8+\mathrm T\) (twice) and \(4+5\mathrm T\) (twice), all rates are at least
\(\mathrm T\), and the physical matrix has normal rates \(8+\mathrm T\) and
\(4+\mathrm T\). With \(\mathbf p=\sum_{g\in\{e,f\}}(\mathbf a^b_g+\mathbf a^m_g)\),
the vector \(\mathbf x^\circ+\mathbf p/4\) lies in the kernel of \(H_M\),
\(H_M\mathbf p=8\mathbf p\), and \(R_{\mathrm{ev}}\) annihilates \(\mathbf x^\circ\)
and \(\mathbf p\), so the deterministic held path from \(\mathbf x^\circ\) is
\begin{equation}
 \bar{\mathbf x}(t)=e^{-\mathrm Tt}\Bigl[\mathbf x^\circ+\frac{1-e^{-8t}}4\mathbf p\Bigr],
 \qquad
 r_{\mathrm{int}}(\bar{\mathbf x}(t))=-e^{-(8+\mathrm T)t},
 \qquad
 r_{\mathrm{ext}}(\bar{\mathbf x}(t))=e^{-\mathrm Tt}\Bigl[-\frac32+\frac{e^{-8t}}2\Bigr],
 \label{eq:rg-formation-cycle-center}
\end{equation}
in both channels, where \(r_{\mathrm{int}}\) is a selected and
\(r_{\mathrm{ext}}\) an exterior residual; along it \(\bar k=\bar m\), so the
evaluator force vanishes.

\emph{Episode.} Let \(\mathscr B\) be the event that
\(|\mathbf X_0-\mathbf x^\circ|<\varrho_\circ\), \(S_0=S^{(0)}\), and no proposal
occurs on \([0,1/4]\). The centered Gaussian component has expected squared
norm \(12\sigma_\circ^2=\varrho_\circ^2/4\), so it puts mass at least \(3/4\) on
the ball, and with the mixture weight \(1/2\), the bound
\(1+\mathfrak z\ge1/2\), the assignment mass \(1/16\), and the independent
clock, \(\Pr(\mathscr B)\ge\tfrac3{256}e^{-1/4}\). On \([0,1/4]\) the process is
the held Ornstein--Uhlenbeck process
\(\widetilde{\mathbf X}_t=e^{-K_Mt}\mathbf X_0+\mathbf Z_t\),
\(\mathbf Z_t=\sqrt{2\mathrm T}\int_0^te^{-K_M(t-\sigma)}dW_\sigma\), and afterwards
the two coincide until the first proposal returning another assignment. In an
eigenbasis of \(K_M\), integration by parts gives
\(\sup_{t\le t^\circ_*}|Z^k_t|\le2\sqrt{2\mathrm T}\sup_{t\le t^\circ_*}|W^k_t|\),
and \(\E\sup_{t\le t^\circ_*}|W^k_t|^2\le4t^\circ_*\) by Doob's \(L^2\)
maximal inequality: with \(W^*=\sup|W^k|\), optional stopping at the hitting
time of a level \(\mathrm y>0\) gives
\(\mathrm y\Pr(W^*\ge\mathrm y)\le\E[|W^k_{t^\circ_*}|\mathbf 1_{W^*\ge\mathrm y}]\),
and integrating against \(2 d\mathrm y\) with truncation gives the bound.
Summing twelve coordinates, \(\E\sup|\mathbf Z|^2\le384\mathrm Tt^\circ_*\), so
with \(\delta_\epsilon=\epsilon^2/(8\sqrt3)\), Chebyshev's inequality gives
\(\Pr(\sup_{t\le t^\circ_*}|\mathbf Z_t|>\delta_\epsilon)\le73728 \epsilon^4t^\circ_*\).
On \(\mathscr B\) and the complementary noise event, contraction gives
\(|\widetilde{\mathbf X}_t-\bar{\mathbf x}(t)|\le\varrho_\circ+\delta_\epsilon\).
Every residual covector has norm \(\sqrt3\), so on \(\mathscr B\) each
residual lies within \(\sqrt3\varrho_\circ\) of \(-1\) at time zero and each
internal belief relative entropy exceeds \((1-\sqrt3/100)^2/2>1/4\), which is
the entrance clause of (a); the restoring rates of (a) are those of \(K_M\)
computed above. Since \(\mathrm Tt^\circ_*<1/100\),
for \(1/4\le t\le t^\circ_*\) the selected residuals have magnitude below
\(e^{-2}+\sqrt3(\varrho_\circ+\delta_\epsilon)<1/5\) and the exterior ones above
\(e^{-\mathrm Tt^\circ_*}(3/2-e^{-2}/2)-\sqrt3(\varrho_\circ+\delta_\epsilon)>13/10\).
Each receiver's exterior cost therefore exceeds its partner cost by more than
one, its posterior row gives its exterior source probability at most
\(e^{-1/\mathrm T}\), and a proposal leaves \(S^{(0)}\) with probability at most
\(4e^{-1/\mathrm T}\); for the rate-one clock the expected number of departing
proposals up to \(t^\circ_*\) is at most \(4t^\circ_*e^{-1/\mathrm T}\), with
self-proposals unrestricted. For \(t\ge t_\epsilon\), every residual
annihilates the kernel of \(H_M\), so a selected residual is at most
\(e^{-8t}+\sqrt3\varrho_\circ e^{-4t}+\sqrt3\delta_\epsilon\le
\epsilon^{16}+\sqrt3\varrho_\circ\epsilon^8+\epsilon^2/8\le\epsilon^2\), the belief
relative entropy \((r^b)^2/2\) and the model relative entropy
\((r^m)^2/(2\sigma^2_{\mathrm{ev}})\) are at most \(\epsilon^4\), and every
exterior belief relative entropy exceeds \((13/10)^2/2>1/2\), the model ones
being larger. Let \(\mathscr A_t\) be the prefix event that the noise stays
within \(\delta_\epsilon\) up to \(t\) and every proposal in \([1/4,t]\) returns
\(S^{(0)}\); then \(\Pr(\mathscr A_{t^\circ_*}^c\mid\mathscr B)\le p_\epsilon\).
Since \(\epsilon^4t^\circ_*\) increases on \((0,1/100]\), the first term of
\(p_\epsilon\) is at most \(73728\cdot10^{-8}(2\log100+1)<0.008\), and
\(e^{-1/\mathrm T}<\mathrm T\) makes the second below \(10^{-14}\), so
\(p_\epsilon<1/2\) and
\(\Pr(\mathscr B\cap\mathscr A_{t^\circ_*})\ge\tfrac3{256}e^{-1/4}(1-p_\epsilon)\ge p_0\).
The stopping rule is prefix-measurable, and the same argument holds for
\(S^{(1)}\) and at the reflected center, which proves the probability bound
of (a).

\emph{Stationary comparison.} On each sheet the diffusion with drift
\(-\nabla U_S\) and covariance \(2\mathrm TI\) is reversible for
\(e^{-U_S/\mathrm T}d\mathbf x\), and the symmetric measure
\(q^\star(\mathbf x,S)\beta_{S'}(\mathbf x)\) makes the proposal step
reversible for \(q^\star\), so \(q^\star\) is invariant. Under \(q^\star\),
given \(S_0=M\), the configuration is \(\mathcal N(0,\mathrm TK_M^{-1})\), and
for a matched edge \(e\) with belief covector \(\mathbf a^b_e\),
\(K_M\succeq H_M\succeq2\mathbf a^b_e\mathbf a^{b\top}_e\), so
\(r^b_e=\mathbf a^b_e\cdot\mathbf x\) has variance
\(\mathrm T\mathbf a^{b\top}_eK_M^{-1}\mathbf a^b_e\le\mathrm T/2\) and
\(|r^b_e|>1/\sqrt2\) has probability at most \(2e^{-1/(2\mathrm T)}\). The
event of (a) requires \(S_0\) to be a matching whose matched pairs carry
internal belief relative entropy above \(1/4\) at time zero, that is
\(|r^b_e|>1/\sqrt2\) on a matched edge, so under the process started from
\(q^\star\) it has probability at most \(2\cdot2e^{-1/(2\mathrm T)}\le4e^{-1/(4\mathrm T)}\),
and \(4e^{-1/(4\mathrm T)}<p_0/2\) because \(\mathrm T\le10^{-16}\). This
completes (a).

\emph{Parent channel.} The covector of \(\sum_i\omega_i\) pairs to zero with
\(n^-_g\) and to \(-1/\sqrt2\) with \(n^+_g\), so it annihilates \(Z_M\) and the
loop is a function of \(P_{W_M}\mathbf x\); \(W_M\), of dimension nine, contains
the kernel of \(H_M\) and the fast common-normal direction
\((n^+_e+n^+_f)/\sqrt2\). Both \(H_M\) and \(R_{\mathrm{ev}}\) preserve \(W_M\)
and \(Z_M\), so on a matching fiber the hidden optimizing conditional is
\(\nu_{\mathrm T}=\mathcal N(0,\mathrm TK_Z^{-1})\) with
\(K_Z=\operatorname{diag}(8+\mathrm T,4+5\mathrm T,4+5\mathrm T)\), independent of
the retained state, and the physical one has precision
\(\operatorname{diag}(8+\mathrm T,4+\mathrm T,4+\mathrm T)\). The channel \(C\)
is one deterministic measurable kernel onto the disjoint union
\(\coprod_S\{S\}\times W_S\), with Lebesgue measure on each \(W_S\) and counting
measure on \(S\); its pushforwards of the joint, the physical posterior, the
optimizer, and the recognition law are normalized, it keeps every assignment,
and the chain rule of \Cref{thm:rg-exact-coarse-vfe} splits
\(\KL(q\Vert q^\star)\) into the retained term and the mean discarded
conditional term. This proves (b).

\emph{Discarded conditional free energy.} Under the held process the hidden
coordinates are \(z_t=e^{-K_Zt}z_0+\xi_t\) with
\(\xi_t\sim\mathcal N(0,\mathrm TK_Z^{-1}(I-e^{-2K_Zt}))\) independent of the
retained coordinates and of the initial data, and \(|z_0|\le\varrho_\circ\) on
\(\mathscr B\) because \(\mathbf x^\circ\) has zero hidden component. Given the
retained state, the hidden law is a mixture over \(z_0\) of these Gaussians,
and for one coordinate with rate \(\lambda_k\) and \(a=e^{-2\lambda_kt}\) the
relative entropy to the stationary law is
\(\tfrac12[-a-\log(1-a)+\lambda_ka z_0^2/\mathrm T]\). Convexity of relative
entropy in its first argument, \(4\le\lambda_k\le9\), and
\(-a-\log(1-a)\le a^2/(2(1-a))\) bound the mean conditional relative entropy of
the held law conditioned on \(\mathscr B\) by
\(\frac{9\varrho_\circ^2}{2\mathrm T}e^{-8t}+\frac{3}{4(1-e^{-8})}e^{-16t}\le\eta_\epsilon\)
for \(t\ge t_\epsilon\). To pass to the episode, let \(\widehat Q\) be this held
law, \(\mathscr R=\widehat Q_Y\otimes\nu_{\mathrm T}\), and \(Q_{\mathscr A}\),
\(Q_{\mathscr A^c}\) its conditionals on \(\mathscr A_t\) and its complement, of
probabilities \(1-p\) and \(p\le p_\epsilon\). The relative entropy of the joint
law of \((\mathbf 1_{\mathscr A_t},Y,Z)\) to
\(\mathrm{Bernoulli}(1-p)\otimes\mathscr R\), computed in both orders, gives
\((1-p)\KL(Q_{\mathscr A}\Vert\mathscr R)+p\KL(Q_{\mathscr A^c}\Vert\mathscr R)
=\KL(\widehat Q\Vert\mathscr R)+\mathbb I(\mathbf 1_{\mathscr A_t};(Y,Z))\), and the
mutual information is at most \(\mathcal S_2(p)\). Dropping the second term on
the left and the nonnegative retained-marginal term of
\(\KL(Q_{\mathscr A}\Vert\mathscr R)\), and using that \(\mathcal S_2\) increases
on \([0,1/2]\), bounds the mean discarded optimizing conditional relative
entropy of \(Q_{\mathscr A}\) by \(\bar\eta_\epsilon\); on
\(\mathscr B\cap\mathscr A_t\) the held and native processes coincide, so
\(Q_{\mathscr A}\) is the episode law. The physical parent conditions the
discarded coordinates also on the evaluator samples, so compare with the
normalized law
\(\Pi(dz\mid y) P_{\mathrm{ev}}(dm,dk\mid\mathbf x)
\exp[(\mathfrak d_y\cdot(k-m)-|\mathfrak d_y|^2/2)/(1-\sigma^2_{\mathrm{ev}})]\),
where \(\mathfrak d_y=\mathfrak Dy\) and \(\mathfrak d_z=\mathfrak Dz\) are
orthogonal; the multiplier has expectation one and depends only on retained
variables, so this law has the physical hidden conditional given
\((y,m,k)\). The recognition evaluator differs from it by the mean shift
\(\mathfrak d_z\), at relative entropy
\(|\mathfrak d_z|^2/(2(1-\sigma^2_{\mathrm{ev}}))\), and the chain rule in both
orders bounds the mean discarded physical conditional relative entropy by
\(\KL(q(dz\mid y)\Vert\Pi(dz\mid y))+\E_q|\mathfrak d_z|^2/(2(1-\sigma^2_{\mathrm{ev}}))
=\KL(q(dz\mid y)\Vert\nu_{\mathrm T})+\log\frac{4+5\mathrm T}{4+\mathrm T}\), where
\(\log\frac{4+5\mathrm T}{4+\mathrm T}\le\frac{4\mathrm T}{4+\mathrm T}\le\mathrm T\).
This proves (c).

\emph{Parent dynamics.} The inherited proposal kernel
\[
\beta_C((S,y),\cdot)=\int q^\star(d\mathbf x\mid C=(S,y))\sum_{S'}\beta_{S'}(\mathbf x)
\delta_{(S',P_{W_{S'}}\mathbf x)}
\]
is a probability kernel with the rate-one clock, reversible with respect to
\(C_\#q^\star\) because the symmetric measure
\(q^\star(d\mathbf x,S)\beta_{S'}(\mathbf x)\) is pushed through \((C,C)\); on a
matching stratum the retained diffusion is the exact projection of the held
process, and elsewhere \(C\) is the identity. On a matching fiber, with
\(z_+=((n^+_e+n^+_f)/\sqrt2)\cdot y\) and the hidden coordinates
\(z_1,z_2,z_3\) of \(Z_M\), the selected costs are
\(\mathcal D_e=(z_++z_1)^2+z_2^2\) and \(\mathcal D_f=(z_+-z_1)^2+z_3^2\), and each
exterior common residual \((r^b_g+r^m_g)/2\) annihilates \(Z_M\) and satisfies
\(\mathcal D_g\ge((r^b_g+r^m_g)/2)^2\). If both exterior common residuals have
magnitude at least \(\mathrm r_\circ\), each receiver chooses its exterior
source with probability at most
\(\exp[(\mathcal D_{\mathrm{in}}-\mathcal D_{\mathrm{out}})/\mathrm T]\), and the
Gaussian identity
\(\E e^{(a+Z)^2/\mathrm T}=(1-2\sigma^2/\mathrm T)^{-1/2}\exp[a^2/(\mathrm T(1-2\sigma^2/\mathrm T))]\)
with \(\sigma^2=\mathrm T/(8+\mathrm T)\) and \(\mathrm T/(4+5\mathrm T)\), whose
factors are at most \(\sqrt{4/3}\) and \(\sqrt2\), gives
\(\beta_C((S^{(0)},y),\{S'\ne S^{(0)}\})\le4\sqrt{8/3}\exp[-(\mathrm r_\circ^2-4z_+^2/3)/\mathrm T]\).
On the episode \(\mathrm r_\circ=13/10\) and \(|z_+|\le\epsilon^2\) on
\([t_\epsilon,t^\circ_*]\), so this is at most \(4\sqrt{8/3} e^{-1/\mathrm T}\).
Condition on \(\mathscr B\cap\mathscr A_{t_\epsilon}\) and let
\(\mu_{\mathrm{act}}\) be the law of \(C(\mathbf X_{t_\epsilon},S_{t_\epsilon})\)
under this prefix. The Brownian increments after \(t_\epsilon\) are independent
of the prefix, and the native future equals the held process until a
departure or a tube exit, whose conditional probability is at most
\(\Pr(\mathscr A^c_{t^\circ_*}\mid\mathscr B)/\Pr(\mathscr A_{t_\epsilon}\mid\mathscr B)
\le p_\epsilon/(1-p_\epsilon)\). The inherited parent started from
\(\mu_{\mathrm{act}}\) and driven by the same increments has the retained drift
and diffusion of the projected held process until its first departure; the
retained tube fails with conditional probability at most
\(p_\epsilon/(1-p_\epsilon)\), and on the tube the parent departs during the
unit interval with probability at most \(4\sqrt{8/3} e^{-1/\mathrm T}\). The
union bound and the triangle inequality give (d). \(\square\)

\noindent\emph{Proof of \Cref{prop:phil-seeded-contrast-formation}.}
\emph{Rows.} For a continuous row path \(\pi\) whose entries lie in
\([\mathrm c_0,1-\mathrm c_0]\) for some \(\mathrm c_0\in(0,1/2]\), construct the
kinetics pathwise as in the proof above: the configuration drift
\(-\nabla U_S\) does not involve the prior, and at a clock time \(\tau\)
receiver \(i\) takes source \(i+1\) exactly when a uniform variable on
\([0,1]\), independent and fresh for each clock time and receiver, is at most
its posterior row
\[
\beta_{i,i+1}(\mathbf x;\pi(\tau))
=\frac{\pi_{i,i+1}(\tau)e^{-\mathcal D_{e(i,i+1)}(\mathbf x)/\mathrm T}}
{\sum_{j=i\pm1}\pi_{ij}(\tau)e^{-\mathcal D_{e(i,j)}(\mathbf x)/\mathrm T}}
\]
at \(\mathbf x=\mathbf X_{\tau-}\). Write \(q^\pi_t\) for the law at time
\(t\). Drive the constructions for two such row paths \(\pi,\pi'\) by the same
initial data, Brownian motion, clock, and uniform variables. They coincide
until the first clock time at which the uniform variable of some receiver falls
between its posterior rows under \(\pi\) and \(\pi'\), and, given agreement
before a clock time
\(\tau\), this happens at \(\tau\) with conditional probability at most
\(\sum_i|\beta_{i,i+1}(\mathbf x;\pi(\tau))-\beta_{i,i+1}(\mathbf x;\pi'(\tau))|\)
at the common configuration. With \(\mathrm p=\pi_{i,i+1}\) and
\(\pi_{i,i-1}=1-\mathrm p\), the derivative of a posterior row in \(\mathrm p\)
is \(\beta_{i,i+1}(1-\beta_{i,i+1})/(\mathrm p(1-\mathrm p))
\le1/(4\mathrm c_0(1-\mathrm c_0))\le1/(2\mathrm c_0)\), and the clock has rate
one, so
\[
\bigl|q^\pi_t(S_i=j)-q^{\pi'}_t(S_i=j)\bigr|
\le\frac2{\mathrm c_0}\int_0^t\|\pi(s)-\pi'(s)\|_\infty ds,
\]
where \(\|\cdot\|_\infty\) is the maximum over receivers and sources. A
solution of \cref{eq:phil-source-learning} that starts uniform satisfies
\[
\pi_{ij}(t)-\tfrac12=\epsilon\int_0^te^{-\epsilon(t-s)}
\bigl(q_s(S_i=j)-\tfrac12\bigr)ds,
\qquad
\bigl|\pi_{ij}(t)-\tfrac12\bigr|\le\tfrac12\bigl(1-e^{-\epsilon t}\bigr)
\le\tfrac{\epsilon t}2,
\]
because \(0\le q_s(S_i=j)\le1\); each row sums to one, and on \([0,t_1]\) its
entries lie in \([\mathrm c_0,1-\mathrm c_0]\) with
\(\mathrm c_0=e^{-\epsilon t_1}/2\). Fix \(t_1>0\) and let \(\Gamma\) send a row
path \(\pi\) to the right-hand side of the first identity with \(q_s=q^\pi_s\).
It maps the closed set of continuous row paths on \([0,t_1]\) with
\(|\pi_{ij}(t)-1/2|\le(1-e^{-\epsilon t})/2\) into itself, and
\[
\|\Gamma\pi(t)-\Gamma\pi'(t)\|_\infty
\le\frac{2\epsilon t_1}{\mathrm c_0}\int_0^t\|\pi(s)-\pi'(s)\|_\infty ds,
\qquad t\le t_1,
\]
so the \(n\)th iterate of \(\Gamma\) contracts the supremum distance by the
factor \((2\epsilon t_1^2/\mathrm c_0)^n/n!\), and \(\Gamma\) has exactly one
fixed point in that set. Every solution lies in the set, solutions on nested
intervals agree, and \(s\mapsto q^\pi_s(S_i=j)\) is continuous because the clock
has rate one, so the fixed point is differentiable and solves
\cref{eq:phil-source-learning}. The rows are therefore a deterministic
continuous function of time, and the row bound holds. This proves (a).

\emph{Seed events.} The function
\(\epsilon t^\circ_*=\epsilon(2\log(1/\epsilon)+1)\) increases on
\((0,1/100]\), with value \((1+4\log10)/100<0.1022\) at \(1/100\), so
\(t^\circ_*<2/\epsilon\) and, by (a), every row entry lies in
\([0.4489,0.5511]\) on \([0,t^\circ_*]\), whence
\(\pi_{ik}(t)/\pi_{ij}(t)\le0.5511/0.4489<1.23\) for the two entries of every
row. For a sign
\(\sigma\in\{1,-1\}\) and a perfect matching \(M\in\{S^{(0)},S^{(1)}\}\), let
\(\mathscr B_{\sigma,M}\) be the event that
\(|\mathbf X_0-\sigma\mathbf x^\circ|<\varrho_\circ\), \(S_0=M\), and no proposal
occurs on \([0,1/4]\). The initial configuration, the initial assignment, and
the clock are independent, so
\[
\Pr(\mathscr B_{\sigma,M})=\frac{e^{-1/4}}{16}\mathcal N(\mathsf m_0,I_{12})
\bigl\{|\mathbf x-\sigma\mathbf x^\circ|<\varrho_\circ\bigr\},
\]
and the four events are disjoint because \(|\mathbf x^\circ|=2>\varrho_\circ\).

\emph{Episode.} Condition on \(\mathscr B_{1,S^{(0)}}\). The configuration drift
does not involve the rows, and the rows are deterministic, so the paragraph
\emph{Episode} of the proof above applies line by line from every starting
point in the ball, whatever the law of \(\mathbf X_0\) on it: the held
Ornstein--Uhlenbeck process, whose noise bound involves only the Brownian
motion; the contraction to the held path \cref{eq:rg-formation-cycle-center};
the residual bounds below \(1/5\) and above \(13/10\) on \([1/4,t^\circ_*]\),
which make each receiver's exterior cost exceed its partner cost by more than
one; and the relative-entropy bounds for \(t\ge t_\epsilon\), all deterministic
on the noise event. The rows enter only the exterior source probability: for a
receiver \(i\) with partner \(j\) and exterior source \(k\),
\[
\beta_{ik}(\mathbf x;\pi(t))\le\frac{\pi_{ik}(t)}{\pi_{ij}(t)}
\exp\Bigl[-\frac{\mathcal D_{e(i,k)}(\mathbf x)-\mathcal D_{e(i,j)}(\mathbf x)}{\mathrm T}\Bigr]
<1.23 e^{-1/\mathrm T}
\]
in place of \(e^{-1/\mathrm T}\). A proposal therefore leaves \(S^{(0)}\) with
probability at most \(5e^{-1/\mathrm T}\), the expected number of departing
proposals up to \(t^\circ_*\) is at most \(5t^\circ_*e^{-1/\mathrm T}\), and the
prefix event \(\mathscr A_{t^\circ_*}\) of that paragraph fails with conditional
probability at most
\[
\tilde p_\epsilon=73728 \epsilon^4t^\circ_*+5t^\circ_*e^{-1/\mathrm T}
<0.008+10^{-14}<\tfrac12,
\]
by the same monotonicity of \(\epsilon^4t^\circ_*\) and
\(e^{-1/\mathrm T}<\mathrm T\). On
\(\mathscr B_{1,S^{(0)}}\cap\mathscr A_{t^\circ_*}\), the block \(\{0,1\}\), of
size two and joined to its complement by \(e_1\) and \(e_3\), carries on its
internal edge \(e_0\) a belief relative entropy above
\((1-\sqrt3\varrho_\circ)^2/2>1/4\) at time zero, because
\(|\mathbf X_0-\mathbf x^\circ|<\varrho_\circ\) and every residual covector
has norm \(\sqrt3\); it has internal belief and model relative entropies at
most \(\epsilon^4\) on \([t_\epsilon,t^\circ_*]\), an interval of length one
inside \([0,2/\epsilon]\), and both boundary belief relative entropies above
\(1/2\) on \([1/4,t^\circ_*]\), so contrast formation occurs.

\emph{The other seeds.} A cycle automorphism \(g\), with link signs transformed
under reflections, acts on configurations by an orthogonal signed permutation
\(\mathsf O_g\) with \(U_{gS}(\mathsf O_g\mathbf x)=U_S(\mathbf x)\) and
\(\mathcal D_{e(gi,gj)}(\mathsf O_g\mathbf x)=\mathcal D_{e(i,j)}(\mathbf x)\).
With \((g\pi)_{gi,gj}=\pi_{ij}\), it therefore maps the kinetics with row path
\(\pi\) from \((\mathbf x,S)\) to the kinetics with row path \(g\pi\) from
\((\mathsf O_g\mathbf x,gS)\), driven by the Brownian motion \(\mathsf O_gW\)
and by the receiver uniform variables relabeled by \(g\), the variable of
receiver \(i\) serving receiver \(gi\) and being replaced by one minus itself
when \(g\) is a reflection. Indeed a reflection sends the source \(i+1\) of
receiver \(i\) to the source \(gi-1\) of receiver \(gi\), and the image
receiver takes it exactly when the original variable is at most
\(\beta_{i,i+1}(\mathbf x;\pi)=\beta_{gi,gi-1}(\mathsf O_g\mathbf x;g\pi)
=1-\beta_{gi,gi+1}(\mathsf O_g\mathbf x;g\pi)\), which is the rule of the
image kinetics applied to one minus that variable, except on the null event
that the variable equals its row; under the rotation the sources are carried
along and the variable is kept. The relabeled family is again a family of
independent uniform variables independent of the other inputs, so the image
process has the law of the kinetics with rows \(g\pi\), and \(g\pi\) has the
same ratio bound. The rotation \(i\mapsto i+1\) fixes
\(\mathbf x^\circ\) and maps \(S^{(0)}\) to \(S^{(1)}\), the reflection
\(i\mapsto-i\) maps \((\mathbf x^\circ,S^{(0)})\) to
\((-\mathbf x^\circ,S^{(1)})\), and their composition maps
\((\mathbf x^\circ,S^{(0)})\) to \((-\mathbf x^\circ,S^{(0)})\), so the
conditional bound holds on the other three seed events, with the images of the
block. Summing over the four disjoint
events,
\[
\Pr(\text{contrast formation})\ge(1-\tilde p_\epsilon)\frac{2e^{-1/4}}{16}
\sum_\pm\mathcal N(\mathsf m_0,I_{12})
\bigl\{|\mathbf x\mp\mathbf x^\circ|<\varrho_\circ\bigr\},
\]
which gives the first bound of (b) because \(\tilde p_\epsilon<1/2\). The
Gaussian density on the ball about \(\mathbf x^\circ\) is at least
\((2\pi)^{-6}\exp[-(|\mathbf x^\circ-\mathsf m_0|+\varrho_\circ)^2/2]\), the ball
has volume \(\pi^6\varrho_\circ^{12}/6!\), and \((2\pi)^{-6}\pi^6/6!=1/46080\),
which gives the second bound after dropping the ball about
\(-\mathbf x^\circ\). Frozen rows have ratio one, so the argument and the bound
hold without learning. The sub-event named in (b) is the event of part (a) of
the proposition above with the block of the held matching. For fixed rows
\(\pi\) the proposal step is reversible for \(q^\star_\pi\) by the same
symmetry \(q^\star_\pi(\mathbf x,S)\beta_{S'}(\mathbf x;\pi)
=q^\star_{\pi,\mathbf X}(\mathbf x)\beta_S(\mathbf x;\pi)\beta_{S'}(\mathbf x;\pi)\),
the sheet diffusions do not involve the rows, and the conditional law
\(\mathcal N(0,\mathrm TK_M^{-1})\) of \(q^\star_\pi\) given \(S_0=M\) does not
depend on \(\pi\), so the paragraph \emph{Stationary comparison} above gives
the invariance and the bound \(4e^{-1/(4\mathrm T)}\) for every \(\pi\). This
proves (b).

\emph{Invariant initial laws.} Since \(\mathsf O_g\) is orthogonal, \(g\) maps
\(\mathcal N(\mathsf m_0,I_{12})\) to \(\mathcal N(\mathsf O_g\mathsf m_0,I_{12})\),
so the initial law is invariant exactly when \(\mathsf m_0\) is fixed by the
rotation and the reflection. The rotation fixes exactly the means whose belief,
model, and link blocks are constant, and the reflection reverses every edge and
hence the sign of a constant link block, which leaves
\((\mathrm c_1\mathbf 1_4,\mathrm c_2\mathbf 1_4,0)\). For such a mean, since
the uniform assignment law and the uniform rows are also invariant, \(g\) maps
the recognition law \(q_t\) of the rows \(\pi\) to the recognition law
\(g_\#q_t\) of the rows \(g\pi\) from the same initial law, by the pathwise
image of the paragraph \emph{The other seeds}, whose relabeled driving inputs
have the law of the original ones; and \(g\pi\) solves
\cref{eq:phil-source-learning} with \(g_\#q_t(S_{gi}=gj)=q_t(S_i=j)\), so
\((g\pi,g_\#q)\) is a second solution from the same data and uniqueness in (a)
gives \(g\pi=\pi\). The reflection \(j\mapsto2i-j\) fixes receiver \(i\) and
exchanges its sources \(i\pm1\), so \(\pi_{i,i+1}=\pi_{i,i-1}=1/2\). This
proves (c). \(\square\)

\section{Depth one under exclusive source gating}
\label{sec:rg-formation-depth-one-proof}

\noindent\emph{Proof of \Cref{prop:rg-gated-depth-one}.}
Throughout, \(F=\operatorname{att}(S)\) and \(W=W_F\) once a sheet \(S\) is
fixed, \(P_W\) is the orthogonal projection onto \(W\), and
\(\mathsf Q_W\), \(\mathsf Q_{W^\perp}\) are matrices whose columns form
orthonormal bases of \(W\) and of \(W^\perp\). Since \(H_S\) is the sum of the
positive semidefinite maps \(\mathbf a_{e(i,S_i)}^\top\mathbf a_{e(i,S_i)}\),
its kernel is \(\bigcap_{e\in F}\ker\mathbf a_e=W_F\) and its range is
\(W_F^\perp\), the span of the rows of \(\mathbf a_F\), so
\(\mathrm k_S=\operatorname{rank}\mathbf a_F\). For every edge the Gram
matrix of the two rows of \(\mathbf a_e\) is
\(\bigl(\begin{smallmatrix}3&1\\1&3\end{smallmatrix}\bigr)\), as in
\cref{eq:rg-formation-cycle-normals}, so \(\|\mathbf a_e\|=2\).

\emph{(a) Two bands.} Let \(M=H_S+\mathrm T\Omega\). By Weyl's inequality
\citep[Sec.~4.3]{HornJohnson2013}, with eigenvalues in increasing order,
\(\lambda_k(H_S)+\mathrm T\lambda_{\min}(\Omega)\le\lambda_k(M)
\le\lambda_k(H_S)+\mathrm T\|\Omega\|\) for every \(k\). The first
\(2N+|E|-\mathrm k_S\) eigenvalues of \(H_S\) vanish, the remaining
\(\mathrm k_S\) lie in \([\mathrm h_S,\|H_S\|]\), and
\(\lambda_{\min}(\Omega)=\|\Omega^{-1}\|^{-1}\). The range was identified
above, and \(\pi\) does not enter \(M\).

\emph{(b) Gating wells.} For each receiver, the row sums to one and every
weight is at least \(\min_j\pi_{ij}\), so
\(\min_j\pi_{ij}e^{-\min_j\mathcal D_{e(i,j)}/\mathrm T}
\le\sum_j\pi_{ij}e^{-\mathcal D_{e(i,j)}/\mathrm T}
\le e^{-\min_j\mathcal D_{e(i,j)}/\mathrm T}\); taking \(-\mathrm T\log\) and
summing over \(i\) gives \cref{eq:rg-gated-well-depth}. Since \(V_S\) is a sum
over receivers and \(\pi_S\) a product,
\[
\sum_S\pi_Se^{-V_S/\mathrm T}=\prod_i\sum_{j\in\mathcal N_i}\pi_{ij}
e^{-\mathcal D_{e(i,j)}/\mathrm T}=e^{-V_\pi/\mathrm T},
\]
which is the configuration marginal up to the reference factor. If every edge of a
covering set \(F\) has zero residual at \(\mathbf x\), every receiver has an
option of zero cost, so \(V_{\min}(\mathbf x)=0\); a common translation of the
belief and model means of the agents of one component of \((V,F)\), with the
links fixed, changes no residual of an edge of \(F\), because the two
endpoints of such an edge either both move or both stay, so \(V_{\min}\)
remains zero.

\emph{(c) Rank formula.} In the coordinates \((\bar k,\mathfrak d,\omega)\)
with \(\mathfrak d=\bar k-\bar m\), the rows of \(\mathbf a_e\) for
\(e:j\to i\) are \(r^b_e=\bar k_i-\bar k_j-\omega_e\) and
\(r^m_e=r^b_e-(\mathfrak d_i-\mathfrak d_j)\); replacing the second by
\(r^m_e-r^b_e=-(\partial\mathfrak d)_e\) is an invertible change within the
rows of \(e\). The rows \(r^b_e\), \(e\in F\), read only \((\bar k,\omega)\)
and each has coefficient \(-1\) on its own link coordinate and zero on every
other, so they are linearly independent. The rows \((\partial\mathfrak d)_e\),
\(e\in F\), read only \(\mathfrak d\) and form the oriented incidence map of
\((V,F)\) on \(\R^V\), whose kernel consists of the functions constant on each
component, so their rank is \(N-\operatorname{cc}(F)\). The two families act
on disjoint coordinate blocks, so
\(\operatorname{rank}\mathbf a_F=|F|+N-\operatorname{cc}(F)\).

\emph{Minimum edge covers.} Let \(F=\operatorname{att}(S)\). Every vertex
selects an incident edge, so \(F\) covers \(V\) and every component of
\((V,F)\) contains an edge; one edge from each component is a matching, so
\(\operatorname{cc}(F)\le\mathfrak n\), and a connected graph on
\(\mathrm c\) vertices has at least \(\mathrm c-1\) edges, so
\(|F|\ge N-\operatorname{cc}(F)\). Hence
\(\mathrm k_S\ge2(N-\operatorname{cc}(F))\ge2(N-\mathfrak n)\), with
equality exactly when \(F\) is a forest with \(\mathfrak n\) components and
therefore \(N-\mathfrak n\) edges. The same two inequalities show that every
edge cover has at least \(N-\mathfrak n\) edges, and a maximum matching
together with one incident edge at each unmatched vertex, which exists
because no vertex is isolated, is an edge cover with exactly
\(N-\mathfrak n\) edges, since the other endpoint of each added edge is
matched, or else the matching could be enlarged. An edge cover with
\(N-\mathfrak n\) edges has at least \(\mathfrak n\) components and hence
exactly \(\mathfrak n\), and is a forest; so the covering forests with
\(\mathfrak n\) components are exactly the minimum edge covers. In a
minimum edge cover every edge has an endpoint of degree one, since otherwise
removing it would leave a smaller cover, and a tree in which every edge has
a leaf endpoint is a star, because two vertices of degree at least two would
be joined by a path whose first edge has no leaf endpoint. Conversely, a
star with center \(z\) is attended when every leaf selects \(z\) and \(z\)
selects one leaf, which the source sets allow, so every minimum edge cover
is some \(\operatorname{att}(S)\). Adding to \(F\) an edge that joins two
components raises \(|F|\) by one and lowers \(\operatorname{cc}(F)\) by one,
hence raises the rank by two.

\emph{Sheet weights.} Since \(H_S/\mathrm T+\Omega\succ0\),
\(Z_S=(2\pi)^{(2N+|E|)/2}\det(H_S/\mathrm T+\Omega)^{-1/2}\). In the basis
\((\mathsf Q_{W^\perp},\mathsf Q_W)\) the matrix \(H_S/\mathrm T+\Omega\) has
the blocks \(H_1/\mathrm T+\Omega_{11}\), \(\Omega_{10}\), \(\Omega_{01}\),
\(\Omega_{00}\), with \(H_1=\mathsf Q_{W^\perp}^\top H_S\mathsf Q_{W^\perp}\succeq\mathrm h_SI\)
and \(\Omega_{00}=\Omega|_W\). Eliminating the invertible block
\(\Omega_{00}\) by block row and column operations gives
\(\det(H_S/\mathrm T+\Omega)=\det\Omega_{00}\det(H_1/\mathrm T+\Sigma)\) with
\(\Sigma=\Omega_{11}-\Omega_{10}\Omega_{00}^{-1}\Omega_{01}\), and
\(0\prec\Sigma\preceq\Omega_{11}\preceq\|\Omega\|I\). Then
\(\mathrm T\Sigma\preceq(\mathrm T\|\Omega\|/\mathrm h_S)H_1\), so
\(\det H_1\le\det(H_1+\mathrm T\Sigma)\le(1+\mathrm T\|\Omega\|/\mathrm h_S)^{\mathrm k_S}\det H_1\),
and \(\det H_1=\operatorname{pdet}H_S\). Since
\(\det(H_1/\mathrm T+\Sigma)=\mathrm T^{-\mathrm k_S}\det(H_1+\mathrm T\Sigma)\),
the determinant lies between \(\mathrm T^{-\mathrm k_S}\mathrm z_S^{-2}\) and
\((1+\mathrm T\|\Omega\|/\mathrm h_S)^{\mathrm k_S}\mathrm T^{-\mathrm k_S}\mathrm z_S^{-2}\),
which is \cref{eq:rg-gated-sheet-weight}. The same block form shows that
under \(q_\Omega(\cdot\mid S)\), which is
\(\mathcal N(0,(H_S/\mathrm T+\Omega)^{-1})\), the coordinates
\(y=\mathsf Q_W^\top\mathbf x\) given \(z=\mathsf Q_{W^\perp}^\top\mathbf x\)
are Gaussian with precision \(\Omega_{00}\), the \((W,W)\) block of the joint
precision, which does not depend on \(\mathrm T\).

\emph{Merge cost.} \(q_\Omega(S)=\pi_SZ_S/\sum_{S'}\pi_{S'}Z_{S'}
\le\pi_SZ_S/\sum_{S'\in\mathcal M}\pi_{S'}Z_{S'}\). The upper bound of
\cref{eq:rg-gated-sheet-weight} for \(S\) and its lower bound for each
\(S'\in\mathcal M\), where \(\mathrm k_{S'}=2(N-\mathfrak n)\) and
\(\mathrm h_{S'}\ge\mathrm h_*\), give
\(q_\Omega(S)\le\mathrm T^{(\mathrm k_S-2(N-\mathfrak n))/2}
(1+\mathrm T\|\Omega\|/\mathrm h_*)^{N-\mathfrak n}
\pi_S\mathrm z_S/\sum_{S'\in\mathcal M}\pi_{S'}\mathrm z_{S'}\). If
\(\operatorname{cc}(\operatorname{att}(S))\le\mathfrak n-m\), then
\(\mathrm k_S\ge2(N-\operatorname{cc}(\operatorname{att}(S)))\ge2(N-\mathfrak n)+2m\);
with \(\mathrm T\le1\), summing over such \(S\) gives
\cref{eq:rg-gated-merge-cost}.

\emph{(d) Block law.} The sheets with
\(\operatorname{cc}(\operatorname{att}(S))<\mathfrak n\) carry mass at most
\(\mathrm K_1\mathrm T\) by \cref{eq:rg-gated-merge-cost}. Fix a sheet with
\(\operatorname{cc}(F)=\mathfrak n\). No component of \((V,F)\) contains
two disjoint edges, since replacing that component's edge in the matching of
one edge per component would give \(\mathfrak n+1\) disjoint edges. A
connected graph in which no two edges are disjoint is a star or a triangle:
fix an edge \(uv\); every other edge meets \(\{u,v\}\); if there are edges
\(vw\) and \(uw'\) with \(w\ne u\) and \(w'\ne v\), then \(w=w'\), and a
further edge would have to meet \(\{u,v\}\) and would then be disjoint from
\(vw\) or from \(uw\), so the graph is the triangle \(uvw\); otherwise all
edges share one vertex. Replacing each triangle component by a two-edge path
on the same three vertices gives a covering forest of stars with
\(\mathfrak n\) components, a minimum edge cover by (c), so the partition
of \(V\) into the components of \(F\) is a minimum-edge-cover partition.
The coherent-block partition \(\mathfrak P_{\mathrm a}(\mathbf x)\) differs
from it only if some attended edge has \(|\mathbf a_e\mathbf x|>\mathrm a\) or
some edge joining two components of \(F\) has
\(|\mathbf a_e\mathbf x|\le\mathrm a\). Under \(q_\Omega(\cdot\mid S)\) the
configuration is \(\mathcal N(0,\mathrm T(H_S+\mathrm T\Omega)^{-1})\). For
\(e\in F\), \(M=H_S+\mathrm T\Omega\succeq\mathbf a_e^\top\mathbf a_e\), so
\(\|\mathbf a_eM^{-1/2}\|\le1\) and the covariance
\(\mathrm T\mathbf a_eM^{-1}\mathbf a_e^\top\) of \(\mathbf a_e\mathbf x\) is at
most \(\mathrm TI_2\); each of its two centered Gaussian coordinates has
variance at most \(\mathrm T\), and \(|\mathbf a_e\mathbf x|>\mathrm a\)
forces one of them beyond \(\mathrm a/\sqrt2\), an event of probability at
most \(2\cdot2e^{-\mathrm a^2/(4\mathrm T)}\). For an edge \(e\notin F\)
joining two components, \(\operatorname{rank}\mathbf a_{F\cup\{e\}}
=\operatorname{rank}\mathbf a_F+\operatorname{rank}(\mathbf a_e|_{W})\),
because \(\ker\mathbf a_{F\cup\{e\}}=W\cap\ker\mathbf a_e\); by (c) the
restriction has rank two, so \(\mathbf a^W_e=\mathbf a_e\mathsf Q_W\) has two nonzero
singular values, the least of which is at least \(\mathrm s_*\). Writing
\(\mathbf x=\mathsf Q_{W^\perp}z+\mathsf Q_Wy\), the conditional law of \(y\)
given \(z\) has precision \(\Omega_{00}\preceq\|\Omega\|I\), so \(\mathbf a^W_ey\) is
Gaussian in \(\R^2\) with covariance
\(\mathbf a^W_e\Omega_{00}^{-1}\mathbf a^{W\top}_e\succeq\|\Omega\|^{-1}\mathbf a^W_e\mathbf a^{W\top}_e\succeq(\mathrm s_*^2/\|\Omega\|)I_2\)
and density at most \((2\pi)^{-1}\|\Omega\|/\mathrm s_*^2\). The event
\(|\mathbf a_e\mathbf x|\le\mathrm a\) places \(\mathbf a^W_ey=\mathbf a_e\mathbf x-\mathbf a_e\mathsf Q_{W^\perp}z\)
in a disc of radius \(\mathrm a\) about a point determined by \(z\), so its
conditional probability is at most \(\pi\mathrm a^2\cdot\|\Omega\|/(2\pi\mathrm s_*^2)\).
A union bound over at most \(|E|\) edges of each kind, integrated over the
sheets, gives \cref{eq:rg-gated-block-law}. \(\square\)

\section{Proof of Pairing Selection}
\label{sec:phil-pairing-selection-proof}

\noindent\emph{Proof of \Cref{conj:phil-contrast-formation}.}
Write \(p_i=\pi_{i,i+1}\), encode the assignment by
\(\sigma_i=+1\) when receiver \(i\) selects \(i+1\), and put
\(\varepsilon=\epsilon_{\mathrm L}\). The state remains
\((\mathbf x,\sigma)\in\R^{12}\times\{\pm1\}^4\), with the shared edge
costs of \cref{eq:rg-formation-cycle-cost}. The multiplicity of edge
\(e\) is
\(m_e(\sigma)=1_{\{\sigma_e=+1\}}+1_{\{\sigma_{e+1}=-1\}}\).
In this proof define
\[
 \Omega=I+R_{\mathrm{ev}},\qquad
 \mathsf A_\sigma=H_\sigma+\mathrm T\Omega,\qquad
 H_\sigma=\nabla^2V_\sigma,\qquad
 \mathsf A_\sigma\succeq\mathrm TI.
\]
Let \(\mathsf w_\sigma=\det(\mathsf A_\sigma/\mathrm T)^{-1/2}\),
\(\mathsf Z(p)=\sum_\sigma\pi_\sigma(p)\mathsf w_\sigma\), and
\[
 \mu_p(\mathbf x,\sigma)
 =\frac{\pi_\sigma(p)}{(2\pi)^6\mathsf Z(p)}
   \exp\left[-\frac{\mathbf x^\top\mathsf A_\sigma\mathbf x}
                    {2\mathrm T}\right],\qquad
 \pi_\sigma(p)=\prod_i p_i^{1_{\{\sigma_i=+1\}}}
                         (1-p_i)^{1_{\{\sigma_i=-1\}}}.
\]
Each held-sheet diffusion preserves its Gaussian factor, and a clock ring
resamples \(\sigma\) from its conditional law under \(\mu_p\). Thus
\(\mu_p\) is invariant for the frozen rows, including boundary rows on their
supported sheets. Set
\begin{equation}
 \bar\beta_i(p)=\mu_p(\sigma_i=+1),\qquad
 \mathfrak g_i(p)=\partial_{p_i}\log\mathsf Z(p),\qquad
 \mathscr V_i(p)=\bar\beta_i(p)-p_i
               =p_i(1-p_i)\mathfrak g_i(p).
\label{eq:phil-selection-reduced-field}
\end{equation}
The last identity follows by differentiating the product prior. Since
\(\mathsf Z\) is a positive multilinear polynomial on the closed cube,
\(\mathfrak g\) and \(\mathscr V\) extend analytically to a neighborhood of
that cube and have bounded derivatives there. The well-posedness argument
for \cref{prop:phil-seeded-contrast-formation} applies at every fixed
\(\mathrm T>0\); it gives deterministic continuously differentiable rows and
\begin{equation}
 \min_i\{p_i(t),1-p_i(t)\}\ge\tfrac12e^{-\varepsilon t},\qquad
 \|p(t)-\tfrac12\|_\infty\le\tfrac12\varepsilon t.
\label{eq:phil-selection-row-bounds}
\end{equation}
All constants below may depend on the fixed temperature. Constants in
bounds on a fixed initial law may also depend on its mean, whereas the
window constant in the conclusion will depend only on the temperature.

\subsection*{Frozen convergence and the initial imprint}

Here \(\operatorname{Ent}_\nu(g)=\int g\log(g/\int g d\nu)d\nu\)
for nonnegative \(g\), and total variation is \(\sup_A|q(A)-\nu(A)|\).
For the time-one frozen kernel, the function
\(\mathsf V(\mathbf x)=|\mathbf x|^2\) satisfies
\begin{equation}
 \E_{\mathbf x}|\mathbf X_t|^2
 \le e^{-2\mathrm Tt}|\mathbf x|^2+12(1-e^{-2\mathrm Tt}),\qquad
 P_1^p\mathsf V\le e^{-2\mathrm T}\mathsf V+12.
\label{eq:phil-selection-moment}
\end{equation}
Indeed, jumps leave \(\mathbf X\) unchanged, and the diffusion generator
applied to its squared norm is at most \(24\mathrm T-2\mathrm T|\mathbf X|^2\).
Choose \(R>24/(1-e^{-2\mathrm T})\). For any frozen \(p\), some assignment
\(\sigma_*(p)\) has \(\pi_{\sigma_*}\ge1/16\). Restrict the rate-one
clock to exactly one ring in \([0,1/2]\) and none afterward. Its probability
is \(e^{-1}/2\). Starting with \(|\mathbf x|\le\sqrt R\), the first held
Gaussian motion is in the ball of radius
\(r'=\sqrt R+\sqrt{24\mathrm T}\) at that ring with probability at least
\(1/2\). On the ball, the probability to redraw \(\sigma_*\) is at least
\(16^{-1}e^{-v(r')/\mathrm T}\), where
\(v(r')=\max_{|\mathbf z|\le r',\sigma}V_\sigma(\mathbf z)<\infty\).
The subsequent held Gaussian transition lasts between \(1/2\) and \(1\).
Its means are bounded and its covariance matrices have common positive
lower and finite upper bounds. Its density consequently has a positive
common lower bound on the unit ball. Hence, for some \(\alpha_{\mathrm T}>0\)
independent of \(p\),
\[
 P_1^p((\mathbf x,\sigma),\cdot)
 \ge\alpha_{\mathrm T}\operatorname{Unif}(B_1)
                         \otimes\delta_{\sigma_*(p)}
 \quad\text{when }|\mathbf x|^2\le R.
\]
These are the drift and minorization hypotheses of
\citet[Assumptions 1--2 and Theorems 1.2--1.3]{HairerMattingly2011}.
The explicit dependence in their Theorem 1.3 makes the resulting constants
uniform in the frozen rows, even though the minorizing measure depends on
\(p\). Thus \(\mu_p\) is the unique invariant law and
\begin{equation}
 |P_t^p\phi(\mathbf x,\sigma)-\mu_p\phi|
 \le\mathrm K_{\mathrm T}e^{-\lambda_{\mathrm T}t}
       (1+|\mathbf x|^2)\|\phi\|_\infty.
\label{eq:phil-selection-harris}
\end{equation}
Intermediate times follow by applying the skeleton estimate to
\(P_r^p\phi\), \(0\le r<1\).

Put \(\mathbf v=(1,-1,1,-1)\),
\(f_{\mathbf v}(\sigma)=\frac12\sum_i\mathbf v_i\sigma_i\), and define
\begin{equation}
 \mathsf I_i(\mathsf m_0)
 =\int_0^\infty[q_t^0(\sigma_i=+1)-\tfrac12]dt,\qquad
 \mathsf j_{\mathrm T}(\mathsf m_0)
 =\mathbf v\cdot\mathsf I(\mathsf m_0).
\label{eq:phil-selection-imprint}
\end{equation}
At uniform rows, the stationary law is invariant under cycle rotations and
reflections, so each stationary row is \(1/2\) and
\(\mu_{1/2}f_{\mathbf v}=0\). Estimate
\cref{eq:phil-selection-harris} proves absolute convergence of these
integrals. It also gives a Poisson integral
\begin{equation}
 \mathfrak R_{\mathrm T}(\mathbf x,\sigma)
 =\int_0^\infty P_t^{1/2}f_{\mathbf v}(\mathbf x,\sigma)dt,\qquad
 |\mathfrak R_{\mathrm T}(\mathbf x,\sigma)|
 \le\mathrm K_{\mathrm T}(1+|\mathbf x|^2).
\label{eq:phil-selection-poisson}
\end{equation}
Convolution of this integral with the initial Gaussian expresses
\(\mathsf j_{\mathrm T}\) as an entire function of its complexified mean:
for every compact set of complex means, all differentiated integrands are
bounded by an integrable polynomial times a Gaussian. The same holds for
\(\mathsf I\). A rotation by one changes the sign of
\(\mathsf j_{\mathrm T}\).

For \(t\le1/(10\varepsilon)\), all row entries are in
\([\mathrm c_0,1-\mathrm c_0]\), \(\mathrm c_0=0.45\). A posterior receiver probability has derivative
at most \(1/[4\mathrm c_0(1-\mathrm c_0)]\le1/(2\mathrm c_0)\) with respect to its prior entry.
Coupling the common clocks and their four posterior rows therefore gives
\[
 \|q_t^\varepsilon-q_t^0\|_{\mathrm{TV}}
 \le\frac2{\mathrm c_0}\int_0^t\|p(s)-\tfrac12\|_\infty ds
 \le\frac{\varepsilon t^2}{2\mathrm c_0}.
\]
Variation of constants in the learning equation, followed by
\cref{eq:phil-selection-harris}, gives
\begin{equation}
 \left|p(t)-\tfrac12-\varepsilon\mathsf I(\mathsf m_0)\right|_\infty
 \le\mathrm K_{\mathrm T,\mathsf m_0}
       (\varepsilon e^{-\lambda_{\mathrm T}t}+\varepsilon^2t)
       +\frac{\varepsilon^2t^3}{6\mathrm c_0}.
\label{eq:phil-selection-imprint-expansion}
\end{equation}
In particular, at \(t_1=\varepsilon^{-1/4}\),
\(p(t_1)-1/2=\varepsilon\mathsf I(\mathsf m_0)+o(\varepsilon)\).

\subsection*{Nonvanishing of the matching imprint}

Take \(\mathbf u_\circ=(0_4,0_4,(0,1,0,-1))\), whose two residual vectors
are \((0,-1,0,1)\). In the noiseless switched motion the subspace
\(\bar k=\bar m\) is invariant. After removing the scalar factor
\(e^{-\mathrm Tt}\), its common residual vector satisfies
\begin{equation}
 \dot\rho=-\mathsf G M_\sigma\rho,\qquad
 \mathsf G=4I-A_{C_4},\qquad
 M_\sigma=\operatorname{diag}(m_e(\sigma)).
\label{eq:phil-selection-tied-residual}
\end{equation}
Here the cyclic difference matrix contributes
\(2I-A_{C_4}\) and the shared-link drift contributes \(2I\).
Of the sixteen assignments, six have \(m_1=m_3\). Their first held motion
is \(e^{-ct}(0,-1,0,1)\), with \(c=4m_1\in\{0,4,8\}\).
At their first clock ring every receiver prefers the even-edge matching
\(S^{(0)}\). Holding that matching then leaves the odd residuals constant
and the even residuals zero in the scalar-free coordinates.

Four further assignments have multiplicity patterns
\[
 (1,1,2,0),\quad(1,0,2,1),\quad(2,0,1,1),\quad(2,1,1,0).
\]
The signed swap of coordinates \(1,3\) and the rotation by two preserve
the even-edge matching and carry the initial residual to itself or its
negative. They reduce these four cases to \((1,1,2,0)\). Write its
scalar-free coordinates as \(\rho=(-a,-b,-d,h)\). Then
\[
 \begin{pmatrix}\dot a\\\dot b\\\dot d\end{pmatrix}
 =\begin{pmatrix}-4&1&0\\1&-4&2\\0&1&-8\end{pmatrix}
   \begin{pmatrix}a\\b\\d\end{pmatrix},\qquad
 (a,b,d)(0)=(0,1,0),\qquad \dot h=-(a+2d),\ h(0)=1.
\]
The nonnegative cone is invariant, and the negative row sums give
\(\max(a,b,d)\le1\), hence \(a\le1/4\), \(d\le1/8\).
The half-spaces
\(a-3b/4-11d/25\le0\) and \(d-b/4\le0\) are invariant: their
outward derivatives on the corresponding boundary are
\(-b/400-7d/100\) and \(-a/4-b/8\), respectively. Consequently
\(a\le43b/50\), \(d\le b/4\), and \(b\ge e^{-4t}\).
The attended block is similar to a negative definite symmetric matrix
because \(\mathsf G\succ0\). Solving
\(\rho(\infty)=\rho(0)-\mathsf G z\), with \(z_3=0\) and the first
three final coordinates zero, gives \(h(\infty)=6/7\). Thus
\(h\ge6/7\).

If the first ring occurs at \(t_{\mathrm r}\) and redraws \(S^{(0)}\),
the subsequent scalar-free held motion, in elapsed time \(r\), is
\[
 \begin{aligned}
 a_r&=a_0e^{-8r},&d_r&=d_0e^{-8r},\\
 b_r&=b_0+\tfrac14(a_0+d_0)(1-e^{-8r}),&
 h_r&=h_0-\tfrac14(a_0+d_0)(1-e^{-8r})\ge\tfrac{171}{224}.
 \end{aligned}
\]
Every adjacent odd-edge absolute residual exceeds its even-edge neighbor
by at least \((7/50)e^{-4t_{\mathrm r}}\). Combining these four cases
with the preceding six, the deterministic motion \(\mathbf Y\) from
\(s\mathbf u_\circ\) has, throughout the hypothesized matching hold,
\begin{equation}
 \sqrt{\mathcal D_{\mathrm{odd}}(\mathbf Y_t)}
 -\sqrt{\mathcal D_{\mathrm{even}}(\mathbf Y_t)}
 \ge\frac7{50}s e^{-\mathrm Tt-8t_{\mathrm r}},
 \qquad t\ge t_{\mathrm r}.
\label{eq:phil-selection-lock-margin}
\end{equation}
There are therefore ten good initial assignments. For the remaining six
only \(f_{\mathbf v}\ge-2\) will be used.

Set \(h_s=(\log s)/\mathrm T\), for \(s>e\), and write the initial
configuration as \(s\mathbf u_\circ+\xi\), \(\xi\sim\mathcal N(0,I)\).
Along the actual assignment path decompose
\(\mathbf X=\mathbf Y+\mathbf E+\mathbf Z\), where the first two terms
solve its switched deterministic equation from \(s\mathbf u_\circ\) and
\(\xi\), respectively, while
\[
 d\mathbf Z=-\mathsf A_{\sigma_t}\mathbf Zdt+\sqrt{2\mathrm T}dW_t,
 \qquad\mathbf Z_0=0.
\]
All three are adapted differential equations, and
\(|\mathbf E_t|\le e^{-\mathrm Tt}|\xi|\).
It\^o's formula yields
\[
 \mathscr L e^{|\mathbf Z|^2/4}
 \le\mathrm T(6-|\mathbf Z|^2/4)e^{|\mathbf Z|^2/4}
 \le\mathrm Te^5.
\]
Stopping at radius \(\log s\), and using
\(\E e^{|\xi|^2/4}=64\), bounds the probability that either
\(|\xi|>\log s\) or \(\sup_{t\le h_s}|\mathbf Z_t|>\log s\) by
\((65+\mathrm Te^5h_s)e^{-(\log s)^2/4}\).
On its complement each square-root edge cost changes from its
\(\mathbf Y\) value by at most \(2\mathrm C_N\log s\), where
\(\mathrm C_N\) is the largest norm of a residual-pair map divided by
\(\sqrt2\).

Define
\[
 L_s(t_{\mathrm r})=h_s-\frac8{\mathrm T}t_{\mathrm r}
                         -\frac2{\mathrm T}\log\log s,\qquad
 B_s=(L_s(t_{\mathrm r})-t_{\mathrm r})_+.
\]
For a good initial assignment, \cref{eq:phil-selection-lock-margin}
gives a square-root gap at least \((7/50)(\log s)^2\) on
\([t_{\mathrm r},L_s]\). For sufficiently large \(s\), while both noise bounds hold,
the actual energy gap is at least
\(\kappa_0(\log s)^4\), \(\kappa_0=(7/50)^2/4\).
Each receiver's wrong choice at a ring then has probability at most
\(e^{-\kappa_0(\log s)^4/\mathrm T}\). Apply this bound at the first
ring, then use the rate-one compensator up to the first wrong redraw or
exit from the noise bound, stopping at \(L_s\wedge h_s\).
The endpoint \(L_s\) is known after the first ring, and the gap estimate
holds until that first failure. Thus the matching holds throughout the
nonempty interval \([t_{\mathrm r},L_s]\) except on an event of
probability at most
\begin{equation}
 \delta_s=(65+\mathrm Te^5h_s)e^{-(\log s)^2/4}
       +4(1+h_s)e^{-\kappa_0(\log s)^4/\mathrm T},
 \qquad h_s\delta_s\longrightarrow0.
\label{eq:phil-selection-lock-failure}
\end{equation}
Since \(\E t_{\mathrm r}=1\), a good initial assignment contributes at
least
\[
 -2h_s+4\E B_s-4h_s\delta_s
 \ge2h_s-4(1+8/\mathrm T)-(8/\mathrm T)\log\log s-4h_s\delta_s
\]
to \(\E\int_0^{h_s}f_{\mathbf v}(\sigma_t)dt\). The other six contribute
at least \(-2h_s\). Averaging the sixteen uniform initial assignments gives
\[
 \E\int_0^{h_s}f_{\mathbf v}(\sigma_t)dt
 \ge\tfrac12h_s-\frac5{\mathrm T}\log\log s-O_{\mathrm T}(1).
\]
By \cref{eq:phil-selection-moment},
\(\E|\mathbf X_{h_s}|^2\le14\). The Markov property and
\cref{eq:phil-selection-poisson} bound the integral after \(h_s\) in
absolute value by a constant independent of \(s\). Therefore
\begin{equation}
 \mathsf j_{\mathrm T}(s\mathbf u_\circ)
 \ge\frac{\log s}{2\mathrm T}
       -\frac5{\mathrm T}\log\log s-O_{\mathrm T}(1).
\label{eq:phil-selection-nonvanishing}
\end{equation}
This proves that \(\mathsf j_{\mathrm T}\) is not identically zero at any
fixed \(\mathrm T>0\). Its zero set is Lebesgue-null. For completeness,
this last assertion follows by induction on dimension: a nonzero analytic
function in one variable has isolated zeros; in several variables at least
one coefficient \(\partial_n^kf(\mathbf x',0)\) is nonzero as an analytic
function, or analyticity would make \(f\) vanish everywhere. Outside the
null zero set of that coefficient, each one-dimensional slice has isolated
zeros, and Fubini's theorem concludes the induction.

\subsection*{Entropy balance for the evolving recognition law}

Fix a finite horizon \(h\). All rows are bounded away from zero on this
interval by \cref{eq:phil-selection-row-bounds}. Write \(q_\sigma(t,\mathbf x)\)
for the evolving density, \(Q=\sum_\sigma q_\sigma\), and
\(b_\sigma(t,\mathbf x)=\mu_{p(t)}(\mathbf x,\sigma)/
 \sum_\eta\mu_{p(t)}(\mathbf x,\eta)\).
For the forward Ornstein--Uhlenbeck semigroup \(K_\sigma\), the last-proposal
formula is
\begin{equation}
 q_\sigma(t)=e^{-t}K_\sigma(t)q_\sigma(0)
   +\int_0^te^{-(t-r)}K_\sigma(t-r)[b_\sigma(r)Q(r)]dr.
\label{eq:phil-selection-forward}
\end{equation}
Iterating this positive formula, and bounding each proposal factor by one,
reduces each term to a Gaussian motion along a deterministic sequence of
sheets. Put \(M=\max_\sigma\|\mathsf A_\sigma\|\).
Its means have norm at most \(|\mathsf m_0|\), and its covariances lie between
\(e^{-2Mh}I\) and \((1+2\mathrm Th)I\). Summing the at most \(16^n\)
sequences on each time simplex, of volume \(t^n/n!\), gives a finite upper
Gaussian envelope. The no-proposal term gives a lower Gaussian envelope.
Consequently, uniformly on \([0,h]\),
\begin{equation}
 0<\mathrm c_he^{-\mathrm C_h(1+|\mathbf x|^2)}
 \le q_\sigma(t,\mathbf x)\le\mathrm C_he^{-\mathrm c_h|\mathbf x|^2}.
\label{eq:phil-selection-density-envelopes}
\end{equation}
This bound does not assert Gaussianity conditional on a state-dependent
assignment history.

Differentiating one Ornstein--Uhlenbeck kernel against a function bounded
by a polynomial times a Gaussian gives a bound of the same type with
factor \((t-r)^{-1/2}\). To check it, use the Gaussian kernel covariance
bounds proportional to \(t-r\), absorb its linear derivative into a
Gaussian with twice the covariance, and complete the square in the
convolution. This singularity is integrable in
\cref{eq:phil-selection-forward}, proving a first-derivative Gaussian
bound. Moreover
\[
 \nabla b_\sigma=-\frac{b_\sigma}{\mathrm T}
       \left(H_\sigma\mathbf x-\sum_\eta b_\eta H_\eta\mathbf x\right),
 \qquad |\nabla b_\sigma|\le\mathrm K_{\mathrm T}|\mathbf x|.
\]
The source \(b_\sigma Q\) therefore has a first derivative with the same
type of bound. For the second derivative, transfer one derivative to this
source using
\(\nabla K_\sigma(r)g=e^{r\mathsf A_\sigma}K_\sigma(r)\nabla g\),
and differentiate the other kernel factor. The singularity is again only
\(r^{-1/2}\). The explicit initial Gaussian term supplies the bounds at
zero. The forward equation obtained by taking a time difference quotient
in \cref{eq:phil-selection-forward} then bounds the time derivative too:
\begin{equation}
 |\nabla q_\sigma|+|\nabla^2q_\sigma|+|\partial_tq_\sigma|
 \le\mathrm C_h(1+|\mathbf x|)^k e^{-\mathrm c_h|\mathbf x|^2}.
\label{eq:phil-selection-density-derivatives}
\end{equation}
The derivative transfer follows directly by integration by parts in the
Gaussian transition density; its inverse mean map is uniformly bounded
on this finite interval. These estimates justify the differentiations,
without assuming finite Fisher information.

Set \(f_\sigma=q_\sigma/\mu_{p(t),\sigma}\) and
\(\mathsf H(t)=\KL(q_t\Vert\mu_{p(t)})\).
The logarithms of both densities grow at most quadratically by
\cref{eq:phil-selection-density-envelopes}, and
\(\partial_t\log\mu_p\) is bounded in \(\mathbf x\), since only its prior
weight and normalizer vary. Together with
\cref{eq:phil-selection-density-derivatives}, this proves that
\(\mathsf H\) is absolutely continuous on \([0,h]\) and permits its
integral to be differentiated.
To integrate the diffusion by parts, first multiply by a smooth radial
cutoff \(\chi_R\), equal to one on \(B_R\), zero outside \(B_{2R}\),
with \(|\nabla\chi_R|\le\mathrm C/R\), increasing to one.
On each compact cylinder the densities have positive lower bounds, so
\(\log f_\sigma\) is a legitimate test function. The boundary term is
bounded by
\[
 \frac{\mathrm C_h}{R}\int_{R\le|\mathbf x|\le2R}
       (1+|\mathbf x|)^{k+3}e^{-\mathrm c_h|\mathbf x|^2}d\mathbf x
 \longrightarrow0.
\]
Detailed balance \(\mu_\sigma b_\eta=\mu_\eta b_\sigma\) symmetrizes the
finite jump sum. Both resulting production integrands are nonnegative,
so monotone convergence removes the cutoff and simultaneously proves
finiteness and local time integrability of the production. The result is
\begin{equation}
 \begin{aligned}
 \mathsf H'(t)&=-\mathscr D_{p(t)}(f)
               -\sum_\sigma\int q_\sigma\partial_t\log\mu_{p(t),\sigma},\\
 \mathscr D_p(f)&=\mathrm T\sum_\sigma\int\mu_{p,\sigma}
                         \frac{|\nabla f_\sigma|^2}{f_\sigma}\\
 &\quad+\frac12\sum_{\sigma,\eta}\int\mu_{p,\sigma}b_\eta
                  (f_\eta-f_\sigma)(\log f_\eta-\log f_\sigma).
 \end{aligned}
\label{eq:phil-selection-entropy-identity}
\end{equation}
The equality holds almost everywhere and in integrated form on every finite
interval. The jump term is integrable already from the Gaussian envelopes
and quadratic logarithmic bounds; thus the cutoff argument does not
presuppose a bound on the diffusion production. This also proves
\(\mathsf H(0)<\infty\).

\subsection*{Entropy contraction at frozen rows}

In the original twelve-dimensional coordinates put
\(\mathsf B_\sigma=\mathsf A_\sigma/\mathrm T\), the Gaussian sheet
precision, and set
\[
 \kappa=\min_\sigma\lambda_{\min}(\mathsf B_\sigma)>0,\qquad
 M=\max_\sigma\|\mathsf B_\sigma\|,\qquad
 q_\sigma^*=\mu_p(\sigma)
 =\frac{\pi_\sigma\mathsf w_\sigma}{\mathsf Z},
 \qquad q_{\min}^*=\min_\sigma q_\sigma^*.
\]
Let \(\gamma_\sigma=\mathcal N(0,\mathsf B_\sigma^{-1})\), and, for a
real function \(u\), define
\[
 \mathscr E_c(u)=\mathrm T\sum_\sigma\int\mu_{p,\sigma}|\nabla u_\sigma|^2,
 \qquad
 \mathscr E_j(u)=\frac12\sum_{\sigma,\eta}\int\mu_{p,\sigma}b_\eta
                                    (u_\sigma-u_\eta)^2.
\]
The Gaussian inequalities needed below follow from the explicit
Ornstein--Uhlenbeck semigroup with invariant law \(\gamma_\sigma\).
Its gradient commutes as
\(\nabla P_tg=e^{-t\mathsf B_\sigma}P_t\nabla g\). Differentiating
entropy and applying Cauchy--Schwarz inside this kernel gives
\[
 -\frac d{dt}\operatorname{Ent}_{\gamma_\sigma}(P_tg)
 =\int\frac{|\nabla P_tg|^2}{P_tg}d\gamma_\sigma
 \le e^{-2\kappa t}\int\frac{|\nabla g|^2}{g}d\gamma_\sigma.
\]
For bounded positive smooth \(g\), the explicit Gaussian formula and
dominated convergence make the entropy tend to zero. Integration, with
\(g=u^2\), gives
\begin{equation}
 \operatorname{Ent}_{\gamma_\sigma}(u^2)
 \le\frac2\kappa\int|\nabla u|^2d\gamma_\sigma,
 \qquad
 \operatorname{Var}_{\gamma_\sigma}(u)
 \le\frac1\kappa\int|\nabla u|^2d\gamma_\sigma.
\label{eq:phil-selection-gaussian-inequalities}
\end{equation}
The second estimate follows by differentiating variance in the same
semigroup. Positive regularization, truncation, and smooth cutoffs extend
both estimates to their form domains: the cutoff errors tend to zero by
integrability of the squared function and its gradient, and the entropy
passes to the limit by its variational formula.

Replace \(u\) by \(|u|\), which decreases both energies, and normalize
\(\int u^2d\mu_p=1\). Put
\(z_\sigma=(\int u_\sigma^2d\gamma_\sigma)^{1/2}\).
The entropy chain rule gives
\[
 \operatorname{Ent}_{\mu_p}(u^2)
 =\sum_\sigma q_\sigma^*
             \operatorname{Ent}_{\gamma_\sigma}(u_\sigma^2)
       +\operatorname{Ent}_{q^*}(z^2).
\]
Here \(0\le z_\sigma\le R=(q_{\min}^*)^{-1/2}\).
The function \(\psi(z)=z^2\log z^2-z^2+1\) satisfies
\(\psi(1)=\psi'(1)=0\) and \(\psi''(z)=4+4\log z\), whence
\(\psi(z)\le2(1+\log R)(z-1)^2\) on \([0,R]\).
Since \(\sum q_\sigma^*z_\sigma^2=1\),
\[
 \operatorname{Ent}_{q^*}(z^2)
 \le L_p\operatorname{Var}_{q^*}(z),\qquad
 L_p=4+2\log(1/q_{\min}^*).
\]
Indeed \(\sum q_\sigma^*(z_\sigma-1)^2
=2(1-\E_{q^*}z)\le2\operatorname{Var}_{q^*}(z)\).

Fix a ball \(B=B_\rho\) in these coordinates, and define
\[
 m_\sigma=|B|^{-1}\int_Bu_\sigma,
 \quad m_* =\sum_\sigma\pi_\sigma m_\sigma,
 \quad X_\sigma=
 \left[\frac{(2\pi)^6\mathsf w_\sigma
                   e^{\rho^2\|\mathsf B_\sigma\|/2}}{|B|}-1\right]_+,
 \quad B_{\mathrm T}=\max_\sigma(1+\sqrt{X_\sigma})^2.
\]
Variance about any center is an upper bound on the variance, so
\[
 \operatorname{Var}_{q^*}(z)
 \le2\sum_\sigma q_\sigma^*(z_\sigma-m_\sigma)^2
     +2\sum_\sigma q_\sigma^*(m_\sigma-m_*)^2.
\]
The first sum is at most \((B_{\mathrm T}/(\mathrm T\kappa))\mathscr E_c(u)\):
\(|z_\sigma-\gamma_\sigma u_\sigma|\le
\sqrt{\operatorname{Var}_{\gamma_\sigma}(u_\sigma)}\), and
Cauchy--Schwarz bounds \(|m_\sigma-\gamma_\sigma u_\sigma|\) by the same
standard deviation times
\(\sqrt{\chi^2(\operatorname{Unif}(B)\Vert\gamma_\sigma)}
\le\sqrt{X_\sigma}\).
For the second sum, Jensen's inequality gives
\[
 \sum_\sigma q_\sigma^*(m_\sigma-m_*)^2
 \le\frac{\mathsf w_{\max}}
              {\mathsf Z|B|}
       \sum_{\sigma,\eta}\pi_\sigma\pi_\eta
                          \int_B(u_\sigma-u_\eta)^2.
\]
On \(B\), the explicit Gaussian densities give
\[
 \mu_{p,\sigma}b_\eta
 \ge\frac{\pi_\sigma\pi_\eta e^{-\rho^2M}}
              {(2\pi)^6\mathsf Z}.
\]
The normalizer therefore cancels, and the preceding sum is at most
\(2(2\pi)^6\mathsf w_{\max}e^{\rho^2M}
 \mathscr E_j(u)/|B|\).
Combining these bounds proves
\[
 \operatorname{Ent}_{\mu_p}(u^2)
 \le\frac2{\mathrm T\kappa}(1+L_pB_{\mathrm T})\mathscr E_c(u)
    +\frac{4L_p(2\pi)^6\mathsf w_{\max}e^{\rho^2M}}{|B|}
                                                       \mathscr E_j(u).
\]
For \(u=\sqrt f\), the production is at least
\(4(\mathscr E_c(u)+\mathscr E_j(u))\). Its jump part uses
\((a-b)(\log a-\log b)\ge4(\sqrt a-\sqrt b)^2\), which follows by
Cauchy--Schwarz on the integral of \(r^{-1/2}\) between \(a\) and \(b\).
The inequality extends to the form domain by the preceding approximations.
Finally,
\[
 q_{\min}^*\ge p_{\min}^4
                 \frac{\mathsf w_{\min}}
                      {\mathsf w_{\max}},\qquad
 p_{\min}=\min_i\{p_i,1-p_i\}.
\]
Thus there are finite positive constants depending only on \(\mathrm T\)
with
\begin{equation}
 \mathscr D_p(f)\ge\lambda(\mathrm T,p)
                         \operatorname{Ent}_{\mu_p}(f),\qquad
 \lambda(\mathrm T,p)\ge
 \frac{\mathrm c_{\mathrm T}}
      {\mathrm K_{\mathrm T}+8\log(1/p_{\min})}.
\label{eq:phil-selection-entropy-contraction}
\end{equation}
The factor \(\mathrm T\) in \(\mathscr E_c\) is the same diffusion
coefficient as in \cref{eq:phil-selection-entropy-identity}.

\subsection*{The matching branch of the small-temperature landscape}

The needed sheet weights can be computed using two four-dimensional
determinants. Let \(D\) be the cyclic difference matrix, and
\(C_\sigma=\operatorname{diag}(m_e(\sigma))\). In the orthogonal
coordinates \((\bar k+\bar m)/\sqrt2\),
\((\bar k-\bar m)/\sqrt2\), and \(\omega\), the energy splits into
\(\frac12(Du-\sqrt2\omega)^\top C_\sigma(Du-\sqrt2\omega)
 +\frac12(Dv)^\top C_\sigma Dv\), while the discrepancy coordinate has
reference precision \(5I\). The determinant identity for the first block
then gives
\begin{equation}
 \det\mathsf A_\sigma
 =\mathrm T^4
  \det[\mathrm TI+C_\sigma(DD^\top+2I)]
  \det[5\mathrm TI+D^\top C_\sigma D].
\label{eq:phil-selection-sheet-determinants}
\end{equation}
The two matching assignments, eight end-path assignments, four middle-path
assignments, and two oriented cycles have the respective determinants
\[
 \begin{aligned}
 \mathsf d_M&=25\mathrm T^8(\mathrm T+8)^2(5\mathrm T+4)^2,\\
 \mathsf d_E&=5\mathrm T^6
   (\mathrm T^3+16\mathrm T^2+77\mathrm T+112)
   (125\mathrm T^3+200\mathrm T^2+85\mathrm T+8),\\
 \mathsf d_N&=5\mathrm T^6(\mathrm T+4)(5\mathrm T+2)
   (\mathrm T^2+12\mathrm T+28)(25\mathrm T^2+30\mathrm T+4),\\
 \mathsf d_C&=5\mathrm T^5(\mathrm T+2)(\mathrm T+4)^2(\mathrm T+6)
                                     (5\mathrm T+2)^2(5\mathrm T+4).
 \end{aligned}
\]
These follow in \cref{eq:phil-selection-sheet-determinants} from the
multiplicity representatives \((2,0,2,0)\), \((2,1,1,0)\),
\((1,2,1,0)\), and \((1,1,1,1)\), respectively.
Normalize the matching weight to one and write
\[
 \mathsf r_E=\sqrt{\mathsf d_M/\mathsf d_E},\qquad
 \mathsf r_N=\sqrt{\mathsf d_M/\mathsf d_N},\qquad
 \mathsf r_C=\sqrt{\mathsf d_M/\mathsf d_C}.
\]
Their leading terms are
\(\mathsf r_E,\mathsf r_N=\sqrt{40/7}\mathrm T+O(\mathrm T^2)\) and
\(\mathsf r_C=\sqrt{5/3}\mathrm T^{3/2}(1+O(\mathrm T))\).
Choose \(\mathrm T_0>0\) so that \(0<\mathsf r_E,\mathsf r_N,\mathsf r_C\le1/8\) for
\(0<\mathrm T\le\mathrm T_0\).

Use matching coordinates
\(x=(p_0,1-p_1,p_2,1-p_3)\). The normalized evidence is the
multilinear polynomial whose vertex weights are one at \(\mathbf0,\mathbf1\),
\(\mathsf r_E\) for subsets of cardinality one or three, \(\mathsf r_N\) for the four adjacent
pairs, and \(\mathsf r_C\) for the two opposite pairs. Complementation preserves
these weights, so the centered reduced field is odd. With
\(D_0=1+4\mathsf r_E+2\mathsf r_N+\mathsf r_C\), its Jacobian at \(x=\mathbf1/2\) is
\[
 \frac1{D_0}
 \begin{pmatrix}
 0&1-\mathsf r_C&1-2\mathsf r_N+\mathsf r_C&1-\mathsf r_C\\
 1-\mathsf r_C&0&1-\mathsf r_C&1-2\mathsf r_N+\mathsf r_C\\
 1-2\mathsf r_N+\mathsf r_C&1-\mathsf r_C&0&1-\mathsf r_C\\
 1-\mathsf r_C&1-2\mathsf r_N+\mathsf r_C&1-\mathsf r_C&0
 \end{pmatrix}.
\]
For example, the adjacent Hessian entry of the evidence is
\((1-\mathsf r_C)/2\), its uniform value is \(D_0/8\), and the row mobility is
\(1/4\), giving the displayed entry. The eigenvalues are
\begin{equation}
 \alpha=\frac{3-2\mathsf r_N-\mathsf r_C}{D_0}>0,\qquad
 \beta=\frac{-1-2\mathsf r_N+3\mathsf r_C}{D_0}<0,\qquad
 \eta=\frac{-1+2\mathsf r_N-\mathsf r_C}{D_0}<0\quad\text{twice}.
\label{eq:phil-selection-saddle-spectrum}
\end{equation}
The positive eigendirection is \(\mathbf1\) in matching coordinates,
or \(\mathbf v\) in the original rows. In fact the chosen bounds give
\(\alpha\ge7/5\), \(\beta\le-1/3\), and \(\eta\le-2/5\).

Cyclic invariance makes the matching diagonal \(x=z\mathbf1\) invariant.
Writing \(w=z(1-z)\), its normalized evidence is
\[
 Z_{\mathrm{diag}}(z)
 =1+4(\mathsf r_E-1)w+(2-8\mathsf r_E+4\mathsf r_N+2\mathsf r_C)w^2,
 \qquad
 \dot z=\frac{z(1-z)}{4Z_{\mathrm{diag}}(z)}Z_{\mathrm{diag}}'(z).
\]
The derivative is \((1-2z)[4(\mathsf r_E-1)+2(2-8\mathsf r_E+4\mathsf r_N+2\mathsf r_C)w]\).
The bracket is negative at both endpoints \(w=0,1/4\), where its values
are \(4(\mathsf r_E-1)\) and \(-3+2\mathsf r_N+\mathsf r_C\). It is affine in \(w\), hence negative
throughout. Thus the diagonal motion tends to \(\mathbf1\) from
\(z>1/2\), and to \(\mathbf0\) from \(z<1/2\), reaching every prescribed
endpoint neighborhood in finite slow time. At either vertex the Jacobian
is \((\mathsf r_E-1)I\): a single receiver flip changes weight one to \(\mathsf r_E\), the
mobility vanishes, and its diagonal derivative supplies the factor
\(-(1-\mathsf r_E)\). Both vertices are therefore sinks with positive inward rate
\(\nu=1-\mathsf r_E\).

\subsection*{Tracking, escape, and selection}

Let \(a_i(t)=q_t(\sigma_i=+1)\) and
\(\Delta_i(t)=a_i(t)-\bar\beta_i(p(t))\).
Differentiating the normalized product prior gives
\[
 \partial_{p_i}\log\mu_p
 =\frac{1_{\{\sigma_i=+1\}}-\bar\beta_i(p)}{p_i(1-p_i)}.
\]
Inserting the learning equation and
\cref{eq:phil-selection-reduced-field} in
\cref{eq:phil-selection-entropy-identity} yields
\begin{equation}
 \mathsf H'
 =-\mathscr D_p(f)-\varepsilon\sum_i\mathfrak g_i(p)\Delta_i
                  -\varepsilon\sum_i\frac{\Delta_i^2}{p_i(1-p_i)}
 \le-\lambda(\mathrm T,p)\mathsf H
      +\varepsilon\|\mathfrak g(p)\|_1\sqrt{\mathsf H/2}.
\label{eq:phil-selection-learning-entropy}
\end{equation}
The last bound uses \(|\Delta_i|\le\operatorname{TV}(q_t,\mu_p)
\le\sqrt{\mathsf H/2}\). The nonpositive square is kept until this step;
it removes any inverse-row factor from the forcing.
Writing \(\Lambda(t)=\int_0^t\lambda(\mathrm T,p(r))dr\), comparison
for \(\sqrt{\mathsf H+\delta^2}\), followed by \(\delta\downarrow0\), gives
\begin{equation}
 \sqrt{\mathsf H(t)}
 \le e^{-\Lambda(t)/2}\sqrt{\mathsf H(0)}
    +\frac{\varepsilon}{2\sqrt2}
       \int_0^te^{-(\Lambda(t)-\Lambda(r))/2}
                         \|\mathfrak g(p(r))\|_1dr.
\label{eq:phil-selection-tracking}
\end{equation}
The bounded fitness, \cref{eq:phil-selection-row-bounds}, and
\cref{eq:phil-selection-entropy-contraction} show that, on
\(t\le2\mathrm C_{\mathrm T}\varepsilon^{-1}\log(1/\varepsilon)\),
\begin{equation}
 \operatorname{TV}(q_t,\mu_{p(t)})
 \le e^{-\Lambda(t)/2}\sqrt{\mathsf H(0)/2}
       +O_{\mathrm T}(\varepsilon\log(1/\varepsilon)).
\label{eq:phil-selection-global-tracking}
\end{equation}
Indeed \(\lambda\) is bounded below on this horizon by a positive
constant divided by \(1+\log(1/\varepsilon)\).

For the saddle a finer bound is needed. Let
\(\|\mathfrak g(p)-\mathfrak g(p')\|_1
 \le L_{\mathrm T}\|p-p'\|_\infty\), and suppose
\(\lambda\ge\lambda_\flat>0\) on \([s_0,t]\).
The estimate \(|\dot p_i|\le\varepsilon\) in
\cref{eq:phil-selection-tracking} gives
\begin{equation}
 \operatorname{TV}(q_t,\mu_{p(t)})
 \le e^{-\lambda_\flat(t-s_0)/2}\sqrt{\mathsf H(s_0)/2}
   +\frac{\varepsilon}{2\lambda_\flat}\|\mathfrak g(p(t))\|_1
   +\frac{L_{\mathrm T}\varepsilon^2}{\lambda_\flat^2}.
\label{eq:phil-selection-refined-tracking}
\end{equation}
The constants follow by integrating \(1\) and \(t-r\) against
\(e^{-\lambda_\flat(t-r)/2}\).

Fix \(0<\mathrm T\le\mathrm T_0\) and a mean with
\(\mathsf j_{\mathrm T}(\mathsf m_0)\ne0\), and put
\(\mathbf e=\mathbf v/2\). In slow time
\(\mathfrak s=\varepsilon t\), write
\(y=p-1/2=a\mathbf e+z\), \(z\perp\mathbf e\).
By \cref{eq:phil-selection-imprint-expansion}, at
\(t_1=\varepsilon^{-1/4}\),
\[
 a(t_1)=\varepsilon(\mathsf j_{\mathrm T}(\mathsf m_0)/2+o(1)),
 \qquad |z(t_1)|=O(\varepsilon).
\]
On \([0,t_1]\), symmetry gives \(\mathfrak g(1/2)=0\), so
\(\|\mathfrak g(p(t))\|_1\le L_{\mathrm T}\varepsilon t/2\).
The bound in \cref{eq:phil-selection-tracking} gives
\(\sqrt{\mathsf H(t_1)}=O(\varepsilon^{7/4})\).
Consequently, until exit from any sufficiently small fixed ball,
\cref{eq:phil-selection-refined-tracking} implies
\begin{equation}
 |\Delta|_2\le\kappa_\varepsilon|y|+b_\varepsilon,
 \qquad \kappa_\varepsilon=O(\varepsilon),\quad
 b_\varepsilon=O(\varepsilon^{7/4})=o(|a(t_1)|).
\label{eq:phil-selection-relative-error}
\end{equation}
The constants here may depend on the fixed mean.

Let \(\beta_\perp=\max(\beta,\eta)<0\) and
\(\gamma=\alpha-\beta_\perp>0\).
Oddness and diagonal invariance of the analytic reduced field imply,
for \(|y|\le r\),
\[
 |\langle\mathscr V(1/2+y),\mathbf e\rangle-\alpha a|
       \le\mathrm C|y|^3,\qquad
 |P_\perp\mathscr V(1/2+y)-D\mathscr V(1/2)z|
       \le\mathrm C|y|^2|z|.
\]
For the second inequality, subtract the zero perpendicular field on
\(a\mathbf e\) and integrate its derivative along the transverse
segment. The derivative differs from its value at zero by
\(O(|y|^2)\) because the field is odd.
Writing \(w=|z|/|a|\), the upper right derivatives of norms give
\begin{equation}
 \begin{aligned}
 \frac{d|a|}{d\mathfrak s}
 &\ge |a|[\alpha-(\mathrm Cr^2+\kappa_\varepsilon)(1+w)]
                -b_\varepsilon,\\
 D^+w
 &\le-[\gamma-\mathrm Cr^2
             -(\mathrm Cr^2+\kappa_\varepsilon)(1+w)]w
        +(1+w)(\kappa_\varepsilon+b_\varepsilon/|a|).
 \end{aligned}
\label{eq:phil-selection-cone}
\end{equation}
Choose \(W\) larger than twice the initial limiting ratio, then choose
\(r>0\) so small that the terms \(\mathrm Cr^2(1+W)\) consume less than
one quarter of both positive rates. For small enough \(\varepsilon\),
\cref{eq:phil-selection-relative-error} controls the remaining terms.
A bootstrap then keeps \(w\le W\), makes \(|a|\) increase without changing
sign, and gives
\[
 w(\mathfrak s)
 \le w(\mathfrak s_1)e^{-\gamma(\mathfrak s-\mathfrak s_1)/2}
  +\frac{2(1+W)}\gamma
       (\kappa_\varepsilon+b_\varepsilon/|a(t_1)|).
\]
The lower growth bound for \(|a|\) yields an exit time no larger than
\[
 \frac{1+O(r^2)+o(1)}\alpha\log(1/\varepsilon)
\]
in slow time. The corresponding upper norm-growth bound forces the exit
time to tend to infinity. Hence \(w\to0\) at exit, and the exit point
tends to \(\operatorname{sign}(\mathsf j_{\mathrm T})r\mathbf e\).

The reduced diagonal orbit from that point enters any fixed neighborhood
of its matching vertex in finite slow time. The tracking error in
\cref{eq:phil-selection-global-tracking} tends to zero after \(t_1\);
finite-time Gr\"onwall comparison therefore carries the evolving rows into
the same neighborhood. At its selected vertex let
\(\zeta_i=|p_i-p_i^M|\). By the derivative computation above and continuity,
a sufficiently small neighborhood has
\((1-\zeta_i)|\mathfrak g_i(p)|\ge\nu/2\), where \(\nu=1-\mathsf r_E>0\), and
the fitness has the sign directed toward that vertex. Thus
\[
 D^+\zeta_i\le-\tfrac\nu2\zeta_i+
                   \operatorname{TV}(q_t,\mu_{p(t)})
\]
in slow time. The neighborhood remains invariant once the tracking floor
is small enough, and \(\zeta_i\) decays exponentially to a vanishing
floor.

Choose any \(\mathrm C_{\mathrm T}>1/\alpha\), independently of the mean,
and then choose the local radius so that its escape coefficient is less
than \(\mathrm C_{\mathrm T}\). The time remaining after the finite
transit tends to infinity uniformly on the slow-time window
\([\mathrm C_{\mathrm T}\log(1/\varepsilon),
  2\mathrm C_{\mathrm T}\log(1/\varepsilon)]\).
The row distances therefore vanish uniformly on that window. Every
normalized sheet weight is at most the matching weight one, so
\[
 \mu_p(M)\ge\pi_M\ge1-\sum_i\zeta_i,
 \qquad
 q_t(S\ne M)\le\sum_i\zeta_i+
                     \operatorname{TV}(q_t,\mu_{p(t)}).
\]
The last total variation tends to zero by
\cref{eq:phil-selection-global-tracking}; its integrated contraction
exponent at the beginning of the window grows at least as a positive
constant times \(\varepsilon^{-1}\).
This proves the first conclusion for every mean outside the null set in
\cref{eq:phil-selection-nonvanishing}. For the frozen arm,
\cref{eq:phil-selection-harris} and rotation symmetry prove the second
conclusion at every \(\mathrm T>0\) and every mean. The convergence is
pointwise in the fixed mean; no uniform bound near the zero set of the
imprint has been used. \(\square\)

\section{Proof of Joint Selection Along a Schedule}
\label{sec:phil-joint-selection-proof}

\noindent\emph{Proof of \Cref{cor:phil-joint-selection}.}
The saddle estimate \cref{eq:phil-selection-saddle-spectrum} gives
\(\alpha(\mathrm T)\ge7/5\) throughout the stated temperature range.
The proof of \Cref{conj:phil-contrast-formation} permits any window
constant greater than \(1/\alpha(\mathrm T)\), independently of the
initial mean. Thus the common choice \(2\) is valid. Write
\[
 \tau_\varepsilon=2\varepsilon^{-1}\log(1/\varepsilon),
 \qquad W_\varepsilon=[\tau_\varepsilon,2\tau_\varepsilon],
 \qquad 0<\varepsilon<e^{-1}.
\]
For a perfect matching \(M\), define
\[
 \nu_{\mathrm T}^M
 =\mathcal N(0,\mathrm T\mathsf A_M^{-1})\otimes\delta_M,
 \qquad
 \mu_{\mathrm T}=\mu_{\mathrm T,1/2}.
\]
Here \(\mu_{\mathrm T,p}\) is the frozen-row law used in the preceding
proof, with its temperature now displayed explicitly. The inverse
exists because \(\mathsf A_M\succeq\mathrm TI\). Let
\(M_{\mathrm T}(\mathsf m_0)\) be the matching selected by the sign of
\(\mathsf j_{\mathrm T}(\mathsf m_0)\), choosing either matching on its
null zero set.

First record the full one-time consequence of the fixed-temperature
proof. Conditional on assignment \(M\), the law
\(\mu_{\mathrm T,p}\) has the Gaussian factor of \(\nu_{\mathrm T}^M\),
independently of \(p\). The normalized sheet weights are at most the
matching weight one. For
\(\zeta_i=|p_i-p_i^M|\), therefore,
\[
 \mu_{\mathrm T,p}(M)\ge\pi_M\ge1-\sum_i\zeta_i,
 \qquad
 \operatorname{TV}(\mu_{\mathrm T,p},\nu_{\mathrm T}^M)
 =1-\mu_{\mathrm T,p}(M)\le\sum_i\zeta_i.
\]
The row convergence and
\cref{eq:phil-selection-global-tracking} imply
\[
 \sup_{t\in W_\varepsilon}
 \operatorname{TV}
 (q_t^{\mathrm T,\varepsilon,\mathsf m_0},
       \nu_{\mathrm T}^{M_{\mathrm T}(\mathsf m_0)})
 \longrightarrow0
\]
at fixed temperature outside the imprint's null zero set.
Moreover \cref{eq:phil-selection-tracking} itself gives convergence
of the learned relative entropy to zero, uniformly on this window.
For frozen rows, the entropy identity and contraction give
\[
 \KL(q_t^{\mathrm T,0,\mathsf m_0}\Vert\mu_{\mathrm T})
 \le e^{-\lambda(\mathrm T,1/2)t}
       \KL(q_0^{\mathrm T,0,\mathsf m_0}\Vert\mu_{\mathrm T}).
\]
The initial entropy is finite by the Gaussian density bounds, and
\(\lambda(\mathrm T,1/2)>0\). Frozen entropy and total variation
therefore tend to zero on the window for every fixed mean.

Use the explicit entropy majorants as follows. Let
\[
 G_{\mathrm T}=\sup_{p\in[0,1]^4}\|\mathfrak g(p)\|_1,\qquad
 \ell_{\mathrm T}(\varepsilon)
 =\frac{\mathrm c_{\mathrm T}}
 {\mathrm K_{\mathrm T}+8\log2+32\log(1/\varepsilon)},\qquad
 \lambda_u(\mathrm T)
 =\frac{\mathrm c_{\mathrm T}}{\mathrm K_{\mathrm T}+8\log2}.
\]
The first constant is finite. The row lower bound and
\cref{eq:phil-selection-entropy-contraction} make
\(\ell_{\mathrm T}(\varepsilon)\) a rate lower bound throughout
\(t\le2\tau_\varepsilon\), and \(\lambda_u(\mathrm T)\) a rate lower
bound at uniform rows. The common initial entropy equals
\[
 H_0(\mathrm T,\mathsf m_0)
 =\log\mathsf Z(1/2)-6+
 \frac{\operatorname{tr}\overline{\mathsf A}
              +\mathsf m_0^\top\overline{\mathsf A}\mathsf m_0}
      {2\mathrm T},
 \qquad
 \overline{\mathsf A}=\frac1{16}\sum_\sigma\mathsf A_\sigma.
\]
Indeed the log density ratio is
\(\log\mathsf Z(1/2)-|\mathbf x-\mathsf m_0|^2/2+
\mathbf x^\top\mathsf A_\sigma\mathbf x/(2\mathrm T)\);
taking its initial expectation proves the formula.
Thus \(H_0\) is continuous and nonnegative. Uniformly on
\(W_\varepsilon\), the entropy tracking estimate gives
\[
 \begin{aligned}
 \KL(q_t^{\mathrm T,\varepsilon,\mathsf m_0}
               \Vert\mu_{\mathrm T,p(t)})
 &\le B_L(\mathrm T,\varepsilon,\mathsf m_0)\\
 &:=\left[
 e^{-\ell_{\mathrm T}(\varepsilon)\tau_\varepsilon/2}\sqrt{H_0}
 +\frac{\varepsilon G_{\mathrm T}}
             {\sqrt2\ell_{\mathrm T}(\varepsilon)}
 \right]^2,\\
 \KL(q_t^{\mathrm T,0,\mathsf m_0}\Vert\mu_{\mathrm T})
 &\le B_F(\mathrm T,\varepsilon,\mathsf m_0)
 :=e^{-\lambda_u(\mathrm T)\tau_\varepsilon}H_0.
 \end{aligned}
\]
These majorants are continuous in positive mobility and mean and tend
to zero at fixed temperature and mean as \(\varepsilon\downarrow0\).

Let \(\mathcal E_{\mathrm T}(\varepsilon,\mathsf m_0)\) be the minimum
of one and the maximum of the following five nonnegative quantities:
\[
 \begin{aligned}
 &\sup_{t\in W_\varepsilon}
          \max_i|p_i(t)-p_i^{M_{\mathrm T}(\mathsf m_0)}|,\\
 &\sup_{t\in W_\varepsilon}
       \operatorname{TV}
       (q_t^{\mathrm T,\varepsilon,\mathsf m_0},
                   \nu_{\mathrm T}^{M_{\mathrm T}(\mathsf m_0)}),\\
 &B_L(\mathrm T,\varepsilon,\mathsf m_0),\\
 &\sup_{t\in W_\varepsilon}
       \operatorname{TV}
          (q_t^{\mathrm T,0,\mathsf m_0},\mu_{\mathrm T}),\\
 &B_F(\mathrm T,\varepsilon,\mathsf m_0).
 \end{aligned}
\]
The preceding estimates prove
\(\mathcal E_{\mathrm T}(\varepsilon,\mathsf m_0)\to0\)
for almost every fixed mean at each fixed temperature.

We will take a supremum over mobility, so measurability must be checked.
Fix a temperature and restrict the positive mobility and the mean to
compact sets, with the mobility bounded away from zero. The corresponding
horizons are bounded, and
\cref{eq:phil-selection-row-bounds} bounds every row away from zero and
one. On this row set the posterior redraw probabilities are uniformly
Lipschitz in the rows, independently of the configuration.
Maximally couple the initial laws and then use common Brownian motions
and proposal clocks until the first different redraw. The law difference
is bounded by the initial total variation plus a constant times the
integral of row differences. Subtracting the learning equations bounds
the row difference by integrals of row and law differences and by the
mobility difference. Gr\"onwall's inequality proves total-variation
continuity in the initial mean and positive mobility on each bounded
time interval. Time continuity follows from
\cref{eq:phil-selection-forward}.

Parametrizing the window by \(t=u\tau_\varepsilon\),
\(1\le u\le2\), proves continuity in positive mobility of each row or
total-variation window supremum. The explicit entropy majorants are
continuous by their formulas. The matching as a function of the mean is
Borel because its imprint is analytic. Thus the error is measurable
in the mean, and its supremum over an interval of positive mobilities
equals its supremum over rational mobilities. No uniform continuity at
zero mobility or separate entropy-continuity assertion is needed.

Fix the sequence \(\mathrm T_n\downarrow0\) in the stated temperature
interval and \(0<\eta_n<1\) with \(\eta_n\to0\). For the prescribed
exponent \(a>7/2\), put \(h_n=\mathrm T_n^a\). The following argument
in fact permits arbitrary positive caps \(h_n\).
On the mean space \(\R^{12}\), introduce the auxiliary probability
measure \(\rho=\mathcal N(0,I_{12})\). Its density is strictly positive;
it introduces no assumption that the physical mean is random.
Set
\[
 F_n(r,\mathsf m_0)
 =\sup_{0<\varepsilon\le r}
       \mathcal E_{\mathrm T_n}(\varepsilon,\mathsf m_0).
\]
The measurable functions \(F_n\) are bounded by one. At fixed \(n\),
they tend to zero for \(\rho\)-almost every mean as \(r\downarrow0\).
Dominated convergence yields
\[
 \int F_n(r,\mathsf m_0)d\rho(\mathsf m_0)\longrightarrow0.
\]
Choose \(d_n=2^{-k_n}\) recursively, where \(k_n\) is the smallest
positive integer satisfying
\[
 d_n\le\min(e^{-1},2^{-n},h_n,d_{n-1}/2),
 \qquad
 \int F_n(d_n,\mathsf m_0)d\rho(\mathsf m_0)
       \le\eta_n2^{-n},
\]
with the \(d_{n-1}/2\) condition omitted for \(n=1\). The preceding
limit proves existence. This defines deterministic thresholds
independent of the mean; it does not give numerical values or an explicit
power law.

Let \(B_n=\{\mathsf m_0:F_n(d_n,\mathsf m_0)>\eta_n\}\).
Markov's inequality and countable subadditivity give
\[
 \rho(B_n)\le2^{-n},
 \qquad
 \rho\left(\bigcup_{n\ge N}B_n\right)\le2^{1-N}.
\]
The decreasing intersection of these tail unions consequently has
measure zero. Add the countable union of the imprint zero sets. The
resulting exceptional set is Lebesgue-null because \(\rho\) and Lebesgue
measure have the same null sets. Outside it, for every sufficiently large
\(n\), all \(0<\varepsilon\le d_n\) satisfy
\(F_n(d_n,\mathsf m_0)\le\eta_n\). Since \(\eta_n<1\), clipping has
no effect on this inequality: the row and total-variation errors and both entropy majorants are at
most \(\eta_n\). The majorants bound the actual entropies, giving
all five stated conclusions.

The frozen matching probability is, exactly,
\[
 \mu_{\mathrm T}(M)
 =\frac1{2+8\mathsf r_E+4\mathsf r_N+2\mathsf r_C}
 =\frac1{2D_0(\mathrm T)}
 =\frac12+O(\mathrm T).
\]
Indeed \(\mathsf r_E,\mathsf r_N=O(\mathrm T)\) and
\(\mathsf r_C=O(\mathrm T^{3/2})\), as proved from the sheet
determinants. Pushforward cannot increase total variation. If
\(q^{S}\) denotes the assignment marginal, the two established
total-variation estimates therefore imply, uniformly on \(W_\varepsilon\),
\[
 \left|
 \operatorname{TV}
 (q_t^{\mathrm T_n,\varepsilon,\mathsf m_0,S},
                  q_t^{\mathrm T_n,0,\mathsf m_0,S})
 -\left(1-\frac1{2D_0(\mathrm T_n)}\right)
 \right|\le2\eta_n.
\]
Here the comparison distance is exactly
\(\operatorname{TV}(\delta_M,\mu_{\mathrm T_n}^S)
=1-\mu_{\mathrm T_n}(M)\), for either matching.
The assignment-law contrast thus tends to \(1/2\), even if the selected
matching changes with \(n\).

Finally, defining \(\mathrm T(\varepsilon)=\mathrm T_n\) on
\(d_{n+1}<\varepsilon\le d_n\) gives a single stepwise path with both
parameters tending to zero. The proved bounds hold for all sufficiently
small mobilities on this path. Arbitrary mobility caps are admissible;
for example \(h_n=\mathrm T_n^n\) forces
\(\varepsilon_n/\mathrm T_n^a\to0\) for each fixed \(a>0\) whenever
\(\varepsilon_n\le d_n\). This imposes an upper bound and does not certify
a prescribed equality between temperature and mobility.

Nothing in this argument makes the imprint sign eventually constant as
temperature vanishes. Its fixed-temperature nonvanishing and null zero
set give no quantitative lower bound at a fixed mean in that limit.
Accordingly the matching is allowed to depend on \(n\), and the result
concerns complete one-time laws, not complete histories or conditional
first-formation events. \(\square\)





\subsection*{An explicit temperature bound for tracking}

\propositionheading{Temperature-dependent entropy tracking}{prop:phil-temperature-tracking}
In the four-cycle of \cref{conj:phil-contrast-formation}, let
\(0<\mathrm T\le1/100\), \(p\in(0,1)^4\), and
\(p_{\min}=\min_i\{p_i,1-p_i\}\). Define
\[
 \mathrm K_0=210\sqrt{2\pi/5},\qquad
 \mathrm c_* =\min\left\{
 \frac2{1+4\mathrm K_0e^{13/8}},\frac1{\mathrm K_0e^{13/4}}\right\}.
\]
The entropy production of the full frozen law satisfies
\[
 \mathscr D_p(f)\ge
 \frac{\mathrm c_*\mathrm T^{5/2}}
 {4+3\log(1/\mathrm T)+8\log(1/p_{\min})}
 \operatorname{Ent}_{\mu_p}(f).
\]
Write \(\varepsilon=\epsilon_{\mathrm L}\),
\(t_{\mathrm L}=2\varepsilon^{-1}\log(1/\varepsilon)\), and
\[
 \mathrm D(\mathrm T,\varepsilon)
 =4+3\log(1/\mathrm T)+8\log2+32\log(1/\varepsilon),\qquad
 \ell=\frac{\mathrm c_*\mathrm T^{5/2}}
                {\mathrm D(\mathrm T,\varepsilon)}.
\]
For \(0<\varepsilon<e^{-1}\) and \(0\le t\le2t_{\mathrm L}\),
\begin{equation}
 \operatorname{TV}(q_t,\mu_{p(t)})
 \le e^{-\ell t/2}
       \sqrt{\frac{13(12+|\mathsf m_0|^2)}{4\mathrm T}}
    +\frac{2\varepsilon\mathrm D(\mathrm T,\varepsilon)}
                 {\mathrm c_*\mathrm T^{7/2}}.
\label{eq:phil-temperature-tracking}
\end{equation}
In particular, for \(\varepsilon=\mathrm T^a\) with any \(a>7/2\),
this total variation tends to zero uniformly on
\([t_{\mathrm L},2t_{\mathrm L}]\) and on bounded sets of initial means.
This conclusion compares one-time laws; it supplies no imprint lower bound.
\status{ESTABLISHED}

\paragraph{Proof.}
Use \(u=(\bar k+\bar m)/\sqrt2\),
\(v=(\bar k-\bar m)/\sqrt2\), and \(\omega\), and put
\(R=[D,-\sqrt2I]\). The common kernel
\(K=\ker R\oplus\ker D\) has dimension five. Every
\(\mathsf H_\sigma\) vanishes there, and the reference precision preserves
\(K\) and \(N=K^\perp\). Its restrictions have eigenvalues one and five.
The passive process on \(K\) is therefore an independent Gaussian
Ornstein--Uhlenbeck process with entropy contraction at least \(2\mathrm T\).
The active process on \(N\), together with the assignment, is autonomous:
all redraw costs depend only on its residuals. Rescaling its coordinate by
\(\mathrm T^{-1/2}\) gives unit diffusion and drift matrix
\(a_\sigma=(\mathsf H_\sigma+\mathrm T\mathsf\Omega)|_N\).
Since every edge multiplicity is at most two,
\[
 \mathrm TI\preceq a_\sigma\preceq13I.
\]
Indeed \(RR^\top=DD^\top+2I\) has eigenvalues \(2,4,4,6\),
while \(D^\top D\) has eigenvalues \(0,2,2,4\), so
\(\|\mathsf H_\sigma\|\le12\).

Divide the four determinants in
\cref{eq:phil-selection-sheet-determinants} by the passive determinant
\(5\mathrm T^5\). Their positive coefficients, evaluated at
\(\mathrm T=1/100\) for upper bounds, give the active bounds
\[
 \begin{aligned}
 5120\mathrm T^3&\le d_M\le5262\mathrm T^3,&
 896\mathrm T&\le d_E\le1017\mathrm T,\\
 896\mathrm T&\le d_N\le995\mathrm T,&
 3072&\le d_C<5120.
 \end{aligned}
\]
Thus \(w_M\) and \(w_C\) are the largest and smallest active sheet
weights, and \(w_M\le\mathrm T^{-3/2}/\sqrt{5120}\),
\(w_C/w_M\ge\mathrm T^{3/2}\). The finite-label constant in the
entropy proof consequently satisfies
\[
 L_p\le4+3\log(1/\mathrm T)+8\log(1/p_{\min}).
\]

Apply the Gaussian-sheet and prior-weighted-ball argument of
\cref{eq:phil-selection-gaussian-inequalities,eq:phil-selection-entropy-contraction}
in the seven active coordinates, using the ball of radius \(1/2\), whose
volume is \(\pi^3/840\). Its Gaussian comparison factor is at most
\(4\mathrm K_0e^{13/8}\mathrm T^{-3/2}\), because
\[
 \frac{(2\pi)^{7/2}}{(\pi^3/840)\sqrt{5120}}=\mathrm K_0.
\]
More explicitly, the two energy coefficients in that argument obey
\[
 C_c\le(2+8\mathrm K_0e^{13/8})L_p\mathrm T^{-5/2},\qquad
 C_j\le4\mathrm K_0e^{13/4}L_p\mathrm T^{-5/2}.
\]
The continuous production equals four times its quadratic form and the
jump production bounds four times its quadratic form. Hence the active
rate is at least \(\mathrm c_*\mathrm T^{5/2}/L_p\). Entropy decomposes
over the independent passive and active factors. Cauchy--Schwarz bounds
the energy of a marginal by the corresponding integrated energy, so the
full rate is at least the smaller of the two rates. The displayed active
lower bound is less than \(2\mathrm T\), proving the full-state inequality.
The form-domain extension uses the positive regularization and cutoffs
already used for \cref{eq:phil-selection-entropy-contraction}.

A one-receiver flip joins only the class pairs \(M\)--\(E\),
\(N\)--\(E\), and \(C\)--\(E\). The determinant bounds above make every
ratio of such adjacent weights at most \(1/\mathrm T\).
Writing the evidence as \(p_iZ_i^++(1-p_i)Z_i^-\), the ratio
\(Z_i^+/Z_i^-\) is between the extremal adjacent ratios. Therefore
\(|\partial_{p_i}\log Z|\le1/\mathrm T\), and
\(\|\mathfrak g\|_1\le4/\mathrm T\).
The row lower bound gives the rate floor \(\ell\) on the stated horizon.
The initial entropy is at most
\(13(12+|\mathsf m_0|^2)/(2\mathrm T)\): in the Gaussian density ratio,
\(\mathsf B_\sigma\preceq13I/\mathrm T\),
\(\mathsf Z\le1\), and the remaining quadratic term has nonpositive
expectation. Insert these bounds in
\cref{eq:phil-selection-tracking} and use Pinsker's inequality to obtain
\cref{eq:phil-temperature-tracking}.
For \(\varepsilon=\mathrm T^a\), its second term is
\(O(\mathrm T^{a-7/2}\log(1/\mathrm T))\), while
\(\ell t_{\mathrm L}\) grows as a positive constant times
\(\mathrm T^{5/2-a}\). The exponential therefore dominates the initial
\(\mathrm T^{-1/2}\) factor, proving the last assertion. \(\square\)
\section{Coherence and Persistence of the Selected Pairing}
\label{sec:phil-selected-coherence-proof}

\propositionheading{Coherence and persistence after pairing selection}{prop:phil-selected-pair-persistence}
In the four-cycle of \cref{conj:phil-contrast-formation}, let
\(0<\mathrm T\le1/100\), let \(M\) be one matching, and suppose
\([t_0,t_0+\mathsf h]\subseteq[t_{\mathrm L},2t_{\mathrm L}]\), with
\(\mathsf h>0\). Write
\(\nu_{\mathrm T}^M=\mathcal N(0,\mathrm T\mathsf A_M^{-1})\otimes\delta_M\).
Assume the learned and frozen recognition laws satisfy
\[
 \sup_{t_0\le t\le t_0+\mathsf h}
 \operatorname{TV}(q_t^{\mathrm L},\nu_{\mathrm T}^M)\le\delta,
 \qquad
 \operatorname{TV}(q_{t_0}^{0},\mu_{\mathrm T,1/2})\le\delta.
\]
Let \(\mathsf O_{\mathrm T,\mathsf h}^M\) be the stationary
Ornstein--Uhlenbeck path law with drift \(-\mathsf A_M\mathbf x\),
diffusion covariance \(2\mathrm TI\), and constant assignment \(M\).
For the configuration and assignment path laws on this interval,
\begin{align}
 \operatorname{TV}(\mathsf P^{\mathrm L}_{t_0,\mathsf h},
                    \mathsf O_{\mathrm T,\mathsf h}^M)
 &\le(\mathsf h+2)\delta,
 \label{eq:phil-selected-path-coupling}\\
 \operatorname{TV}\left(\mathsf P^{0}_{t_0,\mathsf h},
   \tfrac12\mathsf O_{\mathrm T,\mathsf h}^{M_0}
        +\tfrac12\mathsf O_{\mathrm T,\mathsf h}^{M_1}\right)
 &\le\delta+16\mathrm T+31\mathrm T\mathsf h.
 \label{eq:phil-frozen-path-mixture}
\end{align}
The frozen arm keeps its original uniform rows.

For resolution \(\mathsf a>0\), let
\(\mathscr E_M(\mathsf a,\mathsf h)\) require throughout the interval
that the assignment is \(M\), the stacked residual norm
\(\sqrt{(r^b)^2+(r^m)^2}\) on its two edges is at most
\(\mathsf a\), and that norm on each exterior edge exceeds
\(\mathsf a\).
Define
\[
 \mathsf b(\mathrm T,\mathsf a,\mathsf h)
 =8\sqrt{2/\pi}\mathsf a+
            (8+3400\mathsf h)\mathrm T/\mathsf a^2.
\]
Then
\begin{equation}
 \Pr_{\mathrm L}(\mathscr E_M^c)
 \le(\mathsf h+2)\delta+\mathsf b(\mathrm T,\mathsf a,\mathsf h).
 \label{eq:phil-selected-coherence-persistence}
\end{equation}
In particular, the choice
\(\mathsf a=\mathrm T^{1/4}\),
\(\mathsf h=\mathrm T^{-1/4}\), and \(\delta\le\mathrm T\)
makes this probability tend to zero.
The matching-window label, equal to \(M_j\) when
\(\mathscr E_{M_j}\) holds and to an additional symbol otherwise,
converges to \(\delta_M\) in the learned arm and to the equal matching
mixture in the frozen arm.

At one time, define \(\mathfrak P_{\mathsf a}(\mathbf x)\) as the
connected-component partition of edges whose stacked residual norm
\(\sqrt{(r^b)^2+(r^m)^2}\) is at most \(\mathsf a\).
For times when both
one-time hypotheses above hold, put
\[
 \mathsf e(\mathrm T,\mathsf a)
 =4\sqrt{2/\pi}\mathsf a+4\mathrm T/\mathsf a^2.
\]
The learned and frozen partition laws satisfy
\begin{equation}
 \left|\operatorname{TV}
   (q_t^{\mathrm L}\mathfrak P_{\mathsf a},
                 q_t^0\mathfrak P_{\mathsf a})-\tfrac12\right|
 \le2\delta+16\mathrm T+2\mathsf e(\mathrm T,\mathsf a).
 \label{eq:phil-coherent-partition-contrast}
\end{equation}
Thus the contrast persists for the coherent partition, not only the
source label.

These conclusions apply along \cref{cor:phil-joint-selection} with
tolerance \(\delta\le\mathrm T\), mobility cap
\(\epsilon_{\mathrm L}\le\mathrm T^2\), and \(t_0=t_{\mathrm L}\).
The matching can vary along that schedule. The internal restoring rates
are \(8+\mathrm T\) and \(4+5\mathrm T\), on an event whose probability
tends to one. The resolution therefore has restoring ratio at least
\(4\mathsf a^2/\mathrm T\to\infty\).
If the rows are also within \(\delta\) of their matching vertex,
the stationary law with those rows frozen realizes the same persistence
event with probability at least
\(1-4(\mathsf h+2)\delta-\mathsf b\).
This late coherence is restored equilibrium coherence; it supplies no
entrance or first-formation assertion.
\status{ESTABLISHED}

\noindent\emph{Proof.}
First retain the proposal mechanism in the comparison. Put
\(r(t)=q_t^{\mathrm L}(S\ne M)\), and let
\(b_{\ne M}(t,\mathbf x)\) be the probability of proposing a different
assignment at a clock ring. Since the clock has rate one,
\[
 r'(t)=\E b_{\ne M}(t,\mathbf X_t)-r(t).
\]
If \(N_{\ne M}\) counts proposals to an assignment other than \(M\)
on the interval, its compensator gives
\begin{equation}
 \E N_{\ne M}
 =r(t_0+\mathsf h)-r(t_0)+
                   \int_{t_0}^{t_0+\mathsf h}r(t)dt.
 \label{eq:phil-selected-proposal-occupation}
\end{equation}
Consequently
\(\Pr(\exists t\in[t_0,t_0+\mathsf h]:S_t\ne M)
 \le(\mathsf h+1)\delta\).
Couple the initial law maximally with \(\nu_{\mathrm T}^M\), and use
the same future Brownian motion for the native process and the
stationary matching process. Their paths agree until an initial
mismatch or a proposal to another assignment. The union bound and
\cref{eq:phil-selected-proposal-occupation} prove
\cref{eq:phil-selected-path-coupling}. This argument counts proposals
of the native process and does not condition its diffusion on an
assignment history.

For the frozen comparison, use the normalized weights
\(1,\mathsf r_E,\mathsf r_N,\mathsf r_C\) from
\cref{sec:phil-pairing-selection-proof}. The determinant bounds in the
proof of \cref{prop:phil-temperature-tracking} give
\[
 \mathsf r_E,\mathsf r_N\le\tfrac52\mathrm T,
 \qquad \mathsf r_C\le\tfrac32\mathrm T^{3/2}.
\]
The invariant mass \(\kappa_{\mathrm T}\) of nonmatching assignments
is therefore at most \(16\mathrm T\). Conditional on either matching,
its configuration law is the Gaussian factor of \(\nu_{\mathrm T}^M\).
The invariant law is consequently within \(\kappa_{\mathrm T}\) of
the equal mixture of these two laws.

Detailed balance bounds the stationary flux from \(M\) to all
nonmatching assignments by \(\kappa_{\mathrm T}\).
For the other matching \(M'\), write \(\mu_M(\mathbf x)\) for the
joint invariant density on the matching sheet. Pointwise,
\[
 \mu_M b_{M'}
 =\frac{\mu_M\mu_{M'}}{\sum_S\mu_S}
 \le\tfrac12\sqrt{\mu_M\mu_{M'}}.
\]
Since \((\mathsf A_M+\mathsf A_{M'})/2\) is the oriented-cycle
matrix, Gaussian integration gives
\(\int\sqrt{\mu_M\mu_{M'}}=\mu_{\mathrm T,1/2}(M)\mathsf r_C\).
Dividing the two flux bounds by the matching mass yields
\[
 \E_{\nu_{\mathrm T}^M}b_{\ne M}
 \le8\mathsf r_E+4\mathsf r_N+\tfrac52\mathsf r_C
 \le31\mathrm T.
\]
Thus the stationary flux out of the two matching sheets together is
at most \(31\mathrm T\). Couple the stationary native process at
its initial time with the equal matching mixture and use common
Brownian increments until the first departure from its coupled
matching. Its failure probability is at most
\(16\mathrm T+31\mathrm T\mathsf h\). The frozen native process
started at \(q_{t_0}^0\) differs from its stationary path law by at
most its initial total variation \(\delta\), because the frozen
transition kernel is the same. This proves
\cref{eq:phil-frozen-path-mixture}.

It remains to bound residuals on the comparison path.
The four normal coordinates
\(z_g^\pm=n_g^\pm\cdot\mathbf X\), \(g\in M\), from
\cref{eq:rg-formation-cycle-normals} have rates
\(8+\mathrm T\) and \(4+5\mathrm T\). Their stationary variances
are at most \(\mathrm T/4\), and
\[
 r_g^b=\sqrt2z_g^++z_g^-,\qquad
 r_g^m=\sqrt2z_g^+-z_g^-.
\]
For a scalar Ornstein--Uhlenbeck noise starting at zero, integration
by parts in its stochastic convolution gives
\[
 Z_t=\sqrt{2\mathrm T}
 \left[W_t-\lambda\int_0^te^{-\lambda(t-u)}W_u du\right],
 \qquad
 \E\sup_{t\le\mathsf h}|Z_t|^2\le32\mathrm T\mathsf h.
\]
The last inequality uses
\(\sup|Z_t|\le2\sqrt{2\mathrm T}\sup|W_t|\) and the Brownian
maximal estimate proved in the matching-episode argument of
\cref{sec:rg-formation-cycle-proof}.
For the four-dimensional stationary normal vector \(z\), it follows
that
\[
 \E\sup_{t\le\mathsf h}|z_t|^2
 \le2\mathrm T+256\mathrm T\mathsf h.
\]
Each selected stacked residual norm is at most \(2|z_t|\), since
\((r_g^b)^2+(r_g^m)^2=4(z_g^+)^2+2(z_g^-)^2\).
Thus the probability that any internal stacked norm exceeds
\(\mathsf a\) at any time is at most
\((8+1024\mathsf h)\mathrm T/\mathsf a^2\).

Every exterior link coordinate is absent from \(H_M\), and its
block in \(\mathsf A_M\) is the scalar \(\mathrm T\).
It is therefore an independent \(\mathcal N(0,1)\) coordinate in
the stationary Gaussian. Each exterior belief or model residual
has variance at least one. The probability that any of these four
residuals initially has absolute value at most \(2\mathsf a\)
is at most \(8\sqrt{2/\pi}\mathsf a\).

To control their increments, diagonalize the symmetric positive
matrix \(\mathsf A_M\). In a scalar eigencoordinate with rate
\(\lambda>0\), stationarity and the preceding noise estimate imply
\[
 \E\sup_{t\le\mathsf h}|X_t-X_0|^2
 \le\frac{2\mathrm T}{\lambda}(1-e^{-\lambda\mathsf h})^2
                       +64\mathrm T\mathsf h
 \le66\mathrm T\mathsf h,
\]
since \((1-e^{-v})^2\le v\) for \(v\ge0\).
Summing twelve coordinates gives
\begin{equation}
 \E\sup_{t\le\mathsf h}|\mathbf X_t-\mathbf X_0|^2
 \le792\mathrm T\mathsf h.
 \label{eq:phil-selected-ou-increments}
\end{equation}
Every residual covector has norm \(\sqrt3\). Except with
probability \(2376\mathrm T\mathsf h/\mathsf a^2\), all exterior
increments are therefore at most \(\mathsf a\) in absolute value.
The initial small-ball estimate and the two maximal estimates prove
\(\mathsf O_{\mathrm T,\mathsf h}^M(\mathscr E_M^c)\le\mathsf b\).
Together with \cref{eq:phil-selected-path-coupling}, this proves
\cref{eq:phil-selected-coherence-persistence}.
The frozen path comparison proves the stated matching-window label
limit whenever \(\mathsf b\to0\),
\(\mathsf h\delta\to0\), and \(\mathrm T\mathsf h\to0\), in
particular for the displayed choices.

For a single time no increment estimate is needed. The normal second
moment is at most \(\mathrm T\), so the internal failure probability
is at most \(4\mathrm T/\mathsf a^2\). The exterior small-ball
bound at radius \(\mathsf a\) is
\(4\sqrt{2/\pi}\mathsf a\). Thus the partition induced by
\(\nu_{\mathrm T}^{M_j}\) differs from the matching partition
\(\mathfrak P_j\) with probability at most \(\mathsf e\).
The learned partition law is within \(\delta+\mathsf e\) of
\(\delta_{\mathfrak P_M}\); the frozen one is within
\(\delta+16\mathrm T+\mathsf e\) of
\((\delta_{\mathfrak P_0}+\delta_{\mathfrak P_1})/2\).
These matching partitions are distinct, so the triangle inequality
for total variation proves \cref{eq:phil-coherent-partition-contrast}.

The schedule in \cref{cor:phil-joint-selection} permits tolerance
\(\mathrm T\) and mobility cap \(\mathrm T^2\).
Then \(\mathrm T^{-1/4}\le t_{\mathrm L}\) for sufficiently small
\(\mathrm T\), so the following persistence window lies within the
selection window. All errors just bounded tend to zero.
The normal rates establish the restoring ratio on the native event.
Finally, if \(\max_i|p_i-p_i^M|\le\delta\),
\(\operatorname{TV}(\mu_{\mathrm T,p},\nu_{\mathrm T}^M)
 =1-\mu_{\mathrm T,p}(M)\le4\delta\).
Applying the same occupation and coupling argument to this
stationary frozen-row process gives the stationary lower bound in
the statement. The persistence event is therefore already typical
at equilibrium for those rows and carries no first-entrance claim.
\hfill\(\square\)



\subsection{The first contrast on the original episode window}

\propositionheading{A first-block null under slower source learning}{prop:phil-early-block-null}
In the model of \cref{conj:phil-contrast-formation}, let
\(0<\mathrm T\le10^{-16}\), and let the initial law be
\(\mathcal N(\mathsf m_0,I_{12})\) times uniform assignments.
The learned and frozen arms have the same initial mean and differ only in
the source-learning mobility \(\epsilon_{\mathrm L}\).
Set
\[
 H_{\mathrm T}=2\mathrm T^{-1/8}.
\]
Keep the contrast event of \cref{prop:phil-seeded-contrast-formation}
with its regime parameter equal to \(\mathrm T^{1/8}\): an eligible
proper connected block has an internal belief relative entropy exceeding
\(1/4\) initially, and for a duration-one interval inside
\([0,H_{\mathrm T}]\) all its internal belief and model relative
entropies are at most \(\sqrt{\mathrm T}\), while its boundary belief
relative entropies exceed \(1/2\).
For each block take the infimum of its qualifying completion times,
with value infinity when there is no such interval. Let \(K_{\mathrm T}\)
be the block with least such infimum, using one fixed tie rule in both
arms, and give it a cemetery value when no block qualifies. Write
\(E_{\mathrm T}\) for the event that its value is not the cemetery value.
The infimum need not itself be a qualifying completion time.

Define the positive, temperature-independent constant
\[
 \mathsf c(\mathsf m_0)
 =\frac{e^{-1/4}\varrho_\circ^{12}}{737280}
   \exp\left[-\frac12
       (|\mathbf x^\circ-\mathsf m_0|+\varrho_\circ)^2\right],
 \qquad
 \mathsf d_{\mathrm T}
 =\min\{1,8\epsilon_{\mathrm L}\mathrm T^{-1/4}\}.
\]
Then
\[
 \Pr_0(E_{\mathrm T})\ge\mathsf c(\mathsf m_0),
 \qquad
 \Pr_{\epsilon_{\mathrm L}}(E_{\mathrm T})
       \ge\mathsf c(\mathsf m_0)-\mathsf d_{\mathrm T}.
\]
If \(\mathsf d_{\mathrm T}<\mathsf c(\mathsf m_0)\), the two
conditional block laws satisfy
\[
 \operatorname{TV}\left(
 \operatorname{Law}_{\epsilon_{\mathrm L}}(K_{\mathrm T}\mid E_{\mathrm T}),
 \operatorname{Law}_0(K_{\mathrm T}\mid E_{\mathrm T})\right)
 \le\frac{\mathsf d_{\mathrm T}}{\mathsf c(\mathsf m_0)}.
\]
In particular their difference vanishes for every fixed Gaussian mean
whenever
\(\mathrm T\downarrow0\) and
\(\epsilon_{\mathrm L}\mathrm T^{-1/4}\to0\).
This includes the caps
\(\epsilon_{\mathrm L}\le\mathrm T_n^a\), \(a>7/2\), in
\cref{cor:phil-joint-selection}.
\status{ESTABLISHED}

\noindent\emph{Proof.}
First prove a uniform comparison of the complete native path laws on
any finite horizon \(H\). The learned rows are deterministic and satisfy
\[
 |p_i(t)-1/2|\le\epsilon_{\mathrm L}t/2.
\]
At a common configuration, write \(A,B>0\) for a receiver's likelihood
factors for its two sources. Its posterior positive-source probabilities
in the two arms are
\[
 b(p)=\frac{pA}{pA+(1-p)B},
 \qquad b(1/2)=\frac A{A+B}.
\]
Direct subtraction gives
\[
 |b(p)-b(1/2)|
 =\frac{AB|2p-1|}{(pA+(1-p)B)(A+B)}
 \le|2p-1|.
\]
Indeed the denominator is
\(pA^2+AB+(1-p)B^2\ge AB\). This bound is independent of
temperature and the likelihood ratio, without an interior-row restriction.

Use identical initial samples, a common Brownian motion, and the same
rate-one proposal clock in the two arms. Until the first different redraw,
their configurations and assignments coincide: a held-sheet drift and its
noise do not depend on the prior rows. At each ring, couple each receiver's
two Bernoulli laws maximally. The union bound and the row estimate bound
the probability of any different receiver choice at time \(t\) by
\[
 \sum_i|b_i(p_i(t))-b_i(1/2)|
 \le4\epsilon_{\mathrm L}t.
\]
The compensator of the rate-one clock, stopped at the first mismatch,
bounds its probability before \(H\) by
\(\int_0^H4\epsilon_{\mathrm L}t dt=2\epsilon_{\mathrm L}H^2\).
Consequently,
\begin{equation}
 \operatorname{TV}\left(
 \operatorname{Law}_{\epsilon_{\mathrm L}}(\mathbf X_{[0,H]},S_{[0,H]}),
 \operatorname{Law}_0(\mathbf X_{[0,H]},S_{[0,H]})\right)
 \le\min\{1,2\epsilon_{\mathrm L}H^2\}.
\label{eq:phil-early-path-coupling}
\end{equation}
The bound follows because the complete paths are identical unless the
coupling has had a mismatch. Taking \(H=H_{\mathrm T}\) gives
\(\mathsf d_{\mathrm T}\).

The first-block mark is a measurable function of these paths. To see this
directly, for a fixed block let \(A_I(t)\) be the maximum of its internal
belief and model relative entropies on \([t-1,t]\), and let \(B_I(t)\)
be the minimum boundary belief relative entropy there. Configuration
paths are continuous, so these two window functions are continuous.
Replace the strict condition \(B_I(t)>1/2\) by the countable union of
conditions \(B_I(t)\ge1/2+1/k\), \(k\ge1\). For each \(k\) the
qualifying end times form a closed subset of the compact interval
\([1,H_{\mathrm T}]\). Its first time is measurable, as may be checked
by minimizing
\(\max(A_I(t)-\sqrt{\mathrm T},1/2+1/k-B_I(t))\) on each compact
initial subinterval. The infimum over \(k\), the initial eligibility
condition, and the finite comparison over blocks preserve measurability.
This construction implements the strict boundary convention without
asserting that its infimum is attained.

A measurable map of paths does not increase total variation. Let
\(a_I,b_I\) be the probabilities of the non-cemetery marks in the learned
and frozen arms, \(p=\sum_Ia_I\), \(q=\sum_Ib_I\), and \(d\) the
total-variation distance of the two cemetery-augmented mark laws. Then
\[
 2d=\sum_I|a_I-b_I|+|p-q|,
 \qquad d\le\mathsf d_{\mathrm T}.
\]
If \(p\ge q>0\), the triangle inequality gives
\[
 \frac12\sum_I\left|\frac{a_I}{p}-\frac{b_I}{q}\right|
 \le\frac{\sum_I|a_I-b_I|+(p-q)}{2p}
 =\frac d p.
\]
Interchanging the arms proves the bound \(d/\max(p,q)\).
The frozen matching-seed sub-event in
\cref{prop:phil-seeded-contrast-formation} belongs to \(E_{\mathrm T}\)
and has probability at least \(\mathsf c(\mathsf m_0)\).
Thus \(q\ge\mathsf c(\mathsf m_0)\) and
\(p\ge q-d\ge\mathsf c(\mathsf m_0)-\mathsf d_{\mathrm T}\).
The claimed conditional bound follows, and its limiting conclusion is
immediate. \(\square\)

The proposition keeps the original contrast event and uses its seeded
sub-event only as a probability lower bound. That sub-event already has
the restoring interaction and stationary comparison of
\cref{prop:rg-formation-endogenous-cycle}; the proof does not certify
every history in the broader contrast event under the coherence-episode
rule. Any common measurable admissibility restriction retaining the
seeded sub-event has the same comparison.

The result does not cover the original coupled choice
\(\epsilon_{\mathrm L}=\mathrm T^{1/8}\), for which its bound is
vacuous. It also does not compare different initial means or establish a
conditional parent error. On the slower schedules, late assignment
selection can coexist with a vanishing difference between the conditional
laws of the first contrast on \([0,H_{\mathrm T}]\). These are
different observables on different horizons.



\section{Activation from the Actual Conditioned Recognition Law}
\label{sec:phil-actual-conditioned-parent}

\propositionheading{Conditional information after a completed coherence prefix}{prop:phil-actual-conditioned-parent}
In the learned four-cycle, fix \(0<\mathrm T\le1/100\), a selected
matching \(M\), and \(s=t_{\mathrm L}\ge1\). Suppose, throughout
\([s,2s]\),
\[
 \KL(q_t\Vert\mu_{\mathrm T,p(t)})\le\eta,
 \qquad
 \operatorname{TV}(q_t,\nu_{\mathrm T}^M)\le\delta.
\]
Set \(\mathsf a=\mathrm T^{1/4}\), \(\theta=s+1\), and let
\(E=\mathscr E_M(\mathsf a,1)\) be the actual matching and coherence
event of \cref{prop:phil-selected-pair-persistence} on \([s,s+1]\).
Assume
\[
 \rho=3\delta+\mathsf b(\mathrm T,\mathsf a,1)<\tfrac12,
 \qquad
 h_b(r)=-r\log r-(1-r)\log(1-r),
 \qquad
 \kappa_E=\frac{\eta+h_b(\rho)}{1-\rho}.
\]
For all \(t\ge\theta\), write
\(q_t^E=\operatorname{Law}(\mathbf X_t,S_t\mid E)\), always
conditioning on this same completed prefix. Keep the single marked
channel \(C\) of \cref{prop:rg-formation-endogenous-cycle}(b), including
its identity action off the matching strata and its identity extension
on evaluator samples. The parent activation law is exactly
\[
 \mu_{\mathrm{act}}=C_\#q_\theta^E.
\]

Let \(\mathcal H_{\mathrm{opt}}(q)\) be the mean discarded conditional
relative entropy of \(q\) against the restricted optimizer under \(C\).
Let \(\mathcal H_{\mathrm{phys}}(q)\) be the corresponding mean
discarded conditional relative entropy against the full physical
posterior, after lifting \(q\) with its recognition evaluator and
retaining every evaluator sample. Then, for every \(t\in[\theta,2s]\),
\begin{align}
 \KL(q_t^E\Vert\mu_{\mathrm T,p(t)})&\le\kappa_E,
 \label{eq:phil-conditioned-full-entropy}\\
 \mathcal H_{\mathrm{opt}}(q_t^E)&\le\kappa_E,
 \label{eq:phil-conditioned-opt-entropy}\\
 \mathcal H_{\mathrm{phys}}(q_t^E)
 &\le\kappa_E+\log\frac{4+5\mathrm T}{4+\mathrm T}
 \le\kappa_E+\mathrm T.
 \label{eq:phil-conditioned-physical-entropy}
\end{align}
The two hidden-conditional bounds hold with any strictly positive
alternative source rows in their reference laws; the full-law bound
\cref{eq:phil-conditioned-full-entropy} refers to the native rows
\(p(t)\).

These are unscaled, dimensionless relative entropies. The corresponding
scaled optimizer free-energy loss is at most \(\mathrm T\kappa_E\);
the unscaled physical negative-ELBO loss is at most
\(\kappa_E+\mathrm T\), and its \(\mathrm T\)-scaled value is at most
\(\mathrm T\kappa_E+\mathrm T^2\).
Along the joint selection schedule with \(\eta,\delta\le\mathrm T\),
\[
 \rho=O(\mathrm T^{1/4}),
 \qquad
 \kappa_E=O(\mathrm T^{1/4}\log(1/\mathrm T)).
\]
The assumption \(\rho<1/2\) then holds for all sufficiently small
\(\mathrm T\), so these bounds are asymptotic.
\status{ESTABLISHED}

\noindent\emph{Proof.}
The learned-arm persistence bound
\cref{eq:phil-selected-coherence-persistence}, applied with
\(t_0=s\) and \(\mathsf h=1\), gives \(r:=\Pr(E^c)\le\rho\); its
proof uses only the learned-arm hypothesis. Thus \(E\) is a
\(\theta\)-measurable event of probability at least \(1-\rho\).
The law \(q_\theta^E\) is the actual conditional law and may have
retained--hidden correlations and a non-Gaussian hidden conditional.
Its marked pushforward is normalized by construction. At later times
no further event is conditioned upon. In particular, the deterministic
native row path is not recomputed from \(q_t^E\).

For a fixed \(t\in[\theta,2s]\), put
\(B=\mathbf1_E\) and \(R_t=(\mathbf X_t,S_t)\).
Compute the relative entropy of the joint law of \((B,R_t)\) against
\(\mathrm{Bernoulli}(1-r)\otimes\mu_{\mathrm T,p(t)}\) in two orders:
\begin{equation}
 \begin{aligned}
 &(1-r)\KL(q_t^E\Vert\mu_{\mathrm T,p(t)})
       +r\KL(q_t^{E^c}\Vert\mu_{\mathrm T,p(t)})\\
 &\qquad=\KL(q_t\Vert\mu_{\mathrm T,p(t)})
                       +\mathbb I(B;R_t)
 \le\eta+h_b(r).
 \end{aligned}
 \label{eq:phil-conditioned-event-information}
\end{equation}
When \(r=0\), use just the first term. The second term is nonnegative,
\(h_b\) is increasing on \([0,1/2]\), and \(1-r\ge1-\rho\).
This proves \cref{eq:phil-conditioned-full-entropy}.
It also establishes the needed absolute continuity and finiteness;
directly, \(q_t^E\le q_t/(1-r)\).

The channel's fibers can be checked explicitly. On a matching with
edges \(e,f\), let \(F_M\) have the orthonormal columns
\[
 (n_e^+-n_f^+)/\sqrt2,\qquad n_e^-,\qquad n_f^-.
\]
Write \(Z_M=\operatorname{im}F_M\), \(W_M=Z_M^\perp\), and
\(\mathbf x=y+F_Mz\).
The normal identities in \cref{eq:rg-formation-cycle-normals} give
\[
 \nu_{\mathrm T,Z}=\mathcal N(0,\mathrm TK_Z^{-1}),
 \qquad
 K_Z=\operatorname{diag}(8+\mathrm T,4+5\mathrm T,4+5\mathrm T)
\]
as the hidden optimizer conditional. Both the interaction Hessian
and the evaluator precision preserve \(W_M\) and \(Z_M\).
Consequently this conditional is independent of the retained
coordinate and of the prior rows. It is a reference conditional, not
an assertion that the actual conditioned recognition law has that
form. The same statement holds on the other matching by rotation;
off the two matchings the channel is the identity and has no
discarded conditional.

The target of \(C\) is a finite disjoint union of Euclidean strata,
so disintegration gives the exact entropy chain rule
\begin{equation}
 \KL(q_t^E\Vert\mu_{\mathrm T,p(t)})
 =\KL(C_\#q_t^E\Vert C_\#\mu_{\mathrm T,p(t)})
                         +\mathcal H_{\mathrm{opt}}(q_t^E).
 \label{eq:phil-conditioned-parent-chain}
\end{equation}
Dropping its retained term proves
\cref{eq:phil-conditioned-opt-entropy}. This argument uses the
complete actual law; it neither resets the hidden state nor assumes
Gaussianity conditional on a source history.

For the physical bound, work first on one matching fiber. Its physical
outer hidden Gaussian is
\[
 \Pi_{\mathrm T,Z}=\mathcal N(0,\mathrm T(K_Z^{\mathrm{phys}})^{-1}),
 \qquad
 K_Z^{\mathrm{phys}}
  =\operatorname{diag}(8+\mathrm T,4+\mathrm T,4+\mathrm T).
\]
All displayed matrices are positive definite. Write
\(v=1-\sigma_{\mathrm{ev}}^2=1/2\),
\(d_y=\mathfrak Dy\), and \(d_z=\mathfrak DF_Mz\).
Since \(R_{\mathrm{ev}}=\mathfrak D^\top\mathfrak D/v\)
preserves the split, \(d_y\cdot d_z=0\); its hidden eigenvalues
are \(0,4,4\).

Let \(U=(m,k)\) denote the full evaluator sample vectors.
At configuration \(\mathbf x\), the physical evaluator has
\(m\sim\mathcal N(\bar m,\sigma_{\mathrm{ev}}^2I)\) and independent
increment \(k-m\sim\mathcal N(0,vI)\). Its recognition counterpart,
\cref{eq:rg-formation-evaluator-lift}, has the same first law and
increment mean \(d_y+d_z\). At a fixed retained \(y\), form the
comparison probability law
\begin{equation}
 \overline P_y(dz,dU)
 =\Pi_{\mathrm T,Z}(dz)P_{\mathrm{ev}}(dU\mid y+F_Mz)
  \exp\left[
  \frac{d_y\cdot(k-m)-|d_y|^2/2}{v}\right].
 \label{eq:phil-conditioned-evaluator-tilt}
\end{equation}
Completing the Gaussian square in \(k-m\) proves that the multiplier
integrates to one for every \(z\). It changes the increment mean
to \(d_y\) and keeps \(\Pi_{\mathrm T,Z}\) as its \(z\) marginal.
The multiplier is retained-measurable in \((y,U)\), so
\(\overline P_y\) has the physical hidden conditional given these
retained variables. At a fixed \(z\), the recognition evaluator
differs by the additional increment mean \(d_z\), with relative
entropy \(|d_z|^2/(2v)\).

Compute the joint relative entropy first through \(z\), then through
\(U\), and drop the nonnegative evaluator-marginal term. For the
actual hidden law \(q(dz\mid M,y)\), this gives
\[
 \mathcal H_{\mathrm{phys}}(q\mid M,y)
 \le\KL(q(dz\mid M,y)\Vert\Pi_{\mathrm T,Z})
             +\E_q[|d_z|^2/(2v)\mid M,y].
\]
All terms are well-defined: the native process has finite second
moments, conditioning increases the second-moment bound by a factor
at most \((1-\rho)^{-1}\), and the Gaussian reference logarithms are quadratic.
The optimizer conditional is the exact tilt
\[
 \nu_{\mathrm T,Z}(dz)
 =\frac{\Pi_{\mathrm T,Z}(dz)e^{-|d_z|^2/(2v)}}{Z_{\mathrm{tilt}}},
 \qquad
 Z_{\mathrm{tilt}}=\frac{4+\mathrm T}{4+5\mathrm T}.
\]
For its normalizer, the common coordinate is unchanged and each of
the two discrepancy coordinates changes rate from \(4+\mathrm T\)
to \(4+5\mathrm T\); their two square-root Gaussian determinant
factors multiply to \(Z_{\mathrm{tilt}}\). Therefore
\[
 \KL(q(dz\mid M,y)\Vert\Pi_{\mathrm T,Z})
       +\E_q[|d_z|^2/(2v)\mid M,y]
 =\KL(q(dz\mid M,y)\Vert\nu_{\mathrm T,Z})
                         +\log\frac{4+5\mathrm T}{4+\mathrm T}.
\]
Integrate over the retained law. Matching strata have total
probability at most one, and all other strata have zero discarded
loss. Finally,
\[
 0\le\log\frac{4+5\mathrm T}{4+\mathrm T}
 =\log\left(1+\frac{4\mathrm T}{4+\mathrm T}\right)
 \le\frac{4\mathrm T}{4+\mathrm T}\le\mathrm T.
\]
This proves \cref{eq:phil-conditioned-physical-entropy}.
Only recognition has been conditioned on \(E\); the physical target
still uses the original observed records, with all evaluator samples
retained.

A source-prior factor \(\pi_S\) cancels in any conditional law on a
fiber because \(S\) is retained. The hidden optimizer and physical
reference conditionals are consequently the same for every strictly
positive choice of prior rows. The full-law entropy need not remain
small for different rows, because its retained term may change.
An imposed zero prior on an occupied stratum would remove reference
support and is excluded; ordinary finite-time learning from positive
rows preserves positivity.

The exact free-energy chain rules give the stated scaled and unscaled
losses. With \(\eta,\delta\le\mathrm T\) and
\(\mathsf a=\mathrm T^{1/4}\),
\[
 \rho\le3\mathrm T+8\sqrt{2/\pi}\mathrm T^{1/4}
                         +3408\mathrm T^{1/2}.
\]
Use \(h_b(\rho)\le\rho\log(e/\rho)\) and
\(1-\rho\ge1/2\) eventually to obtain the displayed order for
\(\kappa_E\). \(\square\)

\remarkheading{A high-probability prefix is needed}{rem:phil-conditioned-hidden-sign}
The high probability of the declared prefix has mathematical content.
Under the exact matching Gaussian \(\nu_{\mathrm T}^M\), a hidden
coordinate \(z_1\) is independent of the retained state and has zero
unconditioned hidden relative entropy. Conditioning on
\(\{z_1>0\}\), of probability \(1/2\), instead gives a density ratio
two on that half-space and hidden relative entropy \(\log2\) at every
temperature. Thus a probability merely bounded below does not suffice
for a vanishing conditional loss. Conditioning on the entire realized
fine history is different again: it fixes the hidden endpoint and
gives a Dirac conditional, singular relative to the nondegenerate
Gaussian fiber. The activation here uses the full joint law conditional
on \(E\), not a realized fine history.
\status{ESTABLISHED}

The proposition supplies conditional information at activation and for
the native future marginals under that same completed prefix. It does
not yet compare an independently evolving parent path. Its prefix is
the late restored coherence event, and carries no first-formation or
recursive-coalescence conclusion.



\subsection*{Inherited dynamics from the conditioned activation law}

\propositionheading{Dynamics of the conditioned parent}{prop:phil-actual-parent-dynamics}
Use the selected matching, channel, and conditioned activation of
\cref{prop:phil-actual-conditioned-parent}. More explicitly, let
\(0<\mathrm T\le1/100\), \(s=t_{\mathrm L}\), \(\theta=s+1\), and let \(\mathscr E\)
be the selected coherence event on \([s,\theta]\) at resolution
\(\mathsf a\) from \cref{prop:phil-selected-pair-persistence}.
Assume \(\Pr(\mathscr E^c)\le\rho<1\), and, throughout
\([s,\theta+\mathsf h]\subseteq[t_{\mathrm L},2t_{\mathrm L}]\),
\[
 \operatorname{TV}(q_t,\nu_{\mathrm T}^M)\le\delta,
 \qquad \max_i|p_i(t)-p_i^M|\le\delta.
\]
Write \(\varepsilon=\epsilon_{\mathrm L}\),
\(d=\delta+\varepsilon\mathsf h<1/4\), and use relative future time
\(u=t-\theta\). Initialize the parent with
\[
 q_0^C=C_\#\operatorname{Law}(\mathbf X_\theta,S_\theta\mid\mathscr E),
 \qquad p^C(0)=p(\theta).
\]
On the stratum \(\{S\}\times W_S\), its continuous generator is
\(-\mathsf A_S y\cdot\nabla_{W_S}+\mathrm T\Delta_{W_S}\).
Its rate-one proposal kernel is
\begin{equation}
 J_p((S,y),\cdot)=
 \int\kappa_S(d\mathbf x\mid y)
     \sum_{S'}b_p(S'\mid\mathbf x)
                   \delta_{(S',P_{W_{S'}}\mathbf x)},
 \label{eq:phil-actual-parent-kernel}
\end{equation}
where \(b_p(S'\mid\mathbf x)=\mu_{\mathrm T,p}(S'\mid\mathbf x)\) is the
product of receiver posterior rows at rows \(p\), and
\(\kappa_S\) is the optimizing hidden conditional: the three
independent Gaussian coordinates on a matching stratum, and a point
mass off the matchings. The parent rows evolve from its own law:
\begin{equation}
 \frac d{du}p_i^C(u)
 =\varepsilon\bigl(q_u^C(\sigma_i=+1)-p_i^C(u)\bigr).
 \label{eq:phil-actual-parent-learning}
\end{equation}
This coupled parent system has a unique global law with interior rows.

Let \(\mathsf P_{\theta,\mathsf h}^{\mathrm{fine},C}\) be the law of
\((C(\mathbf X_{\theta+u},S_{\theta+u}))_{0\le u\le\mathsf h}\)
conditioned only on \(\mathscr E\), with the fine process continuing
its original row path. Let \(\mathsf P_{\mathsf h}^C\) be the parent
path law from the displayed activation and its own learning equation.
Then
\begin{equation}
 \operatorname{TV}
 (\mathsf P_{\theta,\mathsf h}^{\mathrm{fine},C},
                                      \mathsf P_{\mathsf h}^C)
 \le2\rho+(\mathsf h+3)\delta+
                         \frac{4\mathsf h d}{1-4d}.
 \label{eq:phil-actual-parent-path-error}
\end{equation}
The comparison includes all retained coordinates and the assignment
history, hence also the loop history. Self-proposals remain part of
the rate-one kernel.

Along \cref{cor:phil-joint-selection}, choose
\(\delta\le\mathrm T\), \(\varepsilon\le\mathrm T^2\),
\(\mathsf a=\mathrm T^{1/4}\), and
\(\mathsf h=\mathrm T^{-1/4}\). One may take
\(\rho=3\delta+\mathsf b(\mathrm T,\mathsf a,1)\), where
\(\mathsf b\) is from \cref{prop:phil-selected-pair-persistence}.
The right side of \cref{eq:phil-actual-parent-path-error} is then
\(O(\mathrm T^{1/4})\). This is a conditional, finite-window
approximation after the selected coherence event; the activation law
is not replaced by a Gaussian law.
\status{ESTABLISHED}

\noindent\emph{Proof.}
For a matching \(M\), the decomposition
\(\R^{12}=W_M\oplus Z_M\) of
\cref{prop:rg-formation-endogenous-cycle} is preserved by
\(H_M\), \(\Omega\), and \(\mathsf A_M\).
The hidden drift matrix is
\(\operatorname{diag}(8+\mathrm T,4+5\mathrm T,4+5\mathrm T)\).
The prior weight \(\pi_S(p)\) cancels from the conditional law at
fixed \((S,y)\), so \(\kappa_S\) is independent of \(p\).
Off the matchings no coordinate is discarded. Thus the displayed
continuous generator is the projection of every held-sheet diffusion,
and \cref{eq:phil-actual-parent-kernel} is a probability kernel.
A self-proposal has \(S'=S\) and
\(P_{W_S}\mathbf x=y\), so it leaves the parent state unchanged.
At a different assignment it uses that assignment's projection,
including its twelve-dimensional stratum when appropriate.

For fixed rows, the pushed optimizing law \(C_\#\mu_{\mathrm T,p}\)
is invariant. Indeed each diffusion is reversible on its Gaussian
stratum, and
\(\mu_{\mathrm T,p}(d\mathbf x,S)b_p(S'\mid\mathbf x)\)
is symmetric in \(S,S'\). Pushing both endpoints through \(C\)
gives detailed balance for \(J_p\).
The loop covector annihilates each \(Z_M\), so this construction
retains the loop through the stratum changes as well.

We next check the existence assertion with evolving parent rows.
For any continuous interior row path, concatenate the linear
diffusions with the rate-one proposals, drawing \(\kappa_S\) at
each proposal. On every finite interval the clock has finitely many
rings and every inserted Gaussian is finite almost surely, so this
construction is nonexplosive. Its second moments remain finite:
projection cannot increase the norm, an inserted hidden draw adds
expected squared norm at most \(5\mathrm T/8\), and the diffusion
generator applied to \(|y|^2\) is at most
\(24\mathrm T-2\mathrm T|y|^2\). The full generator is consequently
bounded above by \(25\mathrm T-2\mathrm T|y|^2\).
The activation law has finite second moment because it conditions a
finite-moment fine law on an event of positive probability.

If two row paths have entries in \([\mathrm c_0,1-\mathrm c_0]\), couple their
parent processes by the same initial state, Brownian increments,
clocks, hidden Gaussian draws, and receiver uniforms until their
first different proposal. The derivative of a posterior receiver
probability with respect to its prior entry is at most
\(1/[4\mathrm c_0(1-\mathrm c_0)]\). Therefore
\[
 \operatorname{TV}(q_u^{C,p},q_u^{C,\widetilde p})
 \le\frac1{\mathrm c_0(1-\mathrm c_0)}
       \int_0^u\|p(v)-\widetilde p(v)\|_\infty dv.
\]
The variation-of-constants map for
\cref{eq:phil-actual-parent-learning} consequently has Lipschitz
constant at most \(\varepsilon u_*^2/[2\mathrm c_0(1-\mathrm c_0)]\) on an interval
of length \(u_*\). On a short enough interval it is a contraction
on continuous row paths in an interior cube. Its image stays in
that cube after further shortening the interval, since the row
velocity has absolute value at most \(\varepsilon\).
Moreover its solution satisfies
\[
 p_i^C(u)\ge p_i^C(0)e^{-\varepsilon u},\qquad
 1-p_i^C(u)\ge(1-p_i^C(0))e^{-\varepsilon u}.
\]
On any prescribed finite horizon these lower bounds give a common
interior cube and a positive contraction step length. Successive
steps give existence and uniqueness on that horizon, and then on
all finite horizons. The source marginal is continuous because its
jump rate is bounded, so the resulting rows are continuously
differentiable. This constructs learning from the parent's own law,
without imposing the fine row path after activation.

Let \(\overline\nu_{\mathrm T}^M\) be the retained Gaussian on
\(W_M\). Write \(\mathsf O^C_{\mathrm T,\mathsf h}\) for its
stationary held-\(M\) path law. The projection of
\(\mathsf O_{\mathrm T,\mathsf h}^M\) from
\cref{prop:phil-selected-pair-persistence} is exactly this law.
Only the learned-arm conclusions of that proposition are used, and
their proof does not use its frozen-arm hypothesis.
Conditioning any probability law on an event of probability
\(1-r\) changes its total variation by at most \(r\) after
any measurable projection. Thus the conditioned native retained
future obeys
\begin{equation}
 \operatorname{TV}
 (\mathsf P_{\theta,\mathsf h}^{\mathrm{fine},C},
                     \mathsf O^C_{\mathrm T,\mathsf h})
 \le\rho+(\mathsf h+2)\delta.
 \label{eq:phil-native-conditioned-held}
\end{equation}
This applies to a prefix event through \(\theta\). No condition is
placed on future coherence or future proposals. The same argument
at the activation time gives
\(\operatorname{TV}(q_0^C,
 \overline\nu_{\mathrm T}^M\otimes\delta_M)\le\rho+\delta\).

For the parent's own rows let \(\zeta_i^C\) be the probability in
row \(i\) assigned to the nonpartner of \(M\). Its equation and the
bound on any probability give
\[
 0\le\zeta_i^C(u)\le\zeta_i^C(0)+\varepsilon u
                  \le\delta+\varepsilon u\le d.
\]
The row-independent normalized sheet weights are at most the matching
weight one for \(0<\mathrm T\le1/100\) by
\cref{prop:phil-temperature-tracking}; hence, for any fixed row vector
with this bound,
\[
 \mu_{\mathrm T,p}(M)\ge\pi_M(p)
               \ge1-\sum_i\zeta_i\ge1-4d.
\]
Define the parent departure probability at a proposal by
\(k_p(y)=J_p((M,y),\{S'\ne M\})\).
Detailed balance, and the fact that the optimizing conditional
retained law on \(M\) is \(\overline\nu_{\mathrm T}^M\)
independently of the rows, give
\begin{equation}
 \int k_p(y)\overline\nu_{\mathrm T}^M(dy)
 \le\frac{1-\mu_{\mathrm T,p}(M)}{\mu_{\mathrm T,p}(M)}
 \le\frac{4d}{1-4d}.
 \label{eq:phil-parent-averaged-departure}
\end{equation}
The first inequality bounds stationary flux from \(M\) to its
complement by the complement's stationary mass.

Maximally couple the parent activation with the stationary retained
matching law and use common Brownian increments, clocks, and parent
proposal draws up to the first parent departure. Before departure,
the continuous paths agree because the matching drift splits, and
self-proposals do not move the retained coordinate. The row path
\(p^C(u)\) constructed above is deterministic, even though it was
computed from the parent's own law. Under the stationary comparison
path, its marginal at every time is
\(\overline\nu_{\mathrm T}^M\). Bound the first-departure
compensator with its stopped indicator by the compensator evaluated
on that whole stationary path. Applying
\cref{eq:phil-parent-averaged-departure} at each deterministic
\(p^C(u)\) gives
\[
 \Pr(\text{parent departure before }\mathsf h
                         \text{ while initially coupled})
 \le\frac{4\mathsf h d}{1-4d}.
\]
No Gaussian assertion is made about the parent's law conditioned on
survival or about the native hidden variables. We have proved
\[
 \operatorname{TV}(\mathsf P_{\mathsf h}^C,
                         \mathsf O^C_{\mathrm T,\mathsf h})
 \le\rho+\delta+\frac{4\mathsf h d}{1-4d}.
\]
Combining this with \cref{eq:phil-native-conditioned-held} proves
\cref{eq:phil-actual-parent-path-error}.

For the displayed joint choices,
\(d\le\mathrm T+\mathrm T^{7/4}\le2\mathrm T\),
\(\mathsf h\delta\le\mathrm T^{3/4}\), and
\(\rho=O(\mathrm T^{1/4})\).
The activation and following window fit inside the selection window
because \(1+\mathrm T^{-1/4}\le t_{\mathrm L}\) for sufficiently
small temperature under the mobility cap. The departure term is
\(O(\mathrm T^{3/4})\), giving the stated total error.
The Gaussian comparison laws served only to bound that error; the
parent's activation and subsequent learning remain the laws in the
statement. \hfill\(\square\)




